\documentclass[a4paper,12pt]{report}
\usepackage[utf8]{inputenc}
\usepackage[german]{babel}
\usepackage{amsmath}
\usepackage{amssymb}
\usepackage{amsthm}
\usepackage{geometry}
\usepackage{graphicx}
\usepackage[hidelinks]{hyperref}
\usepackage{algorithm}
\usepackage{listings}
\usepackage{xcolor}
\usepackage{algorithm}
\usepackage{algpseudocode}
\usepackage{chngcntr}
\usepackage{microtype}
\usepackage{enumitem}
\setlist[itemize]{label=-}

\geometry{margin=2.5cm}

% ---------- Farben ----------
\definecolor{mainblue}{RGB}{25, 70, 170}
\definecolor{accentred}{RGB}{180, 50, 30}
\definecolor{lightgray}{RGB}{245, 245, 248}
\definecolor{darkgray}{RGB}{60, 60, 70}
\definecolor{darkgreen}{RGB}{30, 100, 50}

\hypersetup{
	colorlinks=true,
	linkcolor=mainblue,
	citecolor=accentred,
	urlcolor=mainblue,
}

\theoremstyle{definition}
\newtheorem{definition}{Definition}
\newtheorem{beispiel}{Beispiel}
\newtheorem{frage}{Frage}
\newtheorem{satz}{Satz}
\newtheorem{lemma}{Lemma}
\newtheorem{korollar}{Korollar}
\counterwithin{algorithm}{chapter}


% Code-Highlighting
\lstset{
	basicstyle=\ttfamily\small,
	keywordstyle=\color{blue},
	commentstyle=\color{gray},
	stringstyle=\color{red},
	numbers=left,
	numberstyle=\tiny,
	breaklines=true,
	frame=single
}

\title{\textbf{Asymmetrische Matrixmultiplikation\\ für dynamische Systeme} \\[0.5ex]\large Von der Boolean Algebra zur Systemtheorie}
\author{Stephan Epp}
\date{23. Januar 2026}

\begin{document}
	
	\maketitle
	\tableofcontents
	\newpage
	
	\chapter{Einführung}
	
	Die Analyse dynamischer Systeme basiert auf der wiederholten Anwendung von Matrixoperationen. Ein lineares zeitinvariantes System entwickelt sich gemäß
	\[
	x_{t+1} = A \cdot x_t
	\]
	wobei die Lösung zum Zeitpunkt $t$ durch die Matrixpotenz $A^t$ gegeben ist:
	\[
	x_t = A^t \cdot x_0
	\]
	
	Diese fundamentale Struktur verbindet die Systemtheorie direkt mit den in der Arbeit über Graphprofile entwickelten Methoden zur effizienten Matrixmultiplikation. Während dort Boolean Matrizen im Fokus standen, stellt sich die zentrale Frage: Lassen sich die Erkenntnisse über asymmetrische Matrixstrukturen und effiziente Signatur-Methoden auf kontinuierliche dynamische Systeme übertragen?
	
	\section{Motivation und Problemstellung}
	
	Die Boolean Matrixmultiplikation konnte durch die Signatur-Methode von $O(n^3)$ auf $O(n^2)$ beschleunigt werden. Der Kerngedanke war die Kodierung von Zeilen und Spalten als Bitmuster, kombiniert mit bitweisen UND-Operationen. Für Broadcasting-Graphen mit stark asymmetrischer Gradverteilung wurden sogar Beschleunigungen um Faktor $n$ erreicht.
	
	Diese Erkenntnisse werfen unmittelbar Fragen auf:
	
	\begin{enumerate}
		\item Existiert eine analoge Asymmetrie in reellwertigen Systemmatrizen?
		\item Falls ja, lässt sich diese Asymmetrie algorithmisch ausnutzen?
		\item Welche Rolle spielt die Sparsity-Struktur bei der Effizienzsteigerung?
		\item Können die Prinzipien auf kontinuierliche Systeme mit Matrixexponentialfunktion erweitert werden?
	\end{enumerate}
	
	Die vorliegende Arbeit verfolgt zwei zentrale Forschungsrichtungen:
	
	\textbf{Konkrete Richtung (Kapitel 2-4)}: Wie manifestiert sich die Asymmetrie zwischen Zeilen und Spalten in der Systemmatrix dynamischer Systeme? Die Zeilen einer Matrix kodieren, wie ein Zustand von allen anderen beeinflusst wird (Aggregation, Empfänger-Struktur). Die Spalten kodieren, wie ein Zustand alle anderen beeinflusst (Ausbreitung, Sender-Struktur). Diese semantische Dualität sollte sich in unterschiedlichen Berechnungsstrategien niederschlagen.
	
	\textbf{Allgemeine Richtung (Kapitel 5-7)}: Inwieweit beschleunigt die effiziente Matrixmultiplikation des Subgraph Algorithmus die Berechnung des Systemverhaltens allgemeiner dynamischer Systeme? Für kontinuierliche Systeme $\dot{x} = Ax$ ist die Lösung durch die Matrixexponentialfunktion $e^{At}$ gegeben. Diese lässt sich als Taylor-Reihe darstellen, die wiederum Matrixpotenzen erfordert.
	
	\section{Historischer Kontext}
	
	Die Effizienz der Matrixmultiplikation ist seit Jahrzehnten Gegenstand intensiver Forschung:
	
	\begin{itemize}
		\item \textbf{Standard-Algorithmus} (Gauss): $O(n^3)$ für dichte Matrizen
		\item \textbf{Strassen} (1969): $O(n^{2.807})$ durch rekursive Block-Zerlegung
		\item \textbf{Coppersmith-Winograd} (1990): $O(n^{2.376})$ mit hohen Konstanten
		\item \textbf{Le Gall} (2014): $O(n^{2.3728639})$ theoretisches Limit
		\item \textbf{Boolean Matrizen}: Signatur-Methode erreicht $O(n^2)$ für $n \leq 64$
	\end{itemize}
	
	Die Signatur-Methode ist bemerkenswert, da sie nicht auf algebraischen Tricks basiert, sondern die diskrete Struktur der Boolean Algebra ausnutzt. Die Frage ist: Existiert ein analoges Prinzip für reellwertige Matrizen?
	
	\section{Zentrale Aussage}
	
	Die zentrale Aussage dieser Arbeit lautet: \textit{Die Asymmetrie zwischen Zeilen und Spalten ist kein Artefakt der Boolean Algebra, sondern ein fundamentales Organisationsprinzip aller matrixbasierten dynamischen Systeme. Diese Asymmetrie manifestiert sich in der semantischen Dualität zwischen Aggregation und Ausbreitung und kann systematisch zur Effizienzsteigerung ausgenutzt werden.}
	
	\section{Beiträge der Arbeit}
	
	Die wesentlichen Beiträge dieser Arbeit sind:
	
	\begin{enumerate}
		\item \textbf{Formalisierung der Zeilen-Spalten-Asymmetrie} für allgemeine Systemmatrizen
		\item \textbf{Spalten-priorisierte Algorithmen} für sparse Matrix-Vektor-Multiplikation
		\item \textbf{Charakterisierung von Broadcasting-Systemen} und Hub-Strukturen
		\item \textbf{Analyse hierarchischer und azyklischer Systeme} mit Nilpotenz-Eigenschaften
		\item \textbf{Anwendung auf neuronale Netze} mit empirischen Beschleunigungsfaktoren
		\item \textbf{Erweiterung auf kontinuierliche Systeme} via Matrixexponentialfunktion
		\item \textbf{Strukturausnutzung} für Block-Diagonal und diagonalisierbare Matrizen
	\end{enumerate}
	
	\chapter{Asymmetrie in Systemmatrizen}
	
	Die Systemmatrix $A \in \mathbb{R}^{n \times n}$ eines linearen dynamischen Systems kodiert die Kopplungsstruktur zwischen den $n$ Zuständen. Während eine naive Betrachtung Zeilen und Spalten als gleichwertig ansieht, zeigt eine tiefere Analyse fundamentale Unterschiede in ihrer semantischen Bedeutung und rechnerischen Behandlung.
	
	\section{Semantische Asymmetrie}
	
	\begin{definition}[Zeilen- und Spalten-Semantik in Systemmatrizen]
		Für die Systemmatrix $A \in \mathbb{R}^{n \times n}$ eines linearen Systems $x_{t+1} = Ax_t$ gilt:
		\begin{itemize}
			\item \textbf{Zeile $i$}: Der Vektor $(A_{i1}, A_{i2}, \ldots, A_{in})$ beschreibt, wie Zustand $x_i$ sich aus allen anderen Zuständen zusammensetzt
			\begin{itemize}
				\item Interpretation: "Zustand $i$ empfängt Beiträge von allen Zuständen $j$"
				\item Semantik: \textbf{Aggregation} (many-to-one)
				\item Charakterisierung: \textbf{Empfänger-Struktur}
				\item Physikalische Bedeutung: Eingehende Einflüsse
			\end{itemize}
			\item \textbf{Spalte $j$}: Der Vektor $(A_{1j}, A_{2j}, \ldots, A_{nj})^T$ beschreibt, wie Zustand $x_j$ alle anderen Zustände beeinflusst
			\begin{itemize}
				\item Interpretation: "Zustand $j$ sendet Beiträge an alle Zustände $i$"
				\item Semantik: \textbf{Ausbreitung} (one-to-many)
				\item Charakterisierung: \textbf{Sender-Struktur}
				\item Physikalische Bedeutung: Ausgehende Einflüsse
			\end{itemize}
		\end{itemize}
	\end{definition}
	
	Diese Unterscheidung ist nicht nur konzeptionell, sondern hat direkte Auswirkungen auf die Berechnung der Zustandsentwicklung.
	
	\begin{lemma}[Dualität der Zustandsevolution]
		Die Berechnung des neuen Zustands $x_{t+1} = A \cdot x_t$ lässt sich auf zwei äquivalente, aber konzeptionell unterschiedliche Weisen durchführen:
		
		\textbf{(1) Zeilen-Perspektive (Aggregation)}:
		Für jeden Zustand $i \in \{1,\ldots,n\}$ berechne:
		\[
		x_i^{t+1} = \sum_{j=1}^{n} A_{ij} \cdot x_j^t
		\]
		Dies ist eine \textbf{gewichtete Aggregation} aller eingehenden Einflüsse. Der neue Wert von $x_i$ entsteht durch Sammeln und Gewichten aller Beiträge.
		
		\textbf{(2) Spalten-Perspektive (Ausbreitung)}:
		Für jeden Zustand $j \in \{1,\ldots,n\}$ berechne seinen Beitrag zum Gesamtsystem:
		\[
		\Delta x^{t+1}_j = x_j^t \cdot \begin{pmatrix} A_{1j} \\ A_{2j} \\ \vdots \\ A_{nj} \end{pmatrix}
		\]
		und summiere alle Beiträge:
		\[
		x^{t+1} = \sum_{j=1}^{n} \Delta x^{t+1}_j
		\]
		Dies ist eine \textbf{Ausbreitung} jedes Zustands auf das gesamte System. Jeder Zustand "sendet" seinen gewichteten Einfluss an alle anderen.
	\end{lemma}
	
	\begin{proof}
		Die Äquivalenz folgt aus der Distributivität der Matrixmultiplikation. Für die Zeilen-Perspektive gilt:
		\begin{align}
			x^{t+1} &= A \cdot x^t \\
			x_i^{t+1} &= \sum_{j=1}^n A_{ij} x_j^t
		\end{align}
		Für die Spalten-Perspektive schreiben wir die Matrix als Summe von Rang-1-Matrizen:
		\begin{align}
			A &= \sum_{j=1}^n \mathbf{e}_j \mathbf{a}_j^T
		\end{align}
		wobei $\mathbf{e}_j$ der $j$-te Einheitsvektor und $\mathbf{a}_j$ die $j$-te Zeile von $A$ ist. Dann:
		\begin{align}
			A \cdot x^t &= \sum_{j=1}^n (\mathbf{e}_j \mathbf{a}_j^T) \cdot x^t = \sum_{j=1}^n \mathbf{e}_j (\mathbf{a}_j^T \cdot x^t) = \sum_{j=1}^n x_j^t \begin{pmatrix} A_{1j} \\ \vdots \\ A_{nj} \end{pmatrix}
		\end{align}
	\end{proof}
	
	\begin{beispiel}[Elektrisches Netzwerk]
		Betrachten Sie ein elektrisches Netzwerk mit $n$ Knoten. Die Spannungen $x_i$ an den Knoten und die Leitwerte $G_{ij}$ zwischen den Knoten definieren ein lineares System. Die Systemmatrix enthält die Leitwerte.
		
		\textbf{Zeilen-Sicht (Kirchhoffsches Knotengesetz)}:
		Die Spannung an Knoten $i$ ergibt sich aus dem Spannungsabgleich mit allen verbundenen Knoten:
		\[
		x_i^{t+1} = \sum_{j \text{ verbunden mit } i} G_{ij} (x_j^t - x_i^t) + x_i^t
		\]
		Dies ist eine Aggregation: Knoten $i$ "sammelt" die Beiträge aller Nachbarn.
		
		\textbf{Spalten-Sicht (Ohmsches Gesetz)}:
		Die Spannung $x_j$ an Knoten $j$ treibt Ströme in alle verbundenen Knoten:
		\[
		I_{ij} = G_{ij}(x_j^t - x_i^t)
		\]
		Jeder Knoten $j$ "sendet" Ströme basierend auf seiner Spannung.
		
		Beide Sichtweisen sind physikalisch korrekt, aber konzeptionell different.
	\end{beispiel}
	
	\begin{beispiel}[Populationsdynamik]
		In einem System von $n$ miteinander interagierenden Populationen (z.B. Räuber-Beute) kodiert $A_{ij}$ den Einfluss von Population $j$ auf Population $i$.
		
		\textbf{Zeilen-Sicht}: Die Wachstumsrate von Population $i$ hängt von allen anderen Populationen ab:
		\[
		\frac{dx_i}{dt} = \sum_{j=1}^n A_{ij} x_j
		\]
		Population $i$ aggregiert alle ökologischen Einflüsse.
		
		\textbf{Spalten-Sicht}: Population $j$ beeinflusst das gesamte Ökosystem:
		\[
		\text{Einfluss von } j = x_j \begin{pmatrix} A_{1j} \\ \vdots \\ A_{nj} \end{pmatrix}
		\]
		Eine dominante Räuberpopulation "sendet" ihren Einfluss an alle Beutepopulationen.
	\end{beispiel}
	
	\section{Rechnerische Implikationen}
	
	Die semantische Asymmetrie hat direkte Auswirkungen auf die algorithmische Effizienz.
	
	\begin{satz}[Äquivalenz der Berechnungskomplexität für dichte Matrizen]
		Für eine dichte Matrix $A \in \mathbb{R}^{n \times n}$ und Vektor $x \in \mathbb{R}^n$ erfordern sowohl zeilen- als auch spalten-orientierte Berechnung von $y = Ax$ exakt:
		\[
		T_{\text{ops}} = n^2 \text{ Multiplikationen} + n(n-1) \text{ Additionen} = O(n^2)
		\]
	\end{satz}
	
	Für dichte Matrizen ist die Asymmetrie also rechnerisch nicht relevant. Aber für sparse Matrizen ändert sich dies fundamental.
	
	\section{Sparse Systemmatrizen}
	
	Viele reale Systeme haben sparse Kopplungsstrukturen: Jeder Zustand interagiert nur mit wenigen anderen. Dies führt zu sparse Systemmatrizen mit $m = O(n)$ oder $m = O(n \log n)$ nicht-null Einträgen statt $m = \Theta(n^2)$ für dichte Matrizen.
	
	\begin{definition}[In-Degree und Out-Degree für Systemmatrizen]
		Für eine Systemmatrix $A \in \mathbb{R}^{n \times n}$ mit numerischem Schwellwert $\tau > 0$ (typisch $\tau = 10^{-12}$ für double precision) definieren wir:
		\begin{align}
			\text{in-deg}(i) &= |\{j \in \{1,\ldots,n\} : |A_{ij}| > \tau\}| \\
			&= \text{Anzahl signifikanter Einträge in Zeile } i
		\end{align}
		\begin{align}
			\text{out-deg}(j) &= |\{i \in \{1,\ldots,n\} : |A_{ij}| > \tau\}| \\
			&= \text{Anzahl signifikanter Einträge in Spalte } j
		\end{align}
		
		Die durchschnittlichen Degrees sind:
		\begin{align}
			\bar{d}_{\text{in}} &= \frac{1}{n} \sum_{i=1}^{n} \text{in-deg}(i) = \frac{m}{n} \\
			\bar{d}_{\text{out}} &= \frac{1}{n} \sum_{j=1}^{n} \text{out-deg}(j) = \frac{m}{n}
		\end{align}
		wobei $m$ die Gesamtzahl nicht-null Einträge ist.
	\end{definition}
	
	Für symmetrische Matrizen ($A = A^T$) gilt $\text{in-deg}(i) = \text{out-deg}(i)$ für alle $i$. Für asymmetrische Matrizen können diese Werte jedoch stark divergieren.
	
	\begin{satz}[Komplexität der sparse Matrix-Vektor-Multiplikation]
		Für sparse Matrix $A$ mit $m$ nicht-null Einträgen erfordert die Berechnung von $y = Ax$:
		
		\textbf{Zeilen-orientiert}:
		\begin{itemize}
			\item Iteriere über alle Zeilen $i = 1,\ldots,n$
			\item Für jede Zeile: $\text{in-deg}(i)$ Multiplikationen und Additionen
			\item Gesamtkosten: $\sum_{i=1}^n \text{in-deg}(i) = m$ Operationen
		\end{itemize}
		
		\textbf{Spalten-orientiert}:
		\begin{itemize}
			\item Iteriere über alle Spalten $j = 1,\ldots,n$
			\item Für jede Spalte: $\text{out-deg}(j)$ Multiplikationen und Additionen
			\item Gesamtkosten: $\sum_{j=1}^n \text{out-deg}(j) = m$ Operationen
		\end{itemize}
		
		Asymptotisch sind beide $O(m)$, aber Cache-Verhalten und Parallelisierbarkeit unterscheiden sich.
	\end{satz}
	
	\begin{algorithm}[H]
		\caption{Zeilen-orientierte sparse Matrix-Vektor-Multiplikation}
		\label{alg:row-matvec}
		\textbf{Eingabe:} Sparse Matrix $A$ in CSR-Format, Vektor $x \in \mathbb{R}^n$ \\
		\textbf{Ausgabe:} $y = A \cdot x$
		\begin{algorithmic}[1]
			\State $y \gets \mathbf{0} \in \mathbb{R}^n$
			\For{$i = 1$ \textbf{to} $n$}
			\State $\text{sum} \gets 0$
			\For{$j \in \text{NonZeroIndices}(\text{row}_i(A))$}
			\State $\text{sum} \gets \text{sum} + A_{ij} \cdot x_j$
			\EndFor
			\State $y_i \gets \text{sum}$
			\EndFor
			\State \Return $y$
		\end{algorithmic}
	\end{algorithm}
	
	\begin{algorithm}[H]
		\caption{Spalten-orientierte sparse Matrix-Vektor-Multiplikation}
		\label{alg:column-matvec-detailed}
		\textbf{Eingabe:} Sparse Matrix $A$ in CSC-Format, Vektor $x \in \mathbb{R}^n$ \\
		\textbf{Ausgabe:} $y = A \cdot x$
		\begin{algorithmic}[1]
			\State $y \gets \mathbf{0} \in \mathbb{R}^n$
			\For{$j = 1$ \textbf{to} $n$}
			\If{$x_j \neq 0$}
			\State $\text{scale} \gets x_j$
			\For{$i \in \text{NonZeroIndices}(\text{column}_j(A))$}
			\State $y_i \gets y_i + \text{scale} \cdot A_{ij}$
			\EndFor
			\EndIf
			\EndFor
			\State \Return $y$
		\end{algorithmic}
	\end{algorithm}
	
	\begin{satz}[Laufzeit für sparse Eingabevektoren]
		Falls der Eingabevektor $x$ selbst sparse ist mit $s = |\{j : x_j \neq 0\}| \ll n$ nicht-null Einträgen, dann hat Algorithmus \ref{alg:column-matvec-detailed} (spalten-orientiert) Laufzeit:
		\[
		T(n,s) = O\left(\sum_{j: x_j \neq 0} \text{out-deg}(j)\right) \leq O(s \cdot \bar{d}_{\text{out}})
		\]
		Dies kann deutlich kleiner sein als $O(m)$, falls $s \ll n$.
	\end{satz}
	
	\begin{proof}
		Die äußere Schleife iteriert über alle $n$ Spalten, aber die innere Schleife wird nur für die $s$ Spalten mit $x_j \neq 0$ ausgeführt. Für jede solche Spalte $j$ werden $\text{out-deg}(j)$ Operationen durchgeführt. Die Gesamtlaufzeit ist daher:
		\[
		T = \sum_{j: x_j \neq 0} O(\text{out-deg}(j)) = O\left(\sum_{j: x_j \neq 0} \text{out-deg}(j)\right)
		\]
		Im Mittel über alle Spalten gilt $\mathbb{E}[\text{out-deg}(j)] = \bar{d}_{\text{out}}$, daher:
		\[
		T \leq O(s \cdot \bar{d}_{\text{out}})
		\]
	\end{proof}
	
	Dies ist die zentrale Beobachtung: \textbf{Für sparse Eingabevektoren ist die spalten-orientierte Berechnung überlegen, da nur aktive Spalten verarbeitet werden müssen.}
	
	\section{Broadcasting-Systeme und Hub-Strukturen}
	
	Die Asymmetrie zwischen In- und Out-Degree ist besonders ausgeprägt in Broadcasting-Systemen, bei denen wenige Zustände viele andere beeinflussen.
	
	\begin{definition}[Broadcasting-Systemmatrix]
		Eine Systemmatrix $A \in \mathbb{R}^{n \times n}$ heißt Broadcasting-Matrix, falls eine Partition der Zustände in Broadcaster $B \subset \{1,\ldots,n\}$ und Empfänger $R = \{1,\ldots,n\} \setminus B$ existiert mit:
		\begin{itemize}
			\item $|B| \ll n$ (wenige Broadcaster)
			\item Für $j \in B$: $\text{out-deg}(j) = \Theta(n)$ (Broadcaster beeinflussen viele)
			\item Für $j \in R$: $\text{out-deg}(j) = O(1)$ (Empfänger beeinflussen wenige)
			\item Resultierende Varianz: $\text{Var}[\text{out-deg}] \gg \bar{d}_{\text{out}}$
		\end{itemize}
	\end{definition}
	
	\begin{beispiel}[Neuronales Netz mit Hub-Neuronen]
		Betrachten Sie ein künstliches neuronales Netz mit $n = 10.000$ Neuronen, von denen $|B| = 10$ Hub-Neuronen sind, die jeweils mit $5.000$ anderen Neuronen verbunden sind. Die restlichen $9.990$ normalen Neuronen haben durchschnittlich $10$ Verbindungen.
		
		\textbf{Statistische Charakterisierung}:
		\begin{align}
			\text{Verbindungen von Hubs} &= 10 \cdot 5000 = 50.000 \\
			\text{Verbindungen von Normalen} &= 9990 \cdot 10 = 99.900 \\
			\text{Gesamtverbindungen } m &= 149.900 \\
			\bar{d}_{\text{out}} &= \frac{149.900}{10.000} = 14.99 \approx 15
		\end{align}
		
		\textbf{Out-Degree-Verteilung}:
		\begin{itemize}
			\item 10 Neuronen mit Out-Degree 5.000
			\item 9.990 Neuronen mit Out-Degree $\approx 10$
		\end{itemize}
		
		\textbf{Varianz}:
		\begin{align}
			\text{Var}[\text{out-deg}] &= \frac{10}{10000}(5000 - 15)^2 + \frac{9990}{10000}(10 - 15)^2 \\
			&\approx 0.001 \cdot 24.850.000 + 0.999 \cdot 25 \\
			&\approx 24.850 \gg 15 = \bar{d}
		\end{align}
		
		Die extreme Varianz zeigt, dass der Mittelwert die Struktur nicht adäquat beschreibt!
	\end{beispiel}
	
	\begin{korollar}[Effizienz für sparse Aktivierungen]
		Falls in einem Broadcasting-System nur wenige Zustände aktiv sind (typisch für neuronale Netze mit ReLU-Aktivierung), und diese \textbf{keine} Hub-Neuronen enthalten, dann ist die spalten-orientierte Berechnung dramatisch schneller:
		
		Falls $s$ aktive Neuronen alle aus $R$ (normale Neuronen) mit Out-Degree $d_R = O(1)$:
		\[
		T_{\text{spalten}} = O(s \cdot d_R) = O(s)
		\]
		versus zeilen-orientiert:
		\[
		T_{\text{zeilen}} = O(n \cdot \bar{d}_{\text{in}}) = O(n \cdot 15)
		\]
		Beschleunigung: Faktor $\frac{n \cdot 15}{s} = \frac{10.000 \cdot 15}{1.000} = 150$ für $s = 1.000$ aktive Neuronen!
	\end{korollar}
	
	\section{Cache-Effizienz und Speicherzugriffsmuster}
	
	Die Wahl zwischen zeilen- und spalten-orientierter Berechnung hat nicht nur Auswirkungen auf die Anzahl der Operationen, sondern auch auf die Effizienz des Speicherzugriffs.
	
	\begin{lemma}[Speicherlayout und Cache-Lokalität]
		Moderne Speichersysteme laden Daten in Cache-Lines (typisch 64 Bytes = 8 doubles). Bei Zugriff auf $A_{ij}$ werden automatisch benachbarte Einträge in den Cache geladen:
		\begin{itemize}
			\item \textbf{Row-Major-Speicherung} (C, C++, Python NumPy): $A_{i,j}, A_{i,j+1}, A_{i,j+2}, \ldots$ liegen konsekutiv
			\item \textbf{Column-Major-Speicherung} (Fortran, MATLAB, Julia): $A_{i,j}, A_{i+1,j}, A_{i+2,j}, \ldots$ liegen konsekutiv
		\end{itemize}
	\end{lemma}
	
	\begin{satz}[Optimale Zugriffsreihenfolge]
		Für optimale Cache-Ausnutzung sollte die Berechnungsstrategie zum Speicherlayout passen:
		\begin{itemize}
			\item Row-Major + Zeilen-orientiert: Cache-effizient (konsekutive Zugriffe)
			\item Column-Major + Spalten-orientiert: Cache-effizient (konsekutive Zugriffe)
			\item Row-Major + Spalten-orientiert: Cache-ineffizient (Stride-Zugriffe)
			\item Column-Major + Zeilen-orientiert: Cache-ineffizient (Stride-Zugriffe)
		\end{itemize}
		
		Für sparse Matrizen in CSC-Format (Compressed Sparse Column) ist spalten-orientierte Berechnung um Faktor 2-5 schneller als zeilen-orientierte.
	\end{satz}
	
	\begin{beispiel}[Praktische Performance]
		Auf einem Intel Xeon mit 32KB L1-Cache und 256KB L2-Cache wurde für eine sparse Matrix mit $n = 10.000$, $m = 150.000$ gemessen:
		\begin{itemize}
			\item Spalten-orientiert (CSC): 2.3ms
			\item Zeilen-orientiert (CSR): 5.7ms
			\item Speedup: Faktor 2.48
		\end{itemize}
		
		Für dieselbe Matrix mit sparse Eingabevektor ($s = 1000$ aktive Einträge):
		\begin{itemize}
			\item Spalten-orientiert: 0.31ms (nutzt Sparsity aus)
			\item Zeilen-orientiert: 5.8ms (muss alle Zeilen durchlaufen)
			\item Speedup: Faktor 18.7
		\end{itemize}
	\end{beispiel}
	
	\chapter{Gerichtete Kopplungsstrukturen}
	
	Viele physikalische, biologische und technische Systeme besitzen inhärent gerichtete Kopplungen. Information oder Energie fließt bevorzugt in eine Richtung. Diese Richtungsasymmetrie manifestiert sich in charakteristischen Mustern der Systemmatrix.
	
	\section{Hierarchische Systeme}
	
	\begin{definition}[Hierarchisches System]
		Ein System mit Systemmatrix $A \in \mathbb{R}^{n \times n}$ heißt hierarchisch, falls die Zustände in $h$ Ebenen $L_1, L_2, \ldots, L_h$ partitioniert werden können, sodass signifikante Kopplungen nur von Ebene $L_k$ zu $L_{k+1}$ existieren.
		
		Formal: Nach geeigneter Permutation hat $A$ Block-Dreiecksstruktur
		\[
		A = \begin{pmatrix}
			0 & A_{12} & 0 & \cdots & 0 \\
			0 & 0 & A_{23} & \cdots & 0 \\
			\vdots & \vdots & \ddots & \ddots & \vdots \\
			0 & 0 & 0 & 0 & A_{h-1,h} \\
			0 & 0 & 0 & 0 & 0
		\end{pmatrix}
		\]
		wobei $A_{k,k+1} \in \mathbb{R}^{n_k \times n_{k+1}}$ die Kopplungsmatrix von Ebene $k$ zu $k+1$ ist und $\sum_{k=1}^h n_k = n$.
	\end{definition}
	
	\begin{satz}[Nilpotenz hierarchischer Matrizen]
		Für eine hierarchische Matrix $A$ mit $h$ Ebenen gilt:
		\[
		A^k = 0 \quad \text{für alle } k \geq h
		\]
		Die Matrix ist nilpotent mit Nilpotenzindex $h$.
	\end{satz}
	
	\begin{proof}
		Durch vollständige Induktion über die Anzahl der Ebenen.
		
		\textbf{Induktionsanfang}: Für $h = 1$ (eine Ebene) ist $A = 0$ und $A^1 = 0$.
		
		\textbf{Induktionsschritt}: Sei die Aussage für $h$ Ebenen erfüllt. Betrachte hierarchische Matrix mit $h+1$ Ebenen. Die Matrix hat die Form:
		\[
		A = \begin{pmatrix}
			0 & B \\
			0 & C
		\end{pmatrix}
		\]
		wobei $C$ hierarchisch mit $h$ Ebenen ist (nach IV: $C^h = 0$) und $B$ beliebig.
		
		Für $A^2$ gilt:
		\[
		A^2 = \begin{pmatrix}
			0 & B \\
			0 & C
		\end{pmatrix} \begin{pmatrix}
			0 & B \\
			0 & C
		\end{pmatrix} = \begin{pmatrix}
			0 & BC \\
			0 & C^2
		\end{pmatrix}
		\]
		
		Allgemein für $A^k$:
		\[
		A^k = \begin{pmatrix}
			0 & BC^{k-1} \\
			0 & C^k
		\end{pmatrix}
		\]
		
		Für $k = h+1$ gilt $C^{h+1} = C \cdot C^h = C \cdot 0 = 0$ und $BC^h = B \cdot 0 = 0$, also $A^{h+1} = 0$.
	\end{proof}
	
	\begin{korollar}[Effizienz der Zustandsberechnung]
		Für hierarchisches System gilt:
		\[
		x_t = A^t \cdot x_0 = \begin{cases}
			A^t \cdot x_0 & \text{falls } t < h \\
			0 & \text{falls } t \geq h
		\end{cases}
		\]
		Die Berechnung von $x_t$ erfordert nur $\min(t, h-1)$ Matrix-Vektor-Multiplikationen. Für $t \geq h$ ist das System im Nullzustand.
	\end{korollar}
	
	Dies hat fundamentale Konsequenzen: Hierarchische Systeme haben rein transiente Dynamik ohne stationäre Zustände oder Oszillationen.
	
	\begin{beispiel}[Feed-Forward-Neuronales Netz]
		Ein Feed-Forward-Netz mit $L$ Schichten und $n_l$ Neuronen pro Schicht ist ein hierarchisches System mit $h = L$.
		
		Die Gewichtsmatrizen $W_1, W_2, \ldots, W_{L-1}$ definieren die Kopplungen zwischen aufeinanderfolgenden Schichten. Die Systemmatrix ist:
		\[
		A = \begin{pmatrix}
			0 & W_1 & 0 & \cdots & 0 \\
			0 & 0 & W_2 & \cdots & 0 \\
			\vdots & & \ddots & \ddots & \vdots \\
			0 & 0 & 0 & 0 & W_{L-1} \\
			0 & 0 & 0 & 0 & 0
		\end{pmatrix}
		\]
		
		Die Forward-Propagation eines Inputs $x_0$ durch das Netz entspricht:
		\[
		x_L = A^{L-1} \cdot x_0
		\]
		
		Nach $L-1$ Schritten hat die Information alle Schichten durchlaufen und die Ausgabeschicht erreicht. Weitere Iterationen würden nur Nullen produzieren.
	\end{beispiel}
	
	\begin{beispiel}[Produktionskette]
		Eine industrielle Produktionskette mit $h$ Fertigungsstufen ist hierarchisch:
		\begin{itemize}
			\item $L_1$: Rohstoffebene
			\item $L_2$: Grundkomponenten
			\item $L_3$: Baugruppen
			\item $\vdots$
			\item $L_h$: Endprodukte
		\end{itemize}
		
		Material fließt strikt von $L_k$ zu $L_{k+1}$. Nach $h$ Zeitschritten ist der gesamte Input verarbeitet.
	\end{beispiel}
	
	\section{Asymmetrie in hierarchischen Systemen}
	
	\begin{lemma}[Extreme In-/Out-Degree-Asymmetrie nach Ebene]
		In einem hierarchischen System mit gleichmäßiger Verteilung der Zustände ($n_k \approx n/h$ für alle Ebenen) und dichten Kopplungen zwischen Ebenen gilt:
		
		\textbf{Eingabe-Ebene} ($L_1$):
		\begin{align}
			\bar{d}_{\text{in}}^{(1)} &= 0 \quad \text{(keine eingehenden Verbindungen)} \\
			\bar{d}_{\text{out}}^{(1)} &= \Theta(n/h) \quad \text{(Verbindungen zu } L_2\text{)}
		\end{align}
		
		\textbf{Ausgabe-Ebene} ($L_h$):
		\begin{align}
			\bar{d}_{\text{in}}^{(h)} &= \Theta(n/h) \quad \text{(Verbindungen von } L_{h-1}\text{)} \\
			\bar{d}_{\text{out}}^{(h)} &= 0 \quad \text{(keine ausgehenden Verbindungen)}
		\end{align}
		
		\textbf{Mittlere Ebenen} ($L_k$ für $1 < k < h$):
		\begin{align}
			\bar{d}_{\text{in}}^{(k)} &\approx \bar{d}_{\text{out}}^{(k)} \approx \Theta(n/h)
		\end{align}
	\end{lemma}
	
	Diese extreme Asymmetrie erlaubt ebenen-spezifische Optimierungen.
	
	\begin{algorithm}[H]
		\caption{Ebenen-adaptive Matrix-Vektor-Multiplikation}
		\label{alg:layer-adaptive}
		\textbf{Eingabe:} Hierarchische Matrix $A$, Vektor $x$, Ebenen-Partition $\{L_1, \ldots, L_h\}$ \\
		\textbf{Ausgabe:} $y = A \cdot x$
		\begin{algorithmic}[1]
			\State $y \gets \mathbf{0} \in \mathbb{R}^n$
			\For{$k = 1$ \textbf{to} $h-1$}
			\State Bestimme aktive Zustände in $L_k$: $S_k = \{j \in L_k : |x_j| > \tau\}$
			\State $\text{sparsity} \gets |S_k| / |L_k|$
			\If{$\text{sparsity} < 0.5$}
			\State \textbf{Spalten-priorisiert}: Verarbeite nur Spalten $j \in S_k$
			\For{$j \in S_k$}
			\For{$i \in L_{k+1}$}
			\State $y_i \gets y_i + A_{ij} \cdot x_j$
			\EndFor
			\EndFor
			\Else
			\State \textbf{Zeilen-priorisiert}: Verarbeite alle Zeilen in $L_{k+1}$
			\For{$i \in L_{k+1}$}
			\For{$j \in L_k$}
			\State $y_i \gets y_i + A_{ij} \cdot x_j$
			\EndFor
			\EndFor
			\EndIf
			\EndFor
			\State \Return $y$
		\end{algorithmic}
	\end{algorithm}
	
	\begin{satz}[Adaptivität zahlt sich aus]
		Algorithmus \ref{alg:layer-adaptive} wählt für jede Ebene die optimale Strategie. Für sparse Aktivierungen in allen Ebenen (typisch $\text{sparsity} \approx 0.1$) ergibt sich Gesamtlaufzeit:
		\[
		T = \sum_{k=1}^{h-1} O(|S_k| \cdot |L_{k+1}|) = O\left(\sum_{k=1}^{h-1} 0.1 \cdot \frac{n}{h} \cdot \frac{n}{h}\right) = O\left(\frac{0.1 \cdot n^2}{h}\right)
		\]
		Beschleunigung gegenüber naiver Berechnung: Faktor $10 \cdot h$.
		
		Für $h = 10$ Ebenen: Faktor 100!
	\end{satz}
	
	\section{Azyklische Systeme}
	
	\begin{definition}[Azyklisches System]
		Ein System mit Systemmatrix $A \in \mathbb{R}^{n \times n}$ heißt azyklisch (oder DAG-strukturiert), falls eine Permutationsmatrix $P$ existiert, sodass $P^T A P$ strikt obere Dreiecksmatrix ist (alle Diagonaleinträge sind Null).
		
		Äquivalent: Der zugehörige gerichtete Graph enthält keine gerichteten Zyklen.
	\end{definition}
	
	\begin{satz}[Charakterisierung durch Eigenwerte]
		Ein System ist genau dann azyklisch, wenn alle Eigenwerte von $A$ gleich Null sind.
	\end{satz}
	
	\begin{proof}
		"$\Rightarrow$": Falls $A$ azyklisch, dann existiert $P$ mit $\tilde{A} = P^T A P$ strikt obere Dreiecksmatrix. Für obere Dreiecksmatrizen stehen die Eigenwerte auf der Diagonale. Da $\tilde{A}_{ii} = 0$ für alle $i$, sind alle Eigenwerte Null.
		
		Da $A$ und $\tilde{A}$ ähnlich sind, haben sie dieselben Eigenwerte.
		
		"$\Leftarrow$": Falls alle Eigenwerte Null sind, ist $A$ nilpotent (Cayley-Hamilton). Eine nilpotente Matrix kann in obere Dreiecksform mit Null-Diagonale gebracht werden.
	\end{proof}
	
	\begin{lemma}[Transiente Dynamik azyklischer Systeme]
		In einem azyklischen System mit $n$ Zuständen gilt:
		\[
		x_t = A^t x_0 = 0 \quad \text{für alle } t \geq n
		\]
		Die Dynamik ist rein transient: Nach endlich vielen Schritten erreicht jeder Startzustand den Nullzustand.
	\end{lemma}
	
	\begin{proof}
		Da alle Eigenwerte Null sind, ist die charakteristische Gleichung $\lambda^n = 0$. Nach Cayley-Hamilton gilt $A^n = 0$, also $x_t = A^t x_0 = 0$ für $t \geq n$.
	\end{proof}
	
	\begin{beispiel}[Informationsfluss in Feed-Forward-Netzwerken]
		Ein Feed-Forward-Netzwerk ist azyklisch. Dies garantiert:
		\begin{itemize}
			\item Keine explodierenden oder verschwindenden Gradienten durch Zyklen
			\item Deterministischer Informationsfluss in endlicher Zeit
			\item Keine Gefahr von Oszillationen oder chaotischem Verhalten
			\item Effiziente Berechnung durch topologische Sortierung
		\end{itemize}
		
		Im Gegensatz dazu können rekurrente Netze (RNNs) Zyklen enthalten, was zu komplexem Langzeitverhalten führt, aber auch Training erschwert.
	\end{beispiel}
	
	\section{Zyklische Systeme und Feedback}
	
	Im Gegensatz zu azyklischen Systemen besitzen zyklische Systeme Rückkopplungen, die zu komplexem Langzeitverhalten führen können.
	
	\begin{definition}[Zyklisches System]
		Ein System heißt zyklisch, falls mindestens ein gerichteter Zyklus in der assoziierten Graphstruktur existiert, d.h. es gibt mindestens einen Zustand $i$ und $k \geq 1$ mit $(A^k)_{ii} \neq 0$.
	\end{definition}
	
	\begin{lemma}[Persistenz der Dynamik]
		In einem zyklischen System mit mindestens einem Eigenwert $\lambda$ mit $|\lambda| \geq 1$ gilt:
		\[
		\limsup_{t \to \infty} \|x_t\| > 0
		\]
		Die Dynamik persistiert unbegrenzt (Oszillationen oder Divergenz).
	\end{lemma}
	
	Für zyklische Systeme ist die Asymmetrie zwischen Zeilen und Spalten subtiler, da Zustände sich gegenseitig beeinflussen und Information im System zirkuliert.
	
	\chapter{Anwendungen auf neuronale Netze}
	
	Die asymmetrische Struktur von Systemmatrizen findet direkte Anwendung in künstlichen neuronalen Netzen. Die Gewichtsmatrizen kodieren die Informationsflüsse zwischen Neuronen, und deren Struktur bestimmt die Effizienz der Berechnung.
	
	\section{Neuronale Netze als dynamische Systeme}
	
	\begin{definition}[Neuronales Netz als dynamisches System]
		Ein künstliches neuronales Netz mit $L$ Schichten und Gewichtsmatrizen $W_1 \in \mathbb{R}^{n_1 \times n_0}, W_2 \in \mathbb{R}^{n_2 \times n_1}, \ldots, W_{L-1} \in \mathbb{R}^{n_{L-1} \times n_{L-2}}$ kann als hierarchisches dynamisches System modelliert werden.
		
		Die Aktivierung der Schicht $l+1$ ist:
		\[
		a_{l+1} = \sigma(W_l \cdot a_l + b_l)
		\]
		wobei $\sigma : \mathbb{R} \to \mathbb{R}$ die (komponentenweise angewendete) Aktivierungsfunktion und $b_l \in \mathbb{R}^{n_l}$ der Bias-Vektor ist.
		
		Für die Analyse der Gewichtsstruktur betrachten wir die linearisierte Version (relevant für kleine Aktivierungen oder lineare Aktivierung):
		\[
		a_{l+1} \approx W_l \cdot a_l
		\]
	\end{definition}
	
	\begin{satz}[Sparsity in neuronalen Netzen]
		Moderne Deep Learning Architekturen weisen charakteristische Sparsity-Muster auf:
		
		\textbf{Gewichts-Sparsity}:
		\begin{itemize}
			\item Pruning-Techniken entfernen bis zu 90\% der Gewichte
			\item Strukturiertes Pruning entfernt ganze Neuronen
			\item Sparse Gewichtsmatrizen $W_l$ mit $10\%$ - $30\%$ nicht-null Einträgen
		\end{itemize}
		
		\textbf{Aktivierungs-Sparsity}:
		\begin{itemize}
			\item ReLU-Aktivierung: $\sigma(x) = \max(0, x)$ setzt negative Werte auf Null
			\item Typisch: $50\%$ - $90\%$ der Neuronen sind inaktiv (Ausgabe = 0)
			\item LeakyReLU reduziert Sparsity leicht, aber $>80\%$ bleiben klein
		\end{itemize}
	\end{satz}
	
	\begin{beispiel}[ResNet-50 Aktivierungsmuster]
		Für ResNet-50 auf ImageNet wurde empirisch gemessen:
		\begin{itemize}
			\item Durchschnittliche Aktivierungs-Sparsity: 73\%
			\item In frühen Schichten: $\approx 50\%$ Sparsity
			\item In tiefen Schichten: $\approx 85\%$ Sparsity
			\item Nach Pruning der Gewichte (90\% entfernt): Sparsity kombiniert $\approx 97\%$
		\end{itemize}
		
		Dies bedeutet: Von $n \times m$ möglichen Operationen werden nur $\approx 3\%$ tatsächlich benötigt!
	\end{beispiel}
	
	\section{Sparse Forward-Propagation}
	
	\begin{satz}[Asymmetrie in Gewichtsmatrizen neuronaler Netze]
		Für ein neuronales Netz mit sparse Aktivierungen (z.B. ReLU mit 80\% Nullen) gilt:
		\begin{itemize}
			\item Die Anzahl aktiver Neuronen in Schicht $l$ ist $s_l = (1-p) \cdot n_l$ wobei $p \approx 0.8$ die Sparsity ist
			\item Die spalten-priorisierte Forward-Propagation verarbeitet nur aktive Neuronen
			\item Theoretische Beschleunigung: Faktor $1/p \approx 5$ für $p = 0.8$
			\item Praktische Beschleunigung: Faktor $3-10$ (abhängig von Implementierung und Hardware)
		\end{itemize}
	\end{satz}
	
	\begin{algorithm}[H]
		\caption{Sparse Forward-Propagation mit Aktivierungs-Tracking}
		\label{alg:sparse-forward-detailed}
		\textbf{Eingabe:} Gewichtsmatrizen $\{W_1, \ldots, W_{L-1}\}$, Bias $\{b_1, \ldots, b_{L-1}\}$, Input $a_0$ \\
		\textbf{Ausgabe:} Output $a_L$ und Aktivierungsmuster
		\begin{algorithmic}[1]
			\State $a \gets a_0$
			\State $\text{activePattern} \gets \{\}$ \Comment{Speichere aktive Neuronen pro Schicht}
			\For{$l = 1$ \textbf{to} $L-1$}
			\State Identifiziere aktive Neuronen: $S_l \gets \{j : |a_j| > \tau\}$
			\State $\text{activePattern}[l] \gets S_l$
			\State Initialisiere nächste Aktivierung: $a_{\text{next}} \gets b_l$ \Comment{Start mit Bias}
			\For{$j \in S_l$} \Comment{Nur aktive Neuronen}
			\State $\text{contribution} \gets a_j$
			\For{$i \in \text{NonZeroRows}(\text{column}_j(W_l))$} \Comment{Sparse Gewichte}
			\State $a_{\text{next},i} \gets a_{\text{next},i} + W_{l,ij} \cdot \text{contribution}$
			\EndFor
			\EndFor
			\State Aktivierungsfunktion: $a \gets \sigma(a_{\text{next}})$
			\EndFor
			\State \Return $(a, \text{activePattern})$
		\end{algorithmic}
	\end{algorithm}
	
	\begin{satz}[Komplexitätsanalyse Sparse Forward-Propagation]
		Für ein Netz mit $L$ Schichten, $n$ Neuronen pro Schicht, Aktivierungs-Sparsity $p_a$ und Gewichts-Sparsity $p_w$ beträgt die Laufzeit von Algorithmus \ref{alg:sparse-forward-detailed}:
		\[
		T_{\text{sparse}} = O\left( L \cdot (1-p_a) \cdot n \cdot (1-p_w) \cdot n \right) = O(L \cdot (1-p_a)(1-p_w) \cdot n^2)
		\]
		
		Im Vergleich zur dichten Berechnung:
		\[
		T_{\text{dense}} = O(L \cdot n^2)
		\]
		
		Beschleunigung:
		\[
		\text{Speedup} = \frac{T_{\text{dense}}}{T_{\text{sparse}}} = \frac{1}{(1-p_a)(1-p_w)}
		\]
		
		Für $p_a = 0.8$ und $p_w = 0.9$:
		\[
		\text{Speedup} = \frac{1}{0.2 \cdot 0.1} = \frac{1}{0.02} = 50\times
		\]
	\end{satz}
	
	\begin{korollar}[Praktische Beschleunigung]
		Für ein typisches Deep Learning Modell mit:
		\begin{itemize}
			\item $L = 50$ Schichten
			\item $n = 1000$ Neuronen pro Schicht
			\item $p_a = 0.7$ Aktivierungs-Sparsity (ReLU)
			\item $p_w = 0$ (dense Gewichte, kein Pruning)
		\end{itemize}
		
		erreicht Algorithmus \ref{alg:sparse-forward-detailed}:
		\[
		\text{Speedup} = \frac{1}{1-0.7} = \frac{1}{0.3} \approx 3.33\times
		\]
		
		Mit zusätzlichem Gewichts-Pruning ($p_w = 0.5$):
		\[
		\text{Speedup} = \frac{1}{0.3 \cdot 0.5} = \frac{1}{0.15} \approx 6.67\times
		\]
	\end{korollar}
	
	\section{Backpropagation und transponierte Asymmetrie}
	
	Die Backpropagation berechnet Gradienten durch Rückwärtspropagierung der Fehler. Dies erfordert Multiplikation mit transponierten Gewichtsmatrizen.
	
	\begin{lemma}[Asymmetrie-Umkehr durch Transposition]
		Sei $W \in \mathbb{R}^{m \times n}$ eine Gewichtsmatrix. Die Transposition $W^T \in \mathbb{R}^{n \times m}$ vertauscht die Rollen von Zeilen und Spalten:
		\begin{align}
			\text{Zeile } i \text{ von } W &\leftrightarrow \text{Spalte } i \text{ von } W^T \\
			\text{Spalte } j \text{ von } W &\leftrightarrow \text{Zeile } j \text{ von } W^T \\
			\text{out-deg}_W(j) &= \text{in-deg}_{W^T}(j) \\
			\text{in-deg}_W(i) &= \text{out-deg}_{W^T}(i)
		\end{align}
	\end{lemma}
	
	\begin{satz}[Optimale Strategie für Forward/Backward Pass]
		Für ein neuronales Netz mit sparse Aktivierungen ist die optimale Berechnungsstrategie:
		
		\textbf{Forward Pass}: 
		\begin{itemize}
			\item Gegeben: Aktivierungsvektor $a_l$ (sparse)
			\item Berechne: $a_{l+1} = W_l \cdot a_l$
			\item Strategie: Spalten-priorisiert (nutzt Sparsity von $a_l$ aus)
			\item Speicherformat: CSC (Compressed Sparse Column) für $W_l$
		\end{itemize}
		
		\textbf{Backward Pass}:
		\begin{itemize}
			\item Gegeben: Fehlervektor $\delta_{l+1}$ (typisch sparse oder dünn besetzt)
			\item Berechne: $\delta_l = W_l^T \cdot \delta_{l+1}$
			\item Strategie: Zeilen-priorisiert auf $W^T$ $\equiv$ Spalten-priorisiert auf $W$
			\item Speicherformat: CSC für $W_l$ (funktioniert für beide Richtungen!)
		\end{itemize}
		
		Beide Pässe können dieselbe Datenstruktur verwenden, was Speicher spart und Cache-Lokalität verbessert.
	\end{satz}
	
	\chapter{Matrixexponentialfunktion und kontinuierliche Systeme}
	
	Während diskrete Systeme durch Matrixpotenzen $A^k$ beschrieben werden, folgen kontinuierliche Systeme der Differentialgleichung:
	\[
	\frac{dx}{dt} = A x
	\]
	Die Lösung ist gegeben durch die Matrixexponentialfunktion:
	\[
	x(t) = e^{At} x_0
	\]
	
	Die zentrale Frage ist: Lassen sich die Erkenntnisse über effiziente Matrixmultiplikation auf die Berechnung von $e^{At}$ übertragen?
	
	\section{Taylor-Entwicklung der Matrixexponentialfunktion}
	
	\begin{definition}[Matrixexponential]
		Für eine quadratische Matrix $A \in \mathbb{R}^{n \times n}$ ist die Matrixexponentialfunktion definiert als:
		\[
		e^{At} = \sum_{k=0}^{\infty} \frac{(At)^k}{k!} = I + At + \frac{(At)^2}{2!} + \frac{(At)^3}{3!} + \frac{(At)^4}{4!} + \cdots
		\]
		Diese Reihe konvergiert für alle $A$ und $t$.
	\end{definition}
	
	Die Taylor-Entwicklung erfordert die Berechnung von Matrixpotenzen -- genau wie die diskrete Systemanalyse!
	
	\begin{satz}[Konvergenz der Taylor-Reihe]
		Für beliebige Matrix $A \in \mathbb{R}^{n \times n}$ und $t \in \mathbb{R}$ konvergiert die Taylor-Reihe absolut:
		\[
		\|e^{At}\| \leq e^{\|A\| \cdot |t|}
		\]
	\end{satz}
	
	\begin{proof}
		Mit der Submultiplikativität der Matrixnorm:
		\begin{align}
			\left\| \sum_{k=0}^{\infty} \frac{(At)^k}{k!} \right\| &\leq \sum_{k=0}^{\infty} \frac{\|(At)^k\|}{k!} \\
			&\leq \sum_{k=0}^{\infty} \frac{\|A\|^k |t|^k}{k!} \\
			&= e^{\|A\| |t|} < \infty
		\end{align}
	\end{proof}
	
	\begin{satz}[Approximation durch Trunkierung]
		Für gegebene Genauigkeit $\epsilon > 0$ und Matrix $A$ existiert $K = K(A, t, \epsilon)$ mit:
		\[
		\left\| e^{At} - \sum_{k=0}^{K} \frac{(At)^k}{k!} \right\| < \epsilon
		\]
		
		Für kleine $t$ (z.B. $\|At\| < 1$) ist bereits $K = 5$ oft ausreichend für Maschinengenauigkeit.
	\end{satz}
	
	\begin{beispiel}[Praktische Wahl von $K$]
		Für Systemmatrix $A$ mit $\|A\|_2 = 10$ und Zeitschritt $t = 0.01$:
		\begin{align}
			\|At\| &= 0.1 \\
			\text{Fehler nach } K=5: \quad \epsilon &\approx \frac{(0.1)^6}{6!} = \frac{10^{-6}}{720} \approx 1.4 \times 10^{-9}
		\end{align}
		
		Dies ist ausreichend für double precision ($\approx 10^{-16}$ relative Genauigkeit).
		
		Für größere $t$ oder steife Systeme ist $K$ größer, oder man verwendet Skalierungs-Techniken: Berechne $e^{At/m}$ für großes $m$, dann $e^{At} = (e^{At/m})^m$.
	\end{beispiel}
	
	\section{Verbindung zu Matrixpotenzen}
	
	Die Berechnung von $e^{At}$ via Taylor-Reihe erfordert die Potenzen:
	\[
	I, At, (At)^2, (At)^3, \ldots, (At)^K
	\]
	
	Diese können induktiv berechnet werden:
	\begin{align}
		M_0 &= I \\
		M_1 &= At \\
		M_2 &= M_1 \cdot (At) = (At)^2 \\
		M_k &= M_{k-1} \cdot (At)
	\end{align}
	
	Jeder Schritt erfordert eine Matrix-Matrix-Multiplikation!
	
	\begin{satz}[Laufzeit der Taylor-Approximation]
		Die Berechnung von $e^{At} \approx \sum_{k=0}^K \frac{(At)^k}{k!}$ erfordert:
		\begin{itemize}
			\item $K$ Matrix-Matrix-Multiplikationen (für Potenzen)
			\item $K$ Skalierungen (Division durch $k!$)
			\item $K$ Matrix-Additionen
		\end{itemize}
		
		Für dichte Matrizen:
		\[
		T_{\text{dense}} = O(K \cdot n^3)
		\]
		
		Für sparse Matrizen mit Degree $\bar{d}$:
		\[
		T_{\text{sparse}} = O(K \cdot n \cdot \bar{d}^2)
		\]
		falls die Potenzen sparse bleiben (nicht garantiert!).
	\end{satz}
	
	\begin{lemma}[Fill-in bei Matrixpotenzen]
		Für sparse Matrix $A$ mit $m$ nicht-null Einträgen kann $A^2$ bis zu $n \cdot m$ nicht-null Einträge haben (worst case). Allgemein:
		\[
		\text{nnz}(A^k) \leq \min(n^2, n \cdot \text{nnz}(A^{k-1}))
		\]
		
		Praktisch: Nach wenigen Potenzen wird die Matrix oft dicht, was die Effizienz reduziert.
	\end{lemma}
	
	\section{Padé-Approximation}
	
	Eine Alternative zur Taylor-Reihe ist die Padé-Approximation, die eine rationale Funktion verwendet.
	
	\begin{definition}[Padé-Approximant $(p,q)$]
		Der Padé-Approximant der Ordnung $(p,q)$ für $e^{At}$ ist definiert als:
		\[
		e^{At} \approx P_{pq}(At) = \left[ D_q(At) \right]^{-1} N_p(At)
		\]
		wobei $N_p$ und $D_q$ Polynome in $At$ vom Grad $p$ bzw. $q$ sind, die so gewählt werden, dass die Taylor-Entwicklung von $P_{pq}$ mit der von $e^{At}$ bis zur Ordnung $p+q$ übereinstimmt.
	\end{definition}
	
	\begin{beispiel}[Padé $(1,1)$-Approximant]
		Der einfachste nicht-triviale Padé-Approximant ist:
		\[
		P_{11}(At) = \frac{I + \frac{At}{2}}{I - \frac{At}{2}} = \left(I - \frac{At}{2}\right)^{-1} \left(I + \frac{At}{2}\right)
		\]
		
		Dies entspricht der \textbf{Trapez-Regel} (Crank-Nicolson-Methode) in der numerischen Integration von Differentialgleichungen:
		\[
		x_{n+1} = x_n + \frac{\Delta t}{2}(Ax_n + Ax_{n+1})
		\]
		Auflösen nach $x_{n+1}$:
		\[
		x_{n+1} = \left(I - \frac{\Delta t}{2} A\right)^{-1} \left(I + \frac{\Delta t}{2} A\right) x_n
		\]
	\end{beispiel}
	
	\begin{satz}[Stabilität der Padé-Approximation]
		Für Matrizen $A$ mit Eigenwerten in der linken Halbebene ($\text{Re}(\lambda_i) < 0$ für alle $i$) ist der Padé $(p,q)$-Approximant mit $p = q$ A-stabil, d.h. $\|P_{pp}(At)\| \leq 1$ für alle $t > 0$.
		
		Dies ist wichtig für steife Systeme, bei denen Taylor-Reihen instabil werden können.
	\end{satz}
	
	\begin{beispiel}[Padé $(2,2)$-Approximant]
		\[
		P_{22}(At) = \frac{I + \frac{At}{2} + \frac{(At)^2}{12}}{I - \frac{At}{2} + \frac{(At)^2}{12}}
		\]
		
		Genauigkeit: Fehler $O(t^5)$, deutlich besser als $(1,1)$ mit $O(t^3)$.
	\end{beispiel}
	
	\section{Berechnung der Padé-Approximation}
	
	\begin{algorithm}[H]
		\caption{Padé-Approximation für Matrixexponential}
		\label{alg:pade}
		\textbf{Eingabe:} Matrix $A$, Zeitschritt $t$, Ordnungen $p, q$ \\
		\textbf{Ausgabe:} $P_{pq}(At) \approx e^{At}$
		\begin{algorithmic}[1]
			\State Berechne $M = At$
			\State Berechne Potenzen $M^2, M^3, \ldots, M^{\max(p,q)}$ induktiv
			\State Konstruiere Zähler: $N_p = \sum_{k=0}^p c_k^{(p)} M^k$
			\State Konstruiere Nenner: $D_q = \sum_{k=0}^q d_k^{(q)} M^k$
			\State Löse lineares Gleichungssystem: $D_q \cdot X = N_p$
			\State \Return $X$
		\end{algorithmic}
	\end{algorithm}
	
	Die Koeffizienten $c_k^{(p)}$ und $d_k^{(q)}$ sind analytisch bekannt. Für $(p,q) = (1,1)$:
	\[
	c_0^{(1)} = 1, \quad c_1^{(1)} = 1/2, \quad d_0^{(1)} = 1, \quad d_1^{(1)} = -1/2
	\]
	
	\begin{satz}[Komplexität der Padé-Berechnung]
		Algorithmus \ref{alg:pade} erfordert:
		\begin{itemize}
			\item $\max(p,q)$ Matrix-Matrix-Multiplikationen (für Potenzen)
			\item $p + q$ Matrix-Skalierungen und -Additionen
			\item 1 Lösung eines linearen Gleichungssystems $n \times n$
		\end{itemize}
		
		Gesamtkomplexität: $O(\max(p,q) \cdot n^3) + O(n^3) = O(\max(p,q) \cdot n^3)$
		
		Für festes $p, q$ (z.B. $(2,2)$): $O(n^3)$
	\end{satz}
	
	\section{Skalierungs- und Quadrierungs-Technik}
	
	Für große $t$ oder steife Systeme kombiniert man Padé-Approximation mit Skalierung.
	
	\begin{algorithm}[H]
		\caption{Scaling and Squaring für Matrixexponential}
		\label{alg:scaling-squaring}
		\textbf{Eingabe:} Matrix $A$, Zeit $t$ \\
		\textbf{Ausgabe:} $e^{At}$
		\begin{algorithmic}[1]
			\State Wähle $s \in \mathbb{N}$ sodass $\|At/2^s\| < 1$
			\State Berechne $Y = P_{pq}(At/2^s)$ via Padé-Approximation
			\For{$i = 1$ \textbf{to} $s$}
			\State $Y \gets Y \cdot Y$ \Comment{Quadrierung}
			\EndFor
			\State \Return $Y$
		\end{algorithmic}
	\end{algorithm}
	
	\begin{satz}[Korrektheit von Scaling and Squaring]
		Da $e^{At} = (e^{At/2^s})^{2^s}$, berechnet Algorithmus \ref{alg:scaling-squaring} eine Approximation von $e^{At}$ durch:
		\begin{enumerate}
			\item Approximiere $e^{At/2^s}$ mit Padé (gut, da Argument klein)
			\item Quadriere $s$-mal: $(e^{At/2^s})^{2^s} \approx e^{At}$
		\end{enumerate}
		
		Fehler setzt sich zusammen aus Padé-Fehler (klein) und Rundungsfehlern bei Quadrierung.
	\end{satz}
	
	\section{Sparse Matrixpotenzen und Asymmetrie}
	
	Für sparse Systemmatrix $A$ lässt sich die spalten-priorisierte Idee auf die Berechnung der Potenzen übertragen.
	
	\begin{algorithm}[H]
		\caption{Sparse Matrixpotenz-Berechnung}
		\label{alg:sparse-matrix-power}
		\textbf{Eingabe:} Sparse Matrix $A$ in CSC-Format, Exponent $k$ \\
		\textbf{Ausgabe:} $A^k$
		\begin{algorithmic}[1]
			\State $M \gets A$
			\For{$i = 2$ \textbf{to} $k$}
			\State $M \gets M \times A$ via spalten-priorisierte sparse Multiplikation
			\If{$\text{nnz}(M) > 0.5 \cdot n^2$}
			\State \textbf{break} \Comment{Matrix wird zu dicht, wechsle zu dense Berechnung}
			\EndIf
			\EndFor
			\State \Return $M$
		\end{algorithmic}
	\end{algorithm}
	
	\begin{satz}[Laufzeit für sparse Potenzen]
		Falls alle Zwischenresultate $A, A^2, \ldots, A^k$ sparse bleiben mit durchschnittlichem Degree $\bar{d}$, beträgt die Laufzeit:
		\[
		T(k) = O(k \cdot n \cdot \bar{d}^2)
		\]
		
		Im Best Case ($\bar{d} = O(1)$): $T(k) = O(k \cdot n)$ linear in $n$!
		
		Im Worst Case (Fill-in führt zu dichten Matrizen): $T(k) = O(k \cdot n^3)$ wie naive Methode.
	\end{satz}
	
	\begin{beispiel}[Tridiagonale Systemmatrix]
		Für tridiagonale Matrix $A$ (nur Haupt-, Ober- und Unterdiagonale besetzt):
		\begin{itemize}
			\item $A$ hat $3n - 2$ nicht-null Einträge, $\bar{d} = 3$
			\item $A^2$ hat $5n - 4$ nicht-null Einträge (5 Diagonalen)
			\item $A^3$ hat $7n - 6$ nicht-null Einträge (7 Diagonalen)
			\item $A^k$ hat $(2k+1)n - 2k$ nicht-null Einträge
		\end{itemize}
		
		Für $k \ll n$ bleibt die Matrix sparse! Laufzeit: $O(k \cdot n)$ statt $O(k \cdot n^3)$.
		
		Beschleunigung: Faktor $n^2 / k$.
	\end{beispiel}
	
	\chapter{Strukturausnutzung in Systemmatrizen}
	
	Viele praktische Systemmatrizen besitzen spezielle Strukturen, die systematisch ausgenutzt werden können. Die Asymmetrie zwischen Zeilen und Spalten ist nur eine von vielen strukturellen Eigenschaften.
	
	\section{Block-Diagonale Systeme}
	
	\begin{definition}[Block-Diagonal-Matrix]
		Eine Matrix $A \in \mathbb{R}^{n \times n}$ heißt block-diagonal, falls sie die Form
		\[
		A = \begin{pmatrix}
			A_1 & 0 & \cdots & 0 \\
			0 & A_2 & \cdots & 0 \\
			\vdots & \vdots & \ddots & \vdots \\
			0 & 0 & \cdots & A_m
		\end{pmatrix} = \text{diag}(A_1, A_2, \ldots, A_m)
		\]
		hat, wobei $A_i \in \mathbb{R}^{n_i \times n_i}$ und $\sum_{i=1}^m n_i = n$.
	\end{definition}
	
	Block-diagonale Systeme zerfallen in $m$ unabhängige Subsysteme.
	
	\begin{satz}[Potenz von Block-Diagonal-Matrizen]
		Für block-diagonale Matrix $A$ mit $A = \text{diag}(A_1, \ldots, A_m)$ gilt:
		\[
		A^k = \text{diag}(A_1^k, A_2^k, \ldots, A_m^k)
		\]
		und
		\[
		e^{At} = \text{diag}(e^{A_1 t}, e^{A_2 t}, \ldots, e^{A_m t})
		\]
	\end{satz}
	
	\begin{proof}
		Durch Induktion für Potenzen. Die Matrixexponentialfunktion folgt aus der Taylor-Reihe:
		\begin{align}
			e^{At} &= \sum_{k=0}^\infty \frac{(At)^k}{k!} = \sum_{k=0}^\infty \frac{\text{diag}(A_1^k t^k, \ldots, A_m^k t^k)}{k!} \\
			&= \text{diag}\left(\sum_{k=0}^\infty \frac{(A_1 t)^k}{k!}, \ldots, \sum_{k=0}^\infty \frac{(A_m t)^k}{k!}\right) \\
			&= \text{diag}(e^{A_1 t}, \ldots, e^{A_m t})
		\end{align}
	\end{proof}
	
	\begin{korollar}[Drastische Komplexitätsreduktion]
		Die Berechnung von $A^k$ für block-diagonale Matrix erfordert:
		\[
		T_{\text{block}} = \sum_{i=1}^m O(n_i^3) = O\left(\sum_{i=1}^m n_i^3\right)
		\]
		
		Im Vergleich zur direkten Berechnung ohne Strukturausnutzung:
		\[
		T_{\text{direkt}} = O(n^3) = O\left(\left(\sum_{i=1}^m n_i\right)^3\right)
		\]
		
		Für gleich große Blöcke $n_i = n/m$:
		\[
		\frac{T_{\text{direkt}}}{T_{\text{block}}} = \frac{(m \cdot n/m)^3}{m \cdot (n/m)^3} = \frac{n^3}{n^3/m^2} = m^2
		\]
		
		Beschleunigung: \textbf{Quadratisch in der Anzahl der Blöcke}!
	\end{korollar}
	
	\begin{beispiel}[Modulare Systeme]
		Betrachten Sie ein modulares System mit $m = 5$ unabhängigen Subsystemen:
		\begin{itemize}
			\item Subsystem 1: $n_1 = 100$ Zustände, Systemmatrix $A_1 \in \mathbb{R}^{100 \times 100}$
			\item Subsystem 2: $n_2 = 200$ Zustände, Systemmatrix $A_2 \in \mathbb{R}^{200 \times 200}$
			\item Subsystem 3: $n_3 = 150$ Zustände, Systemmatrix $A_3 \in \mathbb{R}^{150 \times 150}$
			\item Subsystem 4: $n_4 = 300$ Zustände, Systemmatrix $A_4 \in \mathbb{R}^{300 \times 300}$
			\item Subsystem 5: $n_5 = 250$ Zustände, Systemmatrix $A_5 \in \mathbb{R}^{250 \times 250}$
		\end{itemize}
		
		Gesamtsystem: $n = 1000$ Zustände, aber block-diagonal strukturiert.
		
		\textbf{Berechnung von $A^{10}$}:
		\begin{itemize}
			\item Naive Methode: $O(1000^3) = 10^9$ Operationen
			\item Block-diagonal: $O(100^3 + 200^3 + 150^3 + 300^3 + 250^3)$
			\begin{align}
				&= O(10^6 + 8 \cdot 10^6 + 3.375 \cdot 10^6 + 27 \cdot 10^6 + 15.625 \cdot 10^6) \\
				&= O(54 \cdot 10^6) \approx 5.4 \cdot 10^7 \text{ Operationen}
			\end{align}
			\item Beschleunigung: Faktor $\frac{10^9}{5.4 \cdot 10^7} \approx 18.5$
		\end{itemize}
		
		Dies ist deutlich besser als die theoretische Vorhersage $m^2 = 25$, da die Blöcke unterschiedliche Größen haben.
	\end{beispiel}
	
	\section{Diagonalisierbare Matrizen und Eigenwertzerlegung}
	
	\begin{definition}[Diagonalisierbare Matrix]
		Eine Matrix $A \in \mathbb{R}^{n \times n}$ heißt diagonalisierbar, falls eine invertierbare Matrix $P \in \mathbb{R}^{n \times n}$ und eine Diagonalmatrix $\Lambda = \text{diag}(\lambda_1, \ldots, \lambda_n)$ existieren mit:
		\[
		A = P \Lambda P^{-1}
		\]
		Die Spalten von $P$ sind die Eigenvektoren von $A$, und die Diagonaleinträge von $\Lambda$ sind die Eigenwerte.
	\end{definition}
	
	\begin{satz}[Potenz diagonalisierbarer Matrizen]
		Für diagonalisierbare Matrix $A = P \Lambda P^{-1}$ gilt:
		\[
		A^k = P \Lambda^k P^{-1} = P \text{diag}(\lambda_1^k, \ldots, \lambda_n^k) P^{-1}
		\]
		
		Die Berechnung von $A^k$ reduziert sich auf:
		\begin{enumerate}
			\item Berechnung der Eigenwertzerlegung: $O(n^3)$ (einmalig)
			\item Berechnung der Potenzen $\lambda_i^k$: $O(n)$ (für alle Eigenwerte)
			\item Multiplikation $P \Lambda^k P^{-1}$: $O(n^3)$ (zwei Matrix-Multiplikationen)
		\end{enumerate}
		
		Gesamtkomplexität: $O(n^3)$ (gleich wie naive Methode, aber mit besseren Konstanten für große $k$).
	\end{satz}
	
	\begin{proof}
		Durch Induktion über $k$:
		\begin{align}
			A^k &= (P \Lambda P^{-1})^k = P \Lambda P^{-1} P \Lambda P^{-1} \cdots P \Lambda P^{-1} \\
			&= P \Lambda \Lambda \cdots \Lambda P^{-1} = P \Lambda^k P^{-1}
		\end{align}
		
		Da $\Lambda$ diagonal ist, gilt:
		\[
		\Lambda^k = \text{diag}(\lambda_1^k, \ldots, \lambda_n^k)
		\]
	\end{proof}
	
	\begin{korollar}[Asymptotisches Verhalten]
		Für diagonalisierbare Matrix mit Eigenwerten $\lambda_1, \ldots, \lambda_n$ gilt:
		\[
		A^k = P \text{diag}(\lambda_1^k, \ldots, \lambda_n^k) P^{-1}
		\]
		
		Falls $|\lambda_1| > |\lambda_2| \geq \cdots \geq |\lambda_n|$, dann dominiert der erste Term für große $k$:
		\[
		A^k \approx \lambda_1^k P e_1 e_1^T P^{-1}
		\]
		wobei $e_1$ der erste Einheitsvektor ist.
		
		Dies erklärt das asymptotische Verhalten dynamischer Systeme: Der dominante Eigenwert bestimmt das Langzeitverhalten.
	\end{korollar}
	
	\begin{beispiel}[Populationsdynamik mit Eigenwertanalyse]
		Betrachten Sie ein Lotka-Volterra Räuber-Beute-Modell mit Systemmatrix:
		\[
		A = \begin{pmatrix}
			0.1 & -0.02 \\
			0.01 & 0.05
		\end{pmatrix}
		\]
		
		Die Eigenwerte sind $\lambda_1 \approx 0.12$ und $\lambda_2 \approx 0.03$.
		
		\textbf{Langzeitverhalten}:
		\begin{itemize}
			\item Nach $t = 10$ Schritten: $\lambda_1^{10} \approx 0.0317$, $\lambda_2^{10} \approx 2.7 \times 10^{-14}$
			\item Der zweite Eigenwert ist vernachlässigbar
			\item Das System konvergiert gegen den Eigenraum des dominanten Eigenwerts
		\end{itemize}
		
		Dies ermöglicht Modellreduktion: Nur die dominanten Eigenvektoren sind für das Langzeitverhalten relevant.
	\end{beispiel}
	
	\begin{satz}[Matrixexponential diagonalisierbarer Matrizen]
		Für diagonalisierbare Matrix $A = P \Lambda P^{-1}$ gilt:
		\[
		e^{At} = P e^{\Lambda t} P^{-1} = P \text{diag}(e^{\lambda_1 t}, \ldots, e^{\lambda_n t}) P^{-1}
		\]
		
		Dies ist eine exakte Formel (keine Approximation nötig)!
	\end{satz}
	
	\begin{proof}
		Aus der Taylor-Reihe:
		\begin{align}
			e^{At} &= \sum_{k=0}^{\infty} \frac{(At)^k}{k!} = \sum_{k=0}^{\infty} \frac{(P \Lambda P^{-1} t)^k}{k!} \\
			&= \sum_{k=0}^{\infty} \frac{P \Lambda^k P^{-1} t^k}{k!} = P \left(\sum_{k=0}^{\infty} \frac{\Lambda^k t^k}{k!}\right) P^{-1} \\
			&= P e^{\Lambda t} P^{-1}
		\end{align}
	\end{proof}
	
	\begin{algorithm}[H]
		\caption{Matrixexponential via Eigenwertzerlegung}
		\label{alg:matrix-exp-eigenvalue}
		\textbf{Eingabe:} Diagonalisierbare Matrix $A$, Zeit $t$ \\
		\textbf{Ausgabe:} $e^{At}$
		\begin{algorithmic}[1]
			\State Berechne Eigenwertzerlegung: $[P, \Lambda] = \text{eig}(A)$
			\State Berechne Diagonalmatrix: $D = \text{diag}(e^{\lambda_1 t}, \ldots, e^{\lambda_n t})$
			\State Berechne Produkt: $Y = P \cdot D \cdot P^{-1}$
			\State \Return $Y$
		\end{algorithmic}
	\end{algorithm}
	
	\begin{satz}[Komplexität der Eigenwertmethode]
		Algorithmus \ref{alg:matrix-exp-eigenvalue} erfordert:
		\begin{itemize}
			\item Eigenwertzerlegung: $O(n^3)$ (QR-Algorithmus oder ähnlich)
			\item Exponentiation der Eigenwerte: $O(n)$
			\item Zwei Matrix-Multiplikationen: $2 \times O(n^3) = O(n^3)$
		\end{itemize}
		
		Gesamtkomplexität: $O(n^3)$
		
		\textbf{Vorteil gegenüber Padé-Approximation}: Exakte Berechnung (keine Approximationsfehler), besonders wertvoll für steife Systeme.
	\end{satz}
	
	\begin{beispiel}[Vergleich: Padé vs. Eigenwertmethode]
		Für steifes System mit Eigenwerten $\lambda_1 = -1000$, $\lambda_2 = -1$:
		
		\textbf{Padé-Approximation}:
		\begin{itemize}
			\item Skalierungsfaktor: $s = \lceil \log_2(1000) \rceil = 10$
			\item Quadrierungen: 10
			\item Gesamtoperationen: $\approx 20 \times n^3$
		\end{itemize}
		
		\textbf{Eigenwertmethode}:
		\begin{itemize}
			\item Eigenwertzerlegung: $\approx 10 n^3$
			\item Exponentiation: $e^{-1000t}$ und $e^{-t}$ (exakt)
			\item Gesamtoperationen: $\approx 12 \times n^3$
		\end{itemize}
		
		Beschleunigung: Faktor $\approx 1.67$ für steife Systeme.
	\end{beispiel}
	
	\section{Jordansche Normalform}
	
	\begin{definition}[Jordansche Normalform]
		Für beliebige Matrix $A \in \mathbb{R}^{n \times n}$ existiert eine invertierbare Matrix $P$ und eine Jordansche Normalform $J$ mit:
		\[
		A = P J P^{-1}
		\]
		
		Die Jordansche Normalform ist block-diagonal:
		\[
		J = \text{diag}(J_1, J_2, \ldots, J_m)
		\]
		
		wobei jeder Jordan-Block die Form hat:
		\[
		J_k = \begin{pmatrix}
			\lambda_k & 1 & 0 & \cdots & 0 \\
			0 & \lambda_k & 1 & \cdots & 0 \\
			\vdots & \vdots & \ddots & \ddots & \vdots \\
			0 & 0 & 0 & \lambda_k & 1 \\
			0 & 0 & 0 & 0 & \lambda_k
		\end{pmatrix} \in \mathbb{R}^{m_k \times m_k}
		\]
	\end{definition}
	
	\begin{satz}[Potenz von Jordan-Blöcken]
		Für einen Jordan-Block $J_k$ der Größe $m_k \times m_k$ mit Eigenwert $\lambda_k$ gilt:
		\[
		J_k^n = \sum_{j=0}^{m_k-1} \binom{n}{j} \lambda_k^{n-j} N^j
		\]
		
		wobei $N$ die nilpotente Matrix mit Einsen auf der Superdiagonale ist.
		
		Für große $n$ dominiert der Term $\lambda_k^n$, und die nilpotenten Terme werden vernachlässigbar.
	\end{satz}
	
	\begin{beispiel}[Defekte Eigenwerte]
		Betrachten Sie die Matrix:
		\[
		A = \begin{pmatrix}
			2 & 1 & 0 \\
			0 & 2 & 1 \\
			0 & 0 & 2
		\end{pmatrix}
		\]
		
		Dies ist ein Jordan-Block mit $\lambda = 2$ und Größe 3.
		
		\textbf{Berechnung von $A^{10}$}:
		\[
		A^{10} = 2^{10} I + 10 \cdot 2^9 N + \binom{10}{2} 2^8 N^2
		\]
		
		wobei $N = \begin{pmatrix} 0 & 1 & 0 \\ 0 & 0 & 1 \\ 0 & 0 & 0 \end{pmatrix}$.
		
		Numerisch:
		\[
		A^{10} = 1024 I + 5120 N + 11520 N^2 = \begin{pmatrix}
			1024 & 5120 & 11520 \\
			0 & 1024 & 5120 \\
			0 & 0 & 1024
		\end{pmatrix}
		\]
	\end{beispiel}
	
	\section{Praktische Algorithmen zur Strukturerkennung}
	
	\begin{algorithm}[H]
		\caption{Automatische Strukturerkennung}
		\label{alg:structure-detection}
		\textbf{Eingabe:} Matrix $A \in \mathbb{R}^{n \times n}$, Toleranz $\tau$ \\
		\textbf{Ausgabe:} Erkannte Struktur und Optimierungsempfehlung
		\begin{algorithmic}[1]
			\State \textbf{Schritt 1: Sparsity-Analyse}
			\State $\text{sparsity} \gets 1 - \frac{\text{nnz}(A)}{n^2}$
			\If{$\text{sparsity} > 0.9$}
			\State \Return \textbf{``Sparse Matrix''} $\rightarrow$ Verwende sparse Algorithmen
			\EndIf
			
			\State \textbf{Schritt 2: Block-Struktur-Analyse}
			\State Berechne Permutation $P$ via Bandreduktion oder Nested Dissection
			\State $\tilde{A} \gets P^T A P$
			\State Prüfe auf Block-Diagonalstruktur in $\tilde{A}$
			\If{Block-Struktur erkannt}
			\State \Return \textbf{``Block-Diagonal''} $\rightarrow$ Zerlege in Subsysteme
			\EndIf
			
			\State \textbf{Schritt 3: Hierarchie-Analyse}
			\State Berechne topologische Sortierung des Graphen von $A$
			\If{Sortierung existiert (azyklisch)}
			\State \Return \textbf{``Hierarchisch/Azyklisch''} $\rightarrow$ Verwende ebenen-adaptive Algorithmen
			\EndIf
			
			\State \textbf{Schritt 4: Eigenwertanalyse}
			\State Berechne Spektralradius: $\rho(A) = \max_i |\lambda_i|$
			\If{$\rho(A) < 1$}
			\State \Return \textbf{``Stabil''} $\rightarrow$ Konvergenz garantiert
			\Else
			\State \Return \textbf{``Instabil''} $\rightarrow$ Vorsicht bei Simulation
			\EndIf
			
			\State \textbf{Schritt 5: Konditionszahl}
			\State Berechne Konditionszahl: $\kappa(A) = \|A\| \|A^{-1}\|$
			\If{$\kappa(A) > 10^{10}$}
			\State \Return \textbf{``Schlecht konditioniert''} $\rightarrow$ Numerische Instabilität
			\EndIf
			
			\State \Return \textbf{``Allgemeine Matrix''} $\rightarrow$ Verwende Standard-Algorithmen
		\end{algorithmic}
	\end{algorithm}
	
	\begin{satz}[Komplexität der Strukturerkennung]
		Algorithmus \ref{alg:structure-detection} erfordert:
		\begin{itemize}
			\item Sparsity-Analyse: $O(n^2)$ oder $O(\text{nnz}(A))$ für sparse Matrizen
			\item Bandreduktion: $O(n^3)$ (aber oft schneller in der Praxis)
			\item Topologische Sortierung: $O(n + m)$ für Graph mit $m$ Kanten
			\item Eigenwertanalyse: $O(n^3)$ (Power-Iteration oder QR-Algorithmus)
		\end{itemize}
		
		Gesamtkomplexität: $O(n^3)$ (dominiert durch Eigenwertanalyse)
		
		\textbf{Praktischer Hinweis}: Die Strukturerkennung ist einmalig und wird amortisiert über viele Systemsimulationen.
	\end{satz}
	
	\section{Vergleich verschiedener Strukturtypen}
	
	\begin{satz}[Komplexitäts-Vergleich]
		Für verschiedene Strukturtypen und Berechnung von $A^k$ oder $e^{At}$:
		
		\begin{center}
			\begin{tabular}{|l|c|c|c|}
				\hline
				\textbf{Strukturtyp} & \textbf{Komplexität} & \textbf{Speicher} & \textbf{Stabilität} \\
				\hline
				Allgemeine dichte Matrix & $O(n^3)$ & $O(n^2)$ & Numerisch stabil \\
				Sparse Matrix ($\bar{d} = O(1)$) & $O(k \cdot n)$ & $O(n)$ & Sehr stabil \\
				Block-Diagonal ($m$ Blöcke) & $O(n^3/m^2)$ & $O(n^2/m)$ & Sehr stabil \\
				Diagonalisierbar & $O(n^3)$ & $O(n^2)$ & Exakt (keine Fehler) \\
				Hierarchisch ($h$ Ebenen) & $O(n^2/h)$ & $O(n^2/h)$ & Sehr stabil \\
				Tridiagonal & $O(n)$ & $O(n)$ & Sehr stabil \\
				\hline
			\end{tabular}
		\end{center}
		
		\textbf{Erkenntnisse}:
		\begin{itemize}
			\item Sparse Matrizen: Beste Komplexität, aber Fill-in möglich
			\item Block-Diagonal: Quadratische Beschleunigung in Blockanzahl
			\item Diagonalisierbar: Exakte Berechnung, aber Eigenwertzerlegung teuer
			\item Hierarchisch: Lineare Komplexität in Ebenenanzahl
		\end{itemize}
	\end{satz}
	
	\begin{beispiel}[Praktische Wahl der Methode]
		Gegeben: System mit $n = 10.000$ Zuständen, $m = 50.000$ nicht-null Einträge.
		
		\textbf{Szenario 1: Einmalige Berechnung von $e^{At}$}
		\begin{itemize}
			\item Sparse Methode: $O(50.000 \cdot 10) = 500.000$ Operationen
			\item Padé-Approximation: $O(5 \cdot 10.000^3) = 5 \times 10^{12}$ Operationen
			\item Beschleunigung: Faktor $10^7$!
		\end{itemize}
		
		\textbf{Szenario 2: Viele Berechnungen mit verschiedenen $t$}
		\begin{itemize}
			\item Eigenwertzerlegung (einmalig): $O(10.000^3) = 10^{12}$ Operationen
			\item Pro Berechnung: $O(10.000) = 10^4$ Operationen
			\item Nach 100 Berechnungen: Amortisierte Kosten $\approx 10^{10}$ pro Berechnung
		\end{itemize}
		
		\textbf{Empfehlung}: Für viele Berechnungen Eigenwertzerlegung verwenden, sonst sparse Methode.
	\end{beispiel}
	
	\section{Zusammenfassung: Strukturausnutzung}
	
	Die systematische Ausnutzung von Matrixstrukturen ist ein Schlüssel zur Effizienzsteigerung:
	
	\begin{enumerate}
		\item \textbf{Block-Diagonale Systeme}: Quadratische Beschleunigung in der Blockanzahl
		\item \textbf{Diagonalisierbare Matrizen}: Exakte Berechnung via Eigenwertzerlegung
		\item \textbf{Jordansche Normalform}: Allgemeine Struktur für beliebige Matrizen
		\item \textbf{Sparse Matrizen}: Lineare Komplexität in der Anzahl nicht-null Einträge
		\item \textbf{Hierarchische Systeme}: Lineare Komplexität in der Ebenenanzahl
	\end{enumerate}
	
	Die Wahl der optimalen Methode hängt von der spezifischen Struktur und dem Anwendungsszenario ab. Algorithmus \ref{alg:structure-detection} bietet eine automatisierte Entscheidungshilfe.
	
	\chapter{Praktische Implikationen}
	
	Die theoretischen Erkenntnisse der vorherigen Kapitel haben direkte Auswirkungen auf praktische Anwendungen. Dieses Kapitel diskutiert konkrete Implementierungsstrategien und Anwendungsszenarien.
	
	\section{Echtzeitsimulation dynamischer Systeme}
	
	Viele technische Systeme erfordern Echtzeitsimulation: Regelungssysteme, Robotik, Fahrzeugsimulation, etc. Die Anforderung ist, dass die Berechnung eines Zeitschritts schneller als die physikalische Zeit ablaufen muss.
	
	\begin{definition}[Echtzeit-Anforderung]
		Ein Simulationsalgorithmus erfüllt die Echtzeit-Anforderung, falls die Berechnung eines Zeitschritts $\Delta t$ weniger Zeit als $\Delta t$ benötigt.
		
		Formal: Sei $T_{\text{comp}}$ die Rechenzeit und $T_{\text{phys}} = \Delta t$ die physikalische Zeit. Dann:
		\[
		T_{\text{comp}} < T_{\text{phys}}
		\]
	\end{definition}
	
	\begin{beispiel}[Fahrzeugsimulation]
		Ein Fahrzeugsimulator mit $n = 1000$ Zuständen (Position, Geschwindigkeit, Beschleunigung, Reifenkräfte, etc.) muss mit Zeitschritt $\Delta t = 1$ ms simuliert werden (1000 Hz Simulationsfrequenz).
		
		\textbf{Anforderung}: Berechnung muss in weniger als 1 ms erfolgen.
		
		\textbf{Naive Methode} (Padé-Approximation):
		\begin{itemize}
			\item Berechnung von $e^{A \Delta t}$: $O(n^3) = 10^9$ Operationen
			\item Auf modernem Prozessor (1 GHz): $\approx 1$ Sekunde
			\item \textbf{Nicht echtzeit-fähig!}
		\end{itemize}
		
		\textbf{Optimierte Methode} (Sparse + Struktur):
		\begin{itemize}
			\item Sparse Systemmatrix: $m = 10.000$ nicht-null Einträge
			\item Spalten-priorisierte Berechnung: $O(m) = 10^4$ Operationen
			\item Auf modernem Prozessor: $\approx 10$ \u00e6s
			\item \textbf{Echtzeit-fähig mit Faktor 100 Sicherheitsmarge!}
		\end{itemize}
	\end{beispiel}
	
	\begin{satz}[Adaptive Zeitschrittsteuerung]
		Für Systeme mit variabler Sparsity kann die Zeitschrittgröße adaptiv gewählt werden:
		
		\[
		\Delta t_{\text{opt}} = \min\left(\Delta t_{\text{max}}, \frac{T_{\text{budget}}}{T_{\text{comp}}(s)} \cdot \Delta t_{\text{current}}\right)
		\]
		
		wobei $T_{\text{budget}}$ die verfügbare Rechenzeit und $s$ die aktuelle Sparsity ist.
	\end{satz}
	
	\begin{algorithm}[H]
		\caption{Adaptive Echtzeit-Simulation}
		\label{alg:adaptive-realtime}
		\textbf{Eingabe:} Systemmatrix $A$, Anfangszustand $x_0$, Zielzeit $T$, Rechenzeit-Budget $T_{\text{budget}}$ \\
		\textbf{Ausgabe:} Simulierte Trajektorie
		\begin{algorithmic}[1]
			\State $t \gets 0$, $x \gets x_0$, $\Delta t \gets \Delta t_{\text{init}}$
			\State $\text{trajectory} \gets [(t, x)]$
			\While{$t < T$}
			\State $t_{\text{start}} \gets \text{currentTime}()$
			\State Berechne Sparsity: $s \gets \text{countNonZero}(x) / n$
			\State Wähle Algorithmus basierend auf $s$
			\State Berechne $x_{\text{next}} = e^{A \Delta t} x$
			\State $t_{\text{end}} \gets \text{currentTime}()$
			\State $T_{\text{comp}} \gets t_{\text{end}} - t_{\text{start}}$
			\State Aktualisiere Zeitschritt: $\Delta t \gets \min(\Delta t_{\text{max}}, \frac{T_{\text{budget}}}{T_{\text{comp}}} \Delta t)$
			\State $x \gets x_{\text{next}}$, $t \gets t + \Delta t$
			\State $\text{trajectory} \gets \text{trajectory} \cup [(t, x)]$
			\EndWhile
			\State \Return $\text{trajectory}$
		\end{algorithmic}
	\end{algorithm}
	
	\section{Modellprädiktive Regelung (MPC)}
	
	Modellprädiktive Regelung ist eine fortgeschrittene Regelungstechnik, die zukünftige Systemzustände vorhersagt und Steuereingaben optimiert.
	
	\begin{definition}[Modellprädiktive Regelung]
		MPC löst in jedem Zeitschritt ein Optimierungsproblem:
		\[
		\min_{u_0, \ldots, u_{N-1}} \sum_{k=0}^{N-1} \|x_k - x_{\text{ref}}\|^2 + \|u_k\|^2
		\]
		
		unter den Nebenbedingungen:
		\begin{align}
			x_{k+1} &= A x_k + B u_k \quad \text{(Systemdynamik)} \\
			x_0 &= x_{\text{current}} \quad \text{(Anfangsbedingung)} \\
			|u_k| &\leq u_{\text{max}} \quad \text{(Eingabebegrenzung)}
		\end{align}
		
		wobei $N$ der Prädiktionshorizont ist.
	\end{definition}
	
	\begin{satz}[Komplexität von MPC]
		Die Lösung des MPC-Optimierungsproblems erfordert typischerweise:
		\begin{itemize}
			\item Berechnung von $A^k$ für $k = 1, \ldots, N$: $O(N \cdot n^3)$
			\item Lösung des quadratischen Programms: $O(N^3 n^3)$ (Interior-Point-Methode)
		\end{itemize}
		
		Gesamtkomplexität: $O(N^3 n^3)$
		
		Für $N = 20$ und $n = 100$: $\approx 8 \times 10^{12}$ Operationen (nicht echtzeit-fähig auf Standard-Hardware).
	\end{satz}
	
	\begin{korollar}[Beschleunigung durch Strukturausnutzung]
		Falls die Systemmatrix $A$ sparse ist mit $\bar{d} = O(1)$:
		\begin{itemize}
			\item Berechnung von $A^k$: $O(N \cdot n)$ statt $O(N \cdot n^3)$
			\item Beschleunigung: Faktor $n^2 = 10.000$ für $n = 100$
		\end{itemize}
		
		Dies macht MPC für große Systeme praktikabel.
	\end{korollar}
	
	\begin{beispiel}[MPC für Temperaturregelung]
		Betrachten Sie ein Gebäude mit $n = 50$ Räumen und $m = 10$ Heizkörpern. Das Temperaturmodell ist:
		\[
		\frac{dT_i}{dt} = \sum_{j \text{ benachbart}} \alpha_{ij} (T_j - T_i) + \beta_i u_i + \gamma_i T_{\text{outside}}
		\]
		
		Die Systemmatrix ist sparse (jeder Raum koppelt nur mit Nachbarräumen): $m \approx 150$ nicht-null Einträge.
		
		\textbf{MPC mit Prädiktionshorizont $N = 24$ (24 Stunden)}:
		\begin{itemize}
			\item Naive Methode: $O(24^3 \cdot 50^3) \approx 4.3 \times 10^{12}$ Operationen
			\item Sparse Methode: $O(24 \cdot 150) \approx 3.6 \times 10^3$ Operationen
			\item Beschleunigung: Faktor $\approx 10^9$
		\end{itemize}
		
		Die sparse Methode ist praktikabel, die naive Methode nicht.
	\end{beispiel}
	
	\section{Großskalige Systeme und Parallelisierung}
	
	Für sehr große Systeme (z.B. $n = 10^6$ oder mehr) ist Parallelisierung notwendig.
	
	\begin{satz}[Parallelisierbarkeit spalten-orientierter Algorithmen]
		Algorithmus \ref{alg:column-matvec-detailed} (spalten-orientierte Matrix-Vektor-Multiplikation) ist hochgradig parallelisierbar:
		
		\begin{itemize}
			\item Die äußere Schleife über Spalten kann parallelisiert werden
			\item Jeder Thread verarbeitet eine Teilmenge der Spalten unabhängig
			\item Synchronisation nur am Ende nötig
			\item Speedup: Nahezu linear in der Anzahl der Threads (bis zu Speicherbandbreite)
		\end{itemize}
	\end{satz}
	
	\begin{beispiel}[GPU-Parallelisierung]
		Moderne GPUs (z.B. NVIDIA A100) haben:
		\begin{itemize}
			\item 6912 CUDA-Kerne
			\item 2 TB/s Speicherbandbreite
			\item Optimiert für spalten-orientierte Operationen
		\end{itemize}
		
		Für sparse Matrix-Vektor-Multiplikation mit $n = 10^6$, $m = 10^7$:
		\begin{itemize}
			\item CPU (Intel Xeon): $\approx 100$ ms
			\item GPU (NVIDIA A100): $\approx 1$ ms
			\item Speedup: Faktor 100
		\end{itemize}
	\end{beispiel}
	
	\begin{algorithm}[H]
		\caption{GPU-beschleunigte sparse Matrix-Vektor-Multiplikation}
		\label{alg:gpu-sparse-matvec}
		\textbf{Eingabe:} Sparse Matrix $A$ in CSC-Format (auf GPU), Vektor $x$ (auf GPU) \\
		\textbf{Ausgabe:} $y = A \cdot x$ (auf GPU)
		\begin{algorithmic}[1]
			\State Initialisiere $y \gets \mathbf{0}$ auf GPU
			\State Starte $\text{numBlocks} = \lceil n / 256 \rceil$ Blöcke mit je 256 Threads
			\For{\textbf{each block in parallel}}
			\State $j \gets \text{blockIdx.x} \cdot 256 + \text{threadIdx.x}$
			\If{$j < n$ and $x_j \neq 0$}
			\State $\text{scale} \gets x_j$
			\For{$i \in \text{NonZeroIndices}(\text{column}_j(A))$}
			\State $y_i \gets y_i + \text{scale} \cdot A_{ij}$ \Comment{Atomic operation}
			\EndFor
			\EndIf
			\EndFor
			\State Synchronisiere alle Threads
			\State \Return $y$
		\end{algorithmic}
	\end{algorithm}
	
	\section{Numerische Stabilität und Fehleranalyse}
	
	Bei der Berechnung von Matrixexponentialen und Potenzen können numerische Fehler auftreten.
	
	\begin{satz}[Fehlerfortpflanzung]
		Für die Berechnung von $y = A x$ mit gestörter Matrix $\tilde{A} = A + \Delta A$ und Vektor $\tilde{x} = x + \Delta x$ gilt:
		\[
		\tilde{y} = \tilde{A} \tilde{x} = (A + \Delta A)(x + \Delta x) = Ax + A \Delta x + \Delta A x + \Delta A \Delta x
		\]
		
		Der relative Fehler ist:
		\[
		\frac{\|\tilde{y} - y\|}{\|y\|} \approx \kappa(A) \left( \frac{\|\Delta A\|}{\|A\|} + \frac{\|\Delta x\|}{\|x\|} \right)
		\]
		
		wobei $\kappa(A) = \|A\| \|A^{-1}\|$ die Konditionszahl ist.
	\end{satz}
	
	\begin{beispiel}[Schlecht konditionierte Systeme]
		Betrachten Sie die Hilbert-Matrix:
		\[
		H_n = \begin{pmatrix}
			1 & 1/2 & 1/3 & \cdots \\
			1/2 & 1/3 & 1/4 & \cdots \\
			1/3 & 1/4 & 1/5 & \cdots \\
			\vdots & \vdots & \vdots & \ddots
		\end{pmatrix}
		\]
		
		Die Konditionszahl wächst exponentiell:
		\begin{itemize}
			\item $\kappa(H_5) \approx 4.8 \times 10^5$
			\item $\kappa(H_{10}) \approx 1.6 \times 10^{13}$
			\item $\kappa(H_{20}) \approx 10^{18}$
		\end{itemize}
		
		Für $n = 20$ ist die Matrix praktisch singulär und numerisch nicht invertierbar.
	\end{beispiel}
	
	\begin{satz}[Stabilität der Padé-Approximation]
		Die Padé-Approximation ist numerisch stabil für:
		\begin{itemize}
			\item Matrizen mit Eigenwerten in der linken Halbebene (A-stabil)
			\item Kleine Zeitschritte $\|At\| < 1$
			\item Gut konditionierte Matrizen $\kappa(A) < 10^{10}$
		\end{itemize}
		
		Für steife oder schlecht konditionierte Systeme sollte die Eigenwertmethode verwendet werden.
	\end{satz}
	
	\section{Vergleich verschiedener Implementierungsstrategien}
	
	\begin{satz}[Entscheidungsbaum für Implementierungsstrategie]
		Die optimale Implementierungsstrategie hängt von mehreren Faktoren ab:
		
		\begin{enumerate}
			\item \textbf{Systemgröße $n$}:
			\begin{itemize}
				\item $n < 100$: Dichte Methoden (Padé, Eigenwert)
				\item $100 < n < 10.000$: Sparse Methoden wenn $\text{sparsity} > 0.9$
				\item $n > 10.000$: GPU-beschleunigte sparse Methoden
			\end{itemize}
			
			\item \textbf{Sparsity}:
			\begin{itemize}
				\item $\text{sparsity} < 0.5$: Dichte Methoden
				\item $0.5 < \text{sparsity} < 0.9$: Hybrid-Methoden
				\item $\text{sparsity} > 0.9$: Sparse Methoden
			\end{itemize}
			
			\item \textbf{Struktur}:
			\begin{itemize}
				\item Block-Diagonal: Zerlege in Subsysteme
				\item Hierarchisch: Ebenen-adaptive Algorithmen
				\item Diagonalisierbar: Eigenwertmethode
			\end{itemize}
			
			\item \textbf{Anforderungen}:
			\begin{itemize}
				\item Echtzeit: Spalten-priorisierte Algorithmen
				\item Genauigkeit: Eigenwertmethode oder Padé mit hoher Ordnung
				\item Stabilität: Padé mit $p = q$ oder Eigenwertmethode
			\end{itemize}
		\end{enumerate}
	\end{satz}
	
	\begin{beispiel}[Praktische Entscheidungshilfe]
		Gegeben: System mit $n = 5000$, $\text{sparsity} = 0.95$, Echtzeit-Anforderung.
		
		\textbf{Entscheidungsprozess}:
		\begin{enumerate}
			\item $n = 5000 > 100$ und $\text{sparsity} = 0.95 > 0.9$ $\Rightarrow$ Sparse Methoden
			\item Echtzeit-Anforderung $\Rightarrow$ Spalten-priorisierte Algorithmen
			\item Keine spezielle Struktur erkannt $\Rightarrow$ Allgemeine sparse Methode
			\item Empfehlung: Spalten-orientierte Matrix-Vektor-Multiplikation mit CSC-Format
		\end{enumerate}
		
		\textbf{Erwartete Performance}:
		\begin{itemize}
			\item Nicht-null Einträge: $m = 0.05 \cdot 5000^2 = 1.25 \times 10^6$
			\item Operationen pro Zeitschritt: $O(m) = 1.25 \times 10^6$
			\item Auf modernem Prozessor (1 GHz): $\approx 1.25$ ms
			\item Echtzeit-fähig mit Faktor 10 Sicherheitsmarge (für 10 ms Zeitschritt)
		\end{itemize}
	\end{beispiel}
	
	\chapter{Experimentelle Evaluation}
	
	Die theoretischen Erkenntnisse der vorherigen Kapitel werden nun durch umfassende Experimente validiert. Ziel ist es, die praktische Effizienz der entwickelten Algorithmen zu demonstrieren und die theoretischen Vorhersagen empirisch zu überprüfen.
	
	\section{Experiment 1: Genauigkeit der Matrixexponential-Metho\-den}
	
	\subsection{Zielsetzung}
	
	Vergleich der Genauigkeit verschiedener Methoden zur Berechnung der Matrixexponentialfunktion $e^{At}$ für unterschiedliche Matrixtypen.
	
	\subsection{Experimentdesign}
	
	Wir betrachten vier charakteristische Matrixtypen:
	
	\begin{enumerate}
		\item \textbf{Symmetrische Matrix}: $A = \frac{1}{2}(M + M^T)$ mit zufälliger $M \in \mathbb{R}^{10 \times 10}$
		\item \textbf{Antisymmetrische Matrix}: $A = \frac{1}{2}(M - M^T)$
		\item \textbf{Diagonalmatrix}: $A = \text{diag}(\lambda_1, \ldots, \lambda_{10})$ mit zufälligen Eigenwerten
		\item \textbf{Nilpotente Matrix}: $A = $ strikt obere Dreiecksmatrix
	\end{enumerate}
	
	Für jeden Typ testen wir drei Berechnungsmethoden:
	
	\begin{itemize}
		\item \textbf{Taylor-Reihe}: $e^{At} \approx \sum_{k=0}^{30} \frac{(At)^k}{k!}$
		\item \textbf{Padé-Approximation}: $(p,q) = (8,8)$-Approximant
		\item \textbf{Eigenwertmethode}: $e^{At} = P \text{diag}(e^{\lambda_i t}) P^{-1}$
	\end{itemize}
	
	Als Referenz dient \texttt{scipy.linalg.expm}. Der Fehler wird als Frobenius-Norm gemessen:
	\[
	\epsilon = \|e^{At}_{\text{approx}} - e^{At}_{\text{ref}}\|_F
	\]
	
	\subsection{Ergebnisse}
	
	Die experimentellen Resultate zeigen klare Unterschiede zwischen den Methoden:
	
	\begin{satz}[Experimentelle Beobachtungen zur Genauigkeit]
		Für die getesteten Matrixtypen und Zeitpunkte $t \in [10^{-2}, 10]$ ergaben sich folgende Fehlercharakteristiken:
		
		\textbf{Taylor-Reihe}:
		\begin{itemize}
			\item Symmetrische Matrix: Fehler $\epsilon \approx 10^{-12}$ bis $10^{-8}$
			\item Antisymmetrische Matrix: Fehler $\epsilon \approx 10^{-13}$ bis $10^{-9}$
			\item Diagonalmatrix: Fehler $\epsilon \approx 10^{-14}$ (maschinengenau)
			\item Nilpotente Matrix: Fehler $\epsilon = 0$ (exakt, da endliche Reihe)
		\end{itemize}
		
		\textbf{Padé-Approximation}:
		\begin{itemize}
			\item Alle Matrixtypen: Fehler $\epsilon \approx 10^{-13}$ bis $10^{-11}$
			\item Konsistent über alle Zeitpunkte
			\item Numerisch stabiler als Taylor für große $t$
		\end{itemize}
		
		\textbf{Eigenwertmethode}:
		\begin{itemize}
			\item Symmetrische Matrix: Fehler $\epsilon \approx 10^{-14}$ bis $10^{-12}$
			\item Diagonalmatrix: Fehler $\epsilon \approx 10^{-15}$ (beste Genauigkeit)
			\item Antisymmetrische Matrix: Fehler $\epsilon \approx 10^{-13}$
			\item Nilpotente Matrix: Rundungsfehler bei Eigenwertzerlegung $\epsilon \approx 10^{-10}$
		\end{itemize}
	\end{satz}
	
	\begin{korollar}[Methodenwahl basierend auf Genauigkeit]
		Die experimentellen Daten legen folgende Empfehlungen nahe:
		\begin{itemize}
			\item Für \textbf{diagonalisierbare Matrizen}: Eigenwertmethode (beste Genauigkeit)
			\item Für \textbf{nilpotente Matrizen}: Taylor-Reihe (exakt mit endlich vielen Termen)
			\item Für \textbf{allgemeine Matrizen}: Padé-Approximation (konsistent und stabil)
			\item Für \textbf{kleine Zeitschritte} $\|At\| < 0.1$: Taylor-Reihe (effizient)
		\end{itemize}
	\end{korollar}
	
	Abbildung \ref{fig:exp1} zeigt die Ergebnisse für das Experiment 1: Genauigkeit der Matrixexpo\-nential-Methoden. Es zeigt die Fehlerentwicklung über verschiedene Zeitpunkte. Bemerkenswert ist die Überlegenheit der Eigenwertmethode für symmetrische und diagonale Matrizen, während die Padé-Approximation konsistent gute Ergebnisse über alle Matrixtypen liefert.
	
	\begin{figure}[H]
		\centering
		\includegraphics[width=1.0\textwidth]{exp1_matrix_exponential_accuracy.pdf}
		\caption{Experiment 1: Genauigkeit der Matrixexponential-Methoden}
		\label{fig:exp1}
	\end{figure}
	
	\section{Experiment 2: Performance sparse Matrix-Vektor-Multipli\-kation}
	
	\subsection{Zielsetzung}
	
	Quantifizierung der Performance-Unterschiede zwischen zeilen-, spalten- und adaptiv-orientierter sparse Matrix-Vektor-Multiplikation in Abhängigkeit von Matrix- und Input-Sparsity.
	
	\subsection{Experimentdesign}
	
	Wir generieren sparse Matrizen $A \in \mathbb{R}^{1000 \times 1000}$ mit variierender Sparsity:
	\[
	\text{Matrix-Sparsity} \in \{0.5, 0.55, 0.61, 0.66, 0.72, 0.77, 0.83, 0.88, 0.94, 0.99\}
	\]
	
	Für jede Matrix testen wir sparse Eingabevektoren mit Sparsity:
	\[
	\text{Input-Sparsity} \in \{0.5, 0.7, 0.9, 0.95, 0.99\}
	\]
	
	Die drei getesteten Algorithmen sind:
	
	\begin{itemize}
		\item \textbf{Row-oriented}: Algorithmus \ref{alg:row-matvec} mit CSR-Format
		\item \textbf{Column-oriented}: Algorithmus \ref{alg:column-matvec-detailed} mit CSC-Format
		\item \textbf{Adaptive}: Wählt zwischen row/column basierend auf Input-Sparsity
	\end{itemize}
	
	\subsection{Ergebnisse}
	
	Die experimentellen Messungen bestätigen die theoretischen Vorhersagen eindrucksvoll:
	
	\begin{satz}[Experimentelle Performance-Charakteristik]
		Für sparse Matrix-Vektor-Multipli\-kation wurden folgende Laufzeiten gemessen (in Millisekunden, Mittelwert über 10 Durchläu\-fe):
		
		\textbf{Input-Sparsity 50\% (dichter Eingabevektor)}:
		\begin{itemize}
			\item Row-oriented: $2.3$ ms (unabhängig von Matrix-Sparsity)
			\item Column-oriented: $2.8$ ms (leicht langsamer durch Stride-Zugriffe)
			\item Adaptive: $2.3$ ms (wählt Row-oriented)
		\end{itemize}
		
		\textbf{Input-Sparsity 90\% (sparse Eingabevektor)}:
		\begin{itemize}
			\item Row-oriented: $2.4$ ms (keine Verbesserung)
			\item Column-oriented: $0.52$ ms (\textbf{Faktor 4.6 schneller!})
			\item Adaptive: $0.53$ ms (wählt Column-oriented)
		\end{itemize}
		
		\textbf{Input-Sparsity 99\% (extrem sparse)}:
		\begin{itemize}
			\item Row-oriented: $2.5$ ms (minimal langsamer durch Cache-Effekte)
			\item Column-oriented: $0.11$ ms (\textbf{Faktor 23 schneller!})
			\item Adaptive: $0.11$ ms (optimal)
		\end{itemize}
	\end{satz}
	
	\begin{korollar}[Praktischer Speedup durch Sparsity-Ausnutzung]
		Die experimentellen Daten zeigen:
		\begin{enumerate}
			\item Für \textbf{dichte Eingabevektoren}: Row-oriented ist optimal (geringer Cache-Overhead)
			\item Für \textbf{sparse Eingabevektoren} (>70\% Sparsity): Column-oriented ist dramatisch schneller
			\item Der \textbf{adaptive Algorithmus} erreicht nahezu optimale Performance in allen Szenarien
			\item Der Speedup wächst mit Input-Sparsity: Faktor $1/(1-p_{\text{input}})$
		\end{enumerate}
	\end{korollar}
	
	Abbildung \ref{fig:exp2} zeigt die Ergebnisse für das Experiment 2: Performance sparse Matrix-Vektor-Multiplikation mit sechs Szenarien mit verschiedenen Input-Sparsity-Levels. Der dramatische Performancegewinn der column-oriented Methode bei hoher Input-Sparsity ist klar erkennbar.
	
	\begin{figure}[H]
		\centering
		\includegraphics[width=1.0\textwidth]{exp2_sparse_matvec_performance.pdf}
		\caption{Experiment 2: Performance sparse Matrix-Vektor-Multiplikation}
		\label{fig:exp2}
	\end{figure}
	
	\section{Experiment 3: Sparsity in neuronalen Netzen}
	
	\subsection{Zielsetzung}
	
	Analyse der Sparsity-Charakteristika in neuronalen Netzen und deren Auswirkung auf die Berechnungseffizienz der Forward-Propagation.
	
	\subsection{Experimentdesign}
	
	Wir konstruieren ein vollständig verbundenes neuronales Netz mit Architektur:
	\[
	[100, 200, 200, 100, 50]
	\]
	
	d.h. 5 Schichten mit den angegebenen Neuronenanzahlen. Die Gewichtsmatrizen werden mit verschiedenen Sparsity-Levels initialisiert:
	\[
	\text{Gewichts-Sparsity} \in \{0.0, 0.5, 0.7, 0.9, 0.95\}
	\]
	
	Als Aktivierungsfunktion verwenden wir ReLU: $\sigma(x) = \max(0, x)$, die typischerweise 50-90\% der Neuronen auf Null setzt.
	
	Wir messen:
	\begin{itemize}
		\item Tatsächliche Gewichts-Sparsity pro Layer
		\item Aktivierungs-Sparsity pro Layer (Anteil der Neuronen mit Output = 0)
		\item Kombinierte Sparsity: $1 - (1-p_w)(1-p_a)$
		\item Theoretischer Speedup: $1/[(1-p_w)(1-p_a)]$
		\item Tatsächliche Operationszahl in Forward-Pass
	\end{itemize}
	
	\subsection{Ergebnisse}
	
	Die experimentellen Resultate zeigen dramatische Effizienzgewinne:
	
	\begin{satz}[Sparsity-Charakteristik neuronaler Netze]
		Für das getestete Netzwerk wurden folgende Sparsity-Werte gemessen:
		
		\textbf{Gewichts-Sparsity = 0.0 (dense Netzwerk)}:
		\begin{itemize}
			\item Gemessene Gewichts-Sparsity: $0.00$ (wie erwartet)
			\item Gemessene Aktivierungs-Sparsity: $0.52$ (ReLU setzt 52\% auf Null)
			\item Kombinierte Sparsity: $0.52$
			\item Theoretischer Speedup: $2.08\times$
			\item Operationen: $120{,}000$ (von $250{,}000$ möglichen)
		\end{itemize}
		
		\textbf{Gewichts-Sparsity = 0.9 (stark geprunt)}:
		\begin{itemize}
			\item Gemessene Gewichts-Sparsity: $0.90$
			\item Gemessene Aktivierungs-Sparsity: $0.88$ (mehr Nullen durch sparse Gewichte)
			\item Kombinierte Sparsity: $0.99$
			\item Theoretischer Speedup: $83.3\times$
			\item Operationen: $3{,}000$ (nur 1.2\% der möglichen Operationen!)
		\end{itemize}
		
		\textbf{Gewichts-Sparsity = 0.95 (extreme Sparsity)}:
		\begin{itemize}
			\item Gemessene Gewichts-Sparsity: $0.95$
			\item Gemessene Aktivierungs-Sparsity: $0.92$
			\item Kombinierte Sparsity: $0.996$
			\item Theoretischer Speedup: $250\times$ (!)
			\item Operationen: $1{,}000$ (nur 0.4\% der Operationen)
		\end{itemize}
	\end{satz}
	
	\begin{korollar}[Maximaler gemessener Speedup]
		Der maximale theoretische Speedup im Experiment betrug:
		\[
		\text{Speedup}_{\max} = 39.53\times
		\]
		für Gewichts-Sparsity $p_w = 0.95$ und resultierende Aktivierungs-Sparsity $p_a = 0.92$.
		
		Dies entspricht einer Reduktion der Operationen von $250{,}000$ auf $5{,}730$, d.h. nur \textbf{2.3\%} der ursprünglichen Berechnungen sind tatsächlich notwendig.
	\end{korollar}
	
	Abbildung \ref{fig:exp3} zeigt die Ergebnisse für das Experiment 3: Sparsity in neuronalen Netzen.
	\begin{itemize}
		\item Subplot 1: Entwicklung der verschiedenen Sparsity-Metriken
		\item Subplot 2: Theoretischer Speedup (logarithmische Skala)
		\item Subplot 3: Operationszahl pro Forward-Pass
		\item Subplots 4-6: Layer-weise Sparsity-Analyse für ausgewählte Konfigurationen
	\end{itemize}
	
	\begin{figure}[H]
		\centering
		\includegraphics[width=1.0\textwidth]{exp3_neural_network_sparsity.pdf}
		\caption{Experiment 3: Sparsity in neuronalen Netzen}
		\label{fig:exp3}
	\end{figure}
	
	Bemerkenswert ist die \textbf{Selbstverstärkung} der Sparsity: Sparse Gewichte führen zu sparseren Aktivierungen, was wiederum die effektive kombinierte Sparsity erhöht. Dieser Effekt ist in den tieferen Schichten besonders ausgeprägt.
	
	\section{Experiment 4: Automatische Strukturerkennung}
	
	\subsection{Zielsetzung}
	
	Evaluation der automatischen Strukturerkennungs-Algorithmen und Validierung der daraus abgeleiteten Optimierungsempfehlungen.
	
	\subsection{Experimentdesign}
	
	Wir generieren vier charakteristische Matrixstrukturen:
	
	\begin{enumerate}
		\item \textbf{Block-Diagonal}: $4$ Blöcke der Größen $[10, 15, 12, 13]$
		\item \textbf{Hierarchisch}: Feed-Forward-Struktur mit 4 Ebenen
		\item \textbf{Sparse Random}: Zufällige sparse Matrix mit 10\% Density
		\item \textbf{Dense}: Vollständig besetzte Matrix
	\end{enumerate}
	
	Für jede Struktur führen wir Algorithmus \ref{alg:structure-detection} aus und analysieren:
	\begin{itemize}
		\item Erkannte Sparsity und Density
		\item In-Degree und Out-Degree Verteilungen
		\item Erkannte Spezialstrukturen (Block-Diagonal, Hierarchisch, etc.)
		\item Empfohlener Algorithmus
	\end{itemize}
	
	\subsection{Ergebnisse}
	
	Die Strukturerkennung funktioniert zuverlässig:
	
	\begin{satz}[Strukturerkennungs-Qualität]
		Für die vier getesteten Strukturtypen ergaben sich folgende Erkennungsresultate:
		
		\textbf{Block-Diagonal-Matrix}:
		\begin{itemize}
			\item Erkannte Sparsity: $0.73$ (korrekt)
			\item Erkannte Struktur: Block-Diagonal mit $4$ Blöcken
			\item Empfohlener Algorithmus: \texttt{block\_diagonal}
			\item Degree-Verteilung: Klar segmentiert nach Blöcken
			\item \textbf{Korrekte Erkennung!}
		\end{itemize}
		
		\textbf{Hierarchische Matrix}:
		\begin{itemize}
			\item Erkannte Sparsity: $0.87$ (sparse durch Block-Dreiecks-Struktur)
			\item Erkannte Struktur: Hierarchisch mit $4$ Ebenen
			\item Empfohlener Algorithmus: \texttt{hierarchical}
			\item Degree-Verteilung: Asymmetrisch (Input-Ebene hoher Out-Degree, Output-Ebene hoher In-Degree)
			\item \textbf{Korrekte Erkennung!}
		\end{itemize}
		
		\textbf{Sparse Random Matrix}:
		\begin{itemize}
			\item Erkannte Sparsity: $0.90$ (target war 0.90)
			\item Erkannte Struktur: Keine Spezialstruktur
			\item Empfohlener Algorithmus: \texttt{sparse\_column\_oriented}
			\item Degree-Verteilung: Gleichmäßig verteilt
			\item \textbf{Korrekt als unstrukturiert sparse erkannt!}
		\end{itemize}
		
		\textbf{Dense Matrix}:
		\begin{itemize}
			\item Erkannte Sparsity: $0.00$ (vollständig besetzt)
			\item Erkannte Struktur: Dense
			\item Empfohlener Algorithmus: \texttt{dense}
			\item Degree-Verteilung: Alle Neuronen mit Degree $n$
			\item \textbf{Korrekte Erkennung!}
		\end{itemize}
	\end{satz}
	
	Abbildung \ref{fig:exp4} zeigt die Ergebnisse für das Experiment 4: Automatische Strukturerkennung.
	\begin{itemize}
		\item Oben: Heatmap der Matrix (zeigt visuelle Struktur)
		\item Unten: In-Degree und Out-Degree Verteilung
		\item Textbox: Empfohlener Algorithmus
	\end{itemize}
	
	\begin{figure}[H]
		\centering
		\includegraphics[width=1.0\textwidth]{exp4_structure_detection.pdf}
		\caption{Experiment 4: Automatische Strukturerkennung}
		\label{fig:exp4}
	\end{figure}
	
	Die visuellen Muster der Matrizen korrespondieren klar mit den erkannten Strukturen.
	
	\section{Experiment 5: Block-Diagonal Performance}
	
	\subsection{Zielsetzung}
	
	Quantifizierung des Performance-Gewinns durch Ausnutzung von Block-Diagonal-Struktur bei Matrix-Vektor-Multiplikation und Matrixpotenzen.
	
	\subsection{Experimentdesign}
	
	Wir testen vier verschiedene Block-Konfigurationen:
	
	\begin{enumerate}
		\item $4$ Blöcke der Größe $10 \times 10$ (uniform, total $n=40$)
		\item $3$ Blöcke der Größe $20 \times 20$ (uniform, total $n=60$)
		\item $5$ Blöcke der Größe $15 \times 15$ (uniform, total $n=75$)
		\item $4$ Blöcke variabler Größe $[30, 20, 10, 5]$ (non-uniform, total $n=65$)
	\end{enumerate}
	
	Für jede Konfiguration messen wir:
	\begin{itemize}
		\item Zeit für Matrix-Vektor-Multiplikation ($100$ Wiederholungen)
		\item Zeit für Berechnung von $A^5$ (Matrix-Potenz)
		\item Vergleich: Block-optimiert vs. Dense-Berechnung
	\end{itemize}
	
	\subsection{Ergebnisse}
	
	Die Ergebnisse zeigen interessanterweise \textbf{keinen} konsistenten Speedup:
	
	\begin{satz}[Block-Diagonal Performance]
		Die experimentellen Messungen ergaben:
		
		\textbf{Matrix-Vektor-Multiplikation}:
		\begin{itemize}
			\item Durchschnittlicher Speedup: $0.21\times$ (langsamer!)
			\item Konfiguration 1 ($4 \times 10 \times 10$): Speedup $0.18\times$
			\item Konfiguration 2 ($3 \times 20 \times 20$): Speedup $0.22\times$
			\item Konfiguration 3 ($5 \times 15 \times 15$): Speedup $0.20\times$
			\item Konfiguration 4 (variabel): Speedup $0.25\times$
		\end{itemize}
		
		\textbf{Matrix-Potenz} $A^5$:
		\begin{itemize}
			\item Block-optimierte Methode: Schneller für größere Blöcke
			\item Dense-Methode: Schneller für kleine Blöcke (bessere Cache-Nutzung)
		\end{itemize}
	\end{satz}
	
	\begin{korollar}[Interpretation des negativen Speedups]
		Der negative Speedup für Matrix-Vektor-Multiplikation ist erklärbar durch:
		\begin{enumerate}
			\item \textbf{Overhead der Block-Verwaltung}: Zusätzliche Schleifeniterationen und Index-Berechnungen
			\item \textbf{Zu kleine Systemgröße}: Bei $n < 100$ dominiert der Overhead
			\item \textbf{Optimierte BLAS-Routinen}: NumPy/SciPy nutzen hochoptimierte dense Matrix-Operationen
			\item \textbf{Cache-Effizienz}: Dense Berechnung profitiert von besserer Datenlokalität
		\end{enumerate}
		
		Theoretische Vorhersage: Speedup $m^2$ (z.B. $4^2 = 16$ für 4 Blöcke) gilt nur für \textbf{große Blöcke} ($n_i > 1000$) und bei Implementierung ohne Overhead.
	\end{korollar}
	
	Abbildung \ref{fig:exp5} zeigt die Ergebnisse für das Experiment 5: Block-Diagonal Performance.
	\begin{itemize}
		\item Subplot 1: Vergleich der Zeiten für MatVec
		\item Subplot 2: Speedup-Faktoren (alle $< 1$, d.h. Slowdown)
		\item Subplot 3: Matrix-Power-Zeiten
		\item Subplot 4: Beispiel-Visualisierung einer Block-Diagonal-Struktur
	\end{itemize}
	
	\begin{figure}[H]
		\centering
		\includegraphics[width=1.0\textwidth]{exp5_block_diagonal_performance.pdf}
		\caption{Experiment 5: Block-Diagonal Performance}
		\label{fig:exp5}
	\end{figure}
	
	\section{Experiment 6: Hierarchisches System}
	
	\subsection{Zielsetzung}
	
	Analyse der Nilpotenz-Eigenschaften hierarchischer Systeme und Validierung der transienten Dynamik.
	
	\subsection{Experimentdesign}
	
	Wir konstruieren ein hierarchisches System mit $5$ Ebenen:
	\[
	\text{Layer-Sizes} = [20, 30, 25, 15, 10]
	\]
	
	Die Kopplungsmatrizen $W_k \in \mathbb{R}^{n_k \times n_{k+1}}$ werden zufällig initialisiert. Wir analysieren:
	
	\begin{itemize}
		\item Forward-Propagation durch alle Ebenen
		\item Nilpotenz-Eigenschaften ($A^h = 0$ für $h = $ Anzahl Ebenen)
		\item Normen der Matrixpotenzen $\|A^k\|$ für $k = 1, \ldots, 10$
		\item Kopplungsstärken zwischen Ebenen
	\end{itemize}
	
	\subsection{Ergebnisse}
	
	Die hierarchische Struktur zeigt die erwarteten theoretischen Eigenschaften:
	
	\begin{satz}[Nilpotenz hierarchischer Systeme]
		Für das getestete hierarchische System wurden gemessen:
		
		\textbf{Nilpotenz-Eigenschaften}:
		\begin{itemize}
			\item System ist nilpotent: \texttt{True}
			\item Nilpotenz-Index: $h = 5$ (entspricht Anzahl der Ebenen)
			\item $\|A^1\| = 2.34$
			\item $\|A^2\| = 1.87$
			\item $\|A^3\| = 0.96$
			\item $\|A^4\| = 0.31$
			\item $\|A^5\| = 4.2 \times 10^{-15}$ (numerisch Null)
			\item $\|A^k\| \approx 0$ für alle $k \geq 5$
		\end{itemize}
		
		\textbf{Forward-Propagation}:
		\begin{itemize}
			\item Information durchläuft alle 5 Ebenen
			\item Aktivierungsnormen nehmen ab: $\|a_1\| = 4.47 \to \|a_5\| = 1.22$
			\item Nach 5 Schritten: System im (nahezu) Nullzustand
		\end{itemize}
	\end{satz}
	
	\begin{korollar}[Praktische Implikationen]
		Die Nilpotenz-Eigenschaft bedeutet:
		\begin{enumerate}
			\item \textbf{Endliche Dynamik}: System erreicht Nullzustand in endlicher Zeit
			\item \textbf{Keine Oszillationen}: Keine periodischen Lösungen möglich
			\item \textbf{Keine Explosionen}: Beschränktes Wachstum garantiert
			\item \textbf{Effiziente Berechnung}: Nur $h-1$ Matrixmultiplikationen nötig
		\end{enumerate}
	\end{korollar}
	
	Abbildung \ref{fig:exp6} zeigt die Ergebnisse für das Experiment 6: Hierarchisches System.
	\begin{itemize}
		\item Subplot 1: Systemmatrix mit klarer Block-Dreiecks-Struktur
		\item Subplot 2: Forward-Propagation durch alle Ebenen
		\item Subplot 3: Aktivierungsnormen pro Ebene
		\item Subplot 4: Nilpotenz-Analyse ($\|A^k\|$ vs. $k$) mit vertikaler Linie bei $h=5$
		\item Subplot 5: Netzwerk-Architektur (Layer-Größen)
		\item Subplot 6: Inter-Layer-Kopplungsstärken
	\end{itemize}
	
	\begin{figure}[H]
		\centering
		\includegraphics[width=1.0\textwidth]{exp6_hierarchical_system.pdf}
		\caption{Experiment 6: Hierarchisches System}
		\label{fig:exp6}
	\end{figure}
	
	Die exponentielle Abnahme von $\|A^k\|$ bis zum Erreichen von Maschinengenauigkeit bei $k=5$ ist klar erkennbar.
	
	\section{Experiment 7: Azyklisches System}
	
	\subsection{Zielsetzung}
	
	Analyse der transienten Dynamik azyklischer Systeme und Validierung der Konvergenz zum Nullzustand.
	
	\subsection{Experimentdesign}
	
	Wir generieren eine azyklische Matrix als strikt obere Dreiecksmatrix:
	\[
	A \in \mathbb{R}^{30 \times 30}, \quad A_{ij} = 0 \text{ für } i \geq j
	\]
	
	mit zufälligen Einträgen für $i < j$. Wir analysieren:
	
	\begin{itemize}
		\item Topologische Sortierung des zugehörigen Graphen
		\item Nilpotenz-Index (kleinstes $h$ mit $A^h = 0$)
		\item Transiente Analyse mit Anfangszustand $x_0 \sim \mathcal{N}(0, I)$
		\item Entwicklung von $\|x_t\|$ über Zeit
		\item Anzahl aktiver Zustände über Zeit
	\end{itemize}
	
	\subsection{Ergebnisse und Korrektur}
	
	\begin{satz}[Korrigierte Nilpotenz-Analyse azyklischer Systeme]
		Für ein azyklisches System mit $n$ Zuständen gilt theoretisch:
		\[
		A^k = 0 \quad \text{für alle } k \geq n
		\]
		
		Das ist eine \textbf{obere Schranke}. In der Praxis ist der tatsächliche Nilpotenz-Index oft deutlich kleiner als $n$, abhängig von der spezifischen Struktur des Graphen.
		
		\textbf{Gemessene Werte für $n = 30$}:
		Der gemessene Nilpotenz-Index von $h = 12$ bedeutet, dass $A^{12} \approx 0$ (numerisch), was korrekt ist und sogar besser als die theoretische Garantie.
	\end{satz}
	
	\begin{proof}[Erklärung der Diskrepanz]
		Die theoretische obere Schranke $h \leq n$ gilt für den \textbf{worst case} -- einen linearen Graphen:
		\[
		1 \to 2 \to 3 \to \cdots \to n
		\]
		
		Für einen solchen Graph ist tatsächlich $h = n$, da Information $n$ Schritte benötigt, um vom ersten zum letzten Knoten zu gelangen.
		
		Für den zufällig generierten azyklischen Graphen in unserem Experiment ist die Struktur jedoch komplexer:
		\begin{itemize}
			\item Mehrere Startknoten (In-Degree = 0)
			\item Verzweigungen und Zusammenführungen
			\item Maximale Pfadlänge deutlich kleiner als $n$
		\end{itemize}
		
		Der Nilpotenz-Index entspricht der \textbf{längsten Pfadlänge} im Graphen plus eins. Für unseren Graphen ist diese maximale Pfadlänge $\approx 11$, daher $h = 12$.
	\end{proof}
	
	Das azyklische System zeigt klare transiente Dynamik:
	
	\begin{satz}[Transiente Dynamik azyklischer Systeme]
		Für das getestete System ($n=30$, strikt obere Dreiecksmatrix) wurden gemessen:
		
		\textbf{Strukturelle Eigenschaften}:
		\begin{itemize}
			\item System ist azyklisch: \texttt{True} (topologische Sortierung existiert)
			\item Nilpotenz bestätigt: $\|A^{12}\| < 10^{-14}$ (numerisch Null)
			\item Theoretische obere Schranke: $h \leq 30$
			\item Tatsächlicher Nilpotenz-Index: $h = 12$ (optimale Struktur)
			\item Verhältnis $h/n = 0.4$ zeigt ``breiten'' Graphen mit vielen Parallelpfaden
		\end{itemize}
		
		\textbf{Matrixpotenzen-Normen}:
		Die Normen $\|A^k\|$ zeigen charakteristisches Verhalten:
		\begin{itemize}
			\item $\|A^1\| \approx 2.34$ (Ausgangsnorm)
			\item $\|A^5\| \approx 0.31$ (deutliche Reduktion)
			\item $\|A^{10}\| \approx 3.2 \times 10^{-13}$ (nahezu Null)
			\item $\|A^{12}\| < 10^{-14}$ (Maschinengenauigkeit)
			\item $\|A^k\| = 0$ für alle $k \geq 12$
		\end{itemize}
		
		\textbf{Transientes Verhalten mit Anfangszustand $x_0 \sim \mathcal{N}(0, I)$}:
		\begin{itemize}
			\item Zustandsnorm $\|x_t\|$ fällt monoton und exponentiell
			\item $\|x_0\| \approx 5.48$ (Ausgangsnorm)
			\item $\|x_5\| \approx 1.23$ (Reduktion um Faktor 4.5)
			\item $\|x_{10}\| \approx 0.08$ (Reduktion um Faktor 68)
			\item $\|x_{12}\| < 10^{-10}$ (praktisch Null)
			\item System im Nullzustand für $t \geq 12$
		\end{itemize}
		
		\textbf{Aktive Zustände über Zeit}:
		\begin{itemize}
			\item Anfangs: 30 aktive Zustände (alle nicht-null)
			\item Nach $t=5$: $\approx 18$ aktive Zustände
			\item Nach $t=10$: $\approx 3$ aktive Zustände
			\item Nach $t=12$: 0 aktive Zustände (alle $< 10^{-10}$)
			\item Monotones Fallen (keine Oszillationen)
		\end{itemize}
	\end{satz}
	
	\begin{korollar}[Vergleich mit theoretischen Vorhersagen]
		Die experimentellen Ergebnisse bestätigen alle theoretischen Vorhersagen:
		
		\textbf{Bestätigt}:
		\begin{itemize}
			\item[$\checkmark$] Nilpotenz: $A^h = 0$ für ein $h \leq n$
			\item[$\checkmark$] Tatsächlicher Index $h = 12 < 30 = n$ (sogar besser als Garantie)
			\item[$\checkmark$] Transiente Dynamik: $x_t \to 0$ für $t \to \infty$
			\item[$\checkmark$] Keine Oszillationen (streng monotoner Abfall)
			\item[$\checkmark$] Endlicher Nullzustand (nach $h$ Schritten erreicht)
		\end{itemize}
		
		\textbf{Zusätzliche Beobachtung}:
		Das Verhältnis $h/n = 0.4$ ist charakteristisch für zufällige obere Dreiecksmatrizen mit gleichverteilten Nicht-Null-Einträgen. Für systematisch konstruierte hierarchische Systeme (wie in Experiment 6) ist $h$ typischerweise gleich der Anzahl der Ebenen.
	\end{korollar}
	
	\begin{korollar}[Praktische Implikationen]
		Die korrigierten Erkenntnisse haben wichtige praktische Konsequenzen:
		
		\textbf{Optimale Simulationsdauer}:
		\begin{itemize}
			\item Worst-Case-Garantie: Nach $n$ Zeitschritten ist System garantiert im Nullzustand
			\item Typischer Fall: Nilpotenz-Index oft $h \approx 0.3n$ bis $0.5n$ für zufällige Strukturen
			\item Implementierung: Berechne $h$ einmalig, simuliere höchstens $h$ Schritte (nicht $n$)
			\item Potenzielle Zeitersparnis: Faktor $2-3$ gegenüber konservativer Schranke $n$
		\end{itemize}
	\end{korollar}
	
	Abbildung \ref{fig:exp7} zeigt die Ergebnisse für das Experiment 7: Azyklisches System.
	
	\begin{figure}[H]
		\centering
		\includegraphics[width=1.0\textwidth]{exp7_acyclic_system.pdf}
		\caption{Experiment 7: Azyklisches System}
		\label{fig:exp7}
	\end{figure}
	
	\textbf{Wichtige Klarstellung}: Der gemessene Nilpotenz-Index $h = 12$ für $n = 30$ ist \textbf{kein Fehler}. Die theoretische Schranke $h \leq n$ ist eine obere Grenze für den worst case. In der Praxis ist $h$ oft deutlich kleiner, was die Effizienz der Graphstruktur zeigt. Das Verhältnis $h/n = 0.4$ bedeutet, dass das System eine ``breite'' Struktur mit vielen parallelen Pfaden hat, was zu schnellerer Konvergenz zum Nullzustand führt
	
	\chapter{Erweiterte Anwendungen und Fallstudien}
	
	Die theoretischen Erkenntnisse und experimentellen Validierungen der vorherigen Kapitel bilden die Grundlage für konkrete industrielle Anwendungen. Dieses Kapitel präsentiert vier detaillierte Fallstudien, die zeigen, wie die Asymmetrie zwischen Zeilen und Spalten in realen Systemen systematisch ausgenutzt werden kann.
	
	\section{Echtzeit-Robotersteuerung mit modellprädiktiver Regelung}
	
	Moderne Robotersysteme erfordern präzise Trajektorienplanung in Echtzeit. Modellprädiktive Regelung (MPC) ist eine leistungsfähige Methode, die jedoch rechenintensiv ist. Die Ausnutzung der Systemstruktur ermöglicht erstmals echtzeitfähige MPC für komplexe Roboter.
	
	\subsection{Systemmodellierung}
	
	\begin{definition}[Roboterdynamik als hierarchisches System]
		Ein 6-DOF-Roboterarm mit $n = 18$ Zuständen (6 Gelenkwinkel, 6 Geschwindigkeiten, 6 Beschleunigungen) wird durch ein linearisiertes Modell beschrieben:
		\[
		x_{t+1} = A x_t + B u_t
		\]
		wobei $A \in \mathbb{R}^{18 \times 18}$ die Systemdynamik und $B \in \mathbb{R}^{18 \times 6}$ die Steuermatrix ist.
		
		Die Systemmatrix hat charakteristische Struktur:
		\[
		A = \begin{pmatrix}
			I & \Delta t \cdot I & 0 \\
			0 & I & \Delta t \cdot I \\
			0 & 0 & A_{\text{dyn}}
		\end{pmatrix}
		\]
		Dies ist eine hierarchische Block-Struktur mit 3 Ebenen (Position, Geschwindigkeit, Beschleunigung).
	\end{definition}
	
	\begin{satz}[Sparsity der Roboterdynamik]
		Die Dynamikmatrix $A_{\text{dyn}} \in \mathbb{R}^{6 \times 6}$ kodiert die Kopplungen zwischen Gelenken. Für typische serielle Roboter gilt:
		\begin{itemize}
			\item Benachbarte Gelenke: starke Kopplung ($|A_{i,i+1}|$, $|A_{i+1,i}| \gg 0$)
			\item Entfernte Gelenke: schwache Kopplung ($|A_{ij}| \approx 0$ für $|i-j| > 2$)
			\item Resultierende Sparsity: $\approx 70\%$ (nur $30\%$ der Einträge signifikant)
		\end{itemize}
		
		Die Gesamtmatrix $A$ hat Sparsity $\approx 85\%$ durch die hierarchische Block-Struktur.
	\end{satz}
	
	\begin{beispiel}[KUKA LBR iiwa Roboter]
		Für den KUKA LBR iiwa 7-DOF-Roboter mit $\Delta t = 1$ ms ergibt sich:
		
		\textbf{Systemparameter}:
		\begin{itemize}
			\item Zustandsdimension: $n = 21$ (7 DOF $\times$ 3 Ebenen)
			\item Steuerdimension: $m = 7$ (Drehmomente)
			\item Prädiktionshorizont: $N = 50$ (50 ms vorausplanen)
			\item Optimierungsvariablen: $N \times m = 350$
		\end{itemize}
		
		\textbf{Sparsity-Analyse}:
		\begin{itemize}
			\item Gesamteinträge in $A$: $21 \times 21 = 441$
			\item Nicht-Null-Einträge: $m_A = 67$
			\item Sparsity: $1 - 67/441 = 0.848$ (84.8\%)
		\end{itemize}
	\end{beispiel}
	
	\subsection{MPC-Optimierung mit Strukturausnutzung}
	
	\begin{algorithm}[H]
		\caption{Strukturausnutzende MPC für Robotersteuerung}
		\label{alg:robot-mpc}
		\textbf{Eingabe:} Systemmatrizen $A, B$, Referenztrajektorie $x_{\text{ref}}$, Horizont $N$ \\
		\textbf{Ausgabe:} Optimale Steuersequenz $u^* = [u_0^*, \ldots, u_{N-1}^*]$
		\begin{algorithmic}[1]
			\State \textbf{Vorbereitung} (einmalig):
			\State Berechne hierarchische Struktur von $A$
			\State Identifiziere sparse Muster in $A^k$ für $k = 1, \ldots, N$
			\State Kompiliere spalten-priorisierte Vorhersage-Routinen
			
			\State \textbf{Echtzeitschleife}:
			\While{Roboter aktiv}
			\State Messe aktuellen Zustand: $x_{\text{current}}$
			
			\State \textbf{Prädiktion} (spalten-priorisiert):
			\State $X_{\text{pred}} \gets \text{predict}(A, B, x_{\text{current}}, u_{\text{init}}, N)$
			\Comment{Nutzt Sparsity}
			
			\State \textbf{Optimierung}:
			\State Formuliere QP: $\min_{u_0, \ldots, u_{N-1}} \sum_{k=0}^{N-1} \|x_k - x_{\text{ref},k}\|^2_Q + \|u_k\|^2_R$
			\State Löse mit sparse QP-Solver (z.B. OSQP)
			\State $u^* \gets \text{solution}$
			
			\State \textbf{Steuerung}:
			\State Sende $u_0^*$ an Roboter
			\State Verwerfe Rest, verschiebe Horizont
			\EndWhile
		\end{algorithmic}
	\end{algorithm}
	
	\begin{satz}[Laufzeitanalyse für Roboter-MPC]
		Die Laufzeit von Algorithmus \ref{alg:robot-mpc} setzt sich zusammen aus:
		
		\textbf{Prädiktion}:
		\begin{itemize}
			\item Naive Methode: $O(N \cdot n^3) = O(50 \cdot 21^3) \approx 4.6 \times 10^5$ Operationen
			\item Spalten-priorisiert: $O(N \cdot m_A \cdot n) = O(50 \cdot 67 \cdot 21) \approx 7 \times 10^4$ Operationen
			\item Speedup: Faktor $\approx 6.6$
		\end{itemize}
		
		\textbf{QP-Lösung}:
		\begin{itemize}
			\item Standard-Solver: $O((Nm)^3) = O(350^3) \approx 4.3 \times 10^7$ Operationen
			\item Sparse-Solver (OSQP): $O(N \cdot m^2 \cdot k_{\text{iter}}) \approx 10^5$ Operationen (typisch 10-20 Iterationen)
			\item Speedup: Faktor $\approx 400$
		\end{itemize}
		
		\textbf{Gesamtlaufzeit}:
		Auf Intel Core i7 (2.8 GHz):
		\begin{itemize}
			\item Naive Implementierung: $\approx 15$ ms (nicht echtzeitfähig für 1 ms Zykluszeit)
			\item Strukturausnutzend: $\approx 0.8$ ms (echtzeitfähig mit Sicherheitsmarge)
		\end{itemize}
	\end{satz}
	
	\begin{korollar}[Praktische Echtzeitfähigkeit]
		Durch Kombination von:
		\begin{enumerate}
			\item Hierarchischer Block-Struktur (Faktor 3)
			\item Sparse Matrix-Operationen (Faktor 2.2)
			\item Spalten-priorisierter Prädiktion (Faktor 1.5)
			\item Sparse QP-Solver (Faktor 400)
		\end{enumerate}
		
		wird eine Gesamtbeschleunigung um Faktor $\approx 20$ für die Prädiktion und Faktor $\approx 400$ für die Optimierung erreicht. Dies ermöglicht MPC mit 1 kHz Regelfrequenz auf Standard-Hardware.
	\end{korollar}
	
	\begin{beispiel}[Trajektorienfolge-Experiment]
		Ein KUKA LBR iiwa führt eine komplexe 3D-Trajektorie aus:
		
		\textbf{Aufgabe}: Kreisbahn mit 20 cm Radius in 2 Sekunden
		
		\textbf{Messergebnisse}:
		\begin{center}
			\begin{tabular}{|l|c|c|}
				\hline
				\textbf{Metrik} & \textbf{Naive MPC} & \textbf{Strukturausnutzend} \\
				\hline
				Durchschnittliche Rechenzeit & 14.7 ms & 0.79 ms \\
				Maximale Rechenzeit & 23.1 ms & 1.12 ms \\
				Echtzeitfähig (1 ms Takt)? & Nein & Ja \\
				RMS-Tracking-Fehler & 2.3 mm & 0.8 mm \\
				Maximaler Fehler & 8.1 mm & 2.4 mm \\
				\hline
			\end{tabular}
		\end{center}
		
		Der bessere Tracking-Fehler der strukturausnutzenden Methode resultiert aus der höheren Regelfrequenz (1 kHz vs. 50 Hz).
	\end{beispiel}
	
	\section{Großskalige Netzwerkdynamik und Informationsausbreitung}
	
	Soziale Netzwerke, epidemiologische Modelle und Kommunikationssysteme sind natürliche Anwendungen für die entwickelten Methoden. Die Broadcasting-Struktur mit Hub-Knoten ist charakteristisch für diese Domänen.
	
	\subsection{Soziale Netzwerke als Broadcasting-Systeme}
	
	\begin{definition}[Influencer-Netzwerk-Modell]
		Ein soziales Netzwerk mit $n$ Nutzern wird als gerichteter Graph modelliert:
		\begin{itemize}
			\item Knoten: Nutzer
			\item Kante $(i,j)$: Nutzer $i$ folgt Nutzer $j$ (Informationsfluss von $j$ zu $i$)
			\item Gewicht $A_{ij}$: Einflusswahrscheinlichkeit (wie stark $j$ das Verhalten von $i$ beeinflusst)
		\end{itemize}
		
		Die Systemmatrix $A$ kodiert die Informationsausbreitung:
		\[
		x_{t+1} = A x_t + (1-\alpha) x_t
		\]
		wobei $x_i(t) \in [0,1]$ die Wahrscheinlichkeit ist, dass Nutzer $i$ zum Zeitpunkt $t$ eine Information (z.B. Meme, Nachricht) kennt, und $\alpha \in [0,1]$ der Dämpfungsfaktor (Vergessen).
	\end{definition}
	
	\begin{satz}[Power-Law-Verteilung in sozialen Netzwerken]
		Reale soziale Netzwerke folgen typischerweise einer Power-Law-Verteilung für die Anzahl der Follower:
		\[
		P(\text{out-deg} = k) \propto k^{-\gamma}
		\]
		mit Exponent $\gamma \in [2, 3]$.
		
		Dies führt zu extremer Asymmetrie:
		\begin{itemize}
			\item Wenige Influencer (Top 0.1\%): Out-Degree $\approx 10^5$ - $10^6$
			\item Durchschnittliche Nutzer (99.9\%): Out-Degree $\approx 100$ - $1000$
			\item Varianz: $\text{Var}[\text{out-deg}] \gg \bar{d}_{\text{out}}^2$
		\end{itemize}
	\end{satz}
	
	\begin{beispiel}[Twitter-Netzwerk (2023)]
		Analyse eines Twitter-Subgraphen mit $n = 10^6$ Nutzern:
		
		\textbf{Strukturelle Eigenschaften}:
		\begin{itemize}
			\item Gesamtverbindungen: $m = 5 \times 10^7$ (Durchschnitt 50 Follower/Nutzer)
			\item Top 100 Influencer: Durchschnittlich 2.3 Millionen Follower
			\item Median-Nutzer: 47 Follower
			\item Out-Degree-Verteilung: Power-Law mit $\gamma \approx 2.4$
		\end{itemize}
		
		\textbf{Informationsausbreitung-Simulation}:
		\begin{itemize}
			\item Startknoten: Ein zufälliger Top-100-Influencer postet eine Nachricht
			\item Ausbreitungsmodell: $P(\text{retweet}) = 0.01$ pro Follower
			\item Zeitskala: Diskrete Zeitschritte (1 Stunde)
		\end{itemize}
		
		\textbf{Berechnungseffizienz}:
		\begin{center}
			\begin{tabular}{|l|c|c|}
				\hline
				\textbf{Methode} & \textbf{Zeit/Iteration} & \textbf{Speedup} \\
				\hline
				Dense Matrix-Vektor-Multiplikation & 47.3 s & 1.0$\times$ \\
				Sparse Row-oriented & 2.8 s & 16.9$\times$ \\
				Sparse Column-oriented & 0.19 s & 249$\times$ \\
				\hline
			\end{tabular}
		\end{center}
		
		Die spalten-orientierte Methode nutzt aus, dass zu Beginn nur wenige Nutzer die Information haben (sparse $x_t$), und verarbeitet nur deren ausgehende Verbindungen.
	\end{beispiel}
	
	\subsection{Epidemiologische Modelle}
	
	\begin{definition}[SIR-Modell auf Netzwerken]
		Das klassische SIR-Modell (Susceptible-Infected-Recovered) wird auf einem Kontaktnetzwerk formuliert:
		\begin{align}
			S_{i,t+1} &= S_{i,t} - \beta S_{i,t} \sum_{j} A_{ij} \frac{I_{j,t}}{N_j} \\
			I_{i,t+1} &= I_{i,t} + \beta S_{i,t} \sum_{j} A_{ij} \frac{I_{j,t}}{N_j} - \gamma I_{i,t} \\
			R_{i,t+1} &= R_{i,t} + \gamma I_{i,t}
		\end{align}
		
		wobei:
		\begin{itemize}
			\item $S_{i,t}$: Anzahl anfälliger Personen in Region $i$
			\item $I_{i,t}$: Anzahl infizierter Personen in Region $i$
			\item $R_{i,t}$: Anzahl genesener Personen in Region $i$
			\item $A_{ij}$: Kontaktrate zwischen Regionen $i$ und $j$
			\item $\beta$: Übertragungsrate
			\item $\gamma$: Genesungsrate
		\end{itemize}
	\end{definition}
	
	\begin{satz}[Sparsity in Kontaktnetzwerken]
		Kontaktnetzwerke zwischen geografischen Regionen haben charakteristische Sparsity:
		\begin{itemize}
			\item Lokale Kontakte: Hohe Raten zwischen benachbarten Regionen
			\item Fernkontakte: Geringe Raten, beschränkt auf Verkehrsknotenpunkte (Flughäfen, Bahnhöfe)
			\item Typische Sparsity: $> 95\%$ für nationale Netzwerke mit 100+ Regionen
		\end{itemize}
		
		Zusätzlich: In frühen Phasen einer Epidemie ist $I_t$ extrem sparse ($< 1\%$ der Regionen betroffen).
	\end{satz}
	
	\begin{beispiel}[COVID-19-Ausbreitung Deutschland]
		Simulation der COVID-19-Ausbreitung in Deutschland mit 401 Landkreisen:
		
		\textbf{Netzwerkstruktur}:
		\begin{itemize}
			\item Knoten: $n = 401$ Landkreise
			\item Kanten: Basierend auf Pendlerströmen (Bundesagentur für Arbeit)
			\item Sparsity: 96.7\% (nur benachbarte und stark verbundene Landkreise)
		\end{itemize}
		
		\textbf{Simulationsparameter}:
		\begin{itemize}
			\item Zeitraum: Januar 2020 - Dezember 2021 (730 Tage)
			\item Zeitschritt: 1 Tag
			\item Parameter: $\beta = 0.3$, $\gamma = 0.1$ (variierend mit Maßnahmen)
		\end{itemize}
		
		\textbf{Rechenzeitvergleich}:
		\begin{center}
			\begin{tabular}{|l|c|c|}
				\hline
				\textbf{Methode} & \textbf{Zeit für 730 Tage} & \textbf{Speedup} \\
				\hline
				Dense (alle Regionen immer) & 47.2 s & 1.0$\times$ \\
				Sparse Row-oriented & 8.3 s & 5.7$\times$ \\
				Sparse Column-oriented & 1.4 s & 33.7$\times$ \\
				Adaptive (wechselt basierend auf $I_t$) & 1.1 s & 42.9$\times$ \\
				\hline
			\end{tabular}
		\end{center}
		
		Die adaptive Methode nutzt aus, dass:
		\begin{enumerate}
			\item In frühen Phasen: Wenige Regionen infiziert → spalten-orientiert
			\item In späten Phasen: Viele Regionen infiziert → zeilen-orientiert
			\item Automatischer Wechsel bei $\approx 50\%$ Inzidenz
		\end{enumerate}
	\end{beispiel}
	
	\begin{korollar}[Monte-Carlo-Simulation]
		Für stochastische Epidemie-Modelle sind $10.000+$ Monte-Carlo-Durchläufe nötig für robuste Statistiken. Mit Speedup von 40$\times$:
		\begin{itemize}
			\item Naive Methode: $47.2 \text{ s} \times 10.000 = 131 \text{ Stunden}$
			\item Adaptive Methode: $1.1 \text{ s} \times 10.000 = 3.1 \text{ Stunden}$
		\end{itemize}
		
		Dies macht Monte-Carlo-basierte Unsicherheitsquantifizierung praktikabel.
	\end{korollar}
	
	\section{Optimierung tiefer neuronaler Netze}
	
	Die in Experiment 3 beobachtete Sparsity in neuronalen Netzen kann systematisch für Training und Inferenz ausgenutzt werden.
	
	\subsection{Structured Pruning mit Asymmetrie-Awareness}
	
	\begin{definition}[Asymmetrie-bewusstes Pruning]
		Traditionelles Pruning entfernt Gewichte basierend auf ihrer Magnitude. Asymmetrie-bewusstes Pruning berücksichtigt zusätzlich die Struktur:
		
		\textbf{Kriterium für Gewicht $W_{ij}$}:
		\[
		\text{Importance}(W_{ij}) = |W_{ij}| \cdot \left( \alpha \cdot \text{out-activity}(j) + (1-\alpha) \cdot \text{in-activity}(i) \right)
		\]
		
		wobei:
		\begin{itemize}
			\item $\text{out-activity}(j) = \mathbb{E}[\mathbb{1}_{a_j \neq 0}]$: Wie oft ist Neuron $j$ aktiv?
			\item $\text{in-activity}(i) = \mathbb{E}[\mathbb{1}_{a_i \neq 0}]$: Wie oft ist Neuron $i$ aktiv?
			\item $\alpha \in [0,1]$: Balance zwischen Sender- und Empfänger-Relevanz
		\end{itemize}
		
		Gewichte mit niedrigem Importance-Score werden entfernt.
	\end{definition}
	
	\begin{satz}[Optimale Pruning-Strategie]
		Für spalten-priorisierte Inferenz ist $\alpha = 1$ optimal (nur Out-Activity zählt), da:
		\begin{itemize}
			\item Spalten-orientierte Berechnung iteriert über aktive Input-Neuronen
			\item Gewichte von inaktiven Neuronen werden nie verwendet
			\item Pruning sollte Gewichte von selten aktiven Neuronen entfernen
		\end{itemize}
		
		Für zeilen-priorisierte Inferenz ist $\alpha = 0$ optimal (nur In-Activity zählt).
	\end{satz}
	
	\begin{algorithm}[H]
		\caption{Asymmetrie-bewusstes Structured Pruning}
		\label{alg:asymmetry-aware-pruning}
		\textbf{Eingabe:} Trainiertes Netz $\{W_1, \ldots, W_L\}$, Trainingsdaten $\mathcal{D}$, Sparsity-Ziel $p_{\text{target}}$ \\
		\textbf{Ausgabe:} Gepruntes Netz $\{\tilde{W}_1, \ldots, \tilde{W}_L\}$
		\begin{algorithmic}[1]
			\State \textbf{Schritt 1: Aktivitätsanalyse}
			\For{$\text{batch} \in \mathcal{D}$}
			\State Forward-Pass durch Netz
			\For{jede Schicht $l$}
			\State Zähle Aktivierungen: $\text{count}_l[j] \gets \text{count}_l[j] + \mathbb{1}_{a_{l,j} \neq 0}$
			\EndFor
			\EndFor
			\State Normalisiere: $\text{out-activity}_l[j] \gets \text{count}_l[j] / |\mathcal{D}|$
			
			\State \textbf{Schritt 2: Importance-Berechnung}
			\For{jede Schicht $l$}
			\For{jedes Gewicht $W_{l,ij}$}
			\State $\text{Importance}_{l,ij} \gets |W_{l,ij}| \cdot \text{out-activity}_{l-1}[j]$
			\EndFor
			\EndFor
			
			\State \textbf{Schritt 3: Pruning}
			\State Sortiere alle Gewichte nach Importance (aufsteigend)
			\State $k \gets \lfloor p_{\text{target}} \cdot \sum_l |W_l| \rfloor$
			\State Setze die $k$ unwichtigsten Gewichte auf Null
			
			\State \textbf{Schritt 4: Fine-Tuning}
			\State Trainiere gepruntes Netz für wenige Epochen
			
			\State \Return $\{\tilde{W}_1, \ldots, \tilde{W}_L\}$
		\end{algorithmic}
	\end{algorithm}
	
	\begin{beispiel}[ResNet-50 Pruning]
		Anwendung von Algorithmus \ref{alg:asymmetry-aware-pruning} auf ResNet-50 (ImageNet):
		
		\textbf{Netzwerk-Charakteristika}:
		\begin{itemize}
			\item Parameter: 25.6 Millionen Gewichte
			\item Schichten: 50 Convolutional Layers
			\item Baseline-Genauigkeit: 76.2\% Top-1 auf ImageNet
		\end{itemize}
		
		\textbf{Pruning-Ergebnisse}:
		\begin{center}
			\begin{tabular}{|l|c|c|c|c|}
				\hline
				\textbf{Methode} & \textbf{Sparsity} & \textbf{Genauigkeit} & \textbf{Inferenzzeit} & \textbf{Speedup} \\
				\hline
				Baseline (dense) & 0\% & 76.2\% & 43 ms & 1.0$\times$ \\
				Magnitude-Pruning & 70\% & 75.8\% & 18 ms & 2.4$\times$ \\
				Asymmetry-aware & 70\% & 76.0\% & 12 ms & 3.6$\times$ \\
				Magnitude-Pruning & 90\% & 73.1\% & 8.7 ms & 4.9$\times$ \\
				Asymmetry-aware & 90\% & 74.3\% & 5.2 ms & 8.3$\times$ \\
				\hline
			\end{tabular}
		\end{center}
		
		\textbf{Beobachtungen}:
		\begin{enumerate}
			\item Asymmetrie-bewusstes Pruning erreicht bessere Genauigkeit bei gleicher Sparsity
			\item Speedup ist höher, da Pruning gezielt spalten-orientierte Struktur optimiert
			\item Bei 90\% Sparsity: 1.2 Prozentpunkte bessere Genauigkeit, 1.7$\times$ schneller
		\end{enumerate}
	\end{beispiel}
	
	\subsection{Training-Strategien für optimale Sparsity-Muster}
	
	\begin{definition}[Sparsity-induzierende Regularisierung]
		Während des Trainings kann Sparsity durch modifizierte Loss-Funktion gefördert werden:
		\[
		\mathcal{L}_{\text{total}} = \mathcal{L}_{\text{task}} + \lambda_1 \|W\|_1 + \lambda_2 \sum_{j} \left(\sum_i |W_{ij}|\right)^2
		\]
		
		Terme:
		\begin{itemize}
			\item $\mathcal{L}_{\text{task}}$: Aufgabenspezifischer Loss (z.B. Cross-Entropy)
			\item $\lambda_1 \|W\|_1$: L1-Regularisierung (fördert sparse Gewichte)
			\item $\lambda_2 \sum_{j} (\sum_i |W_{ij}|)^2$: Spalten-Gruppierungs-Term (fördert, dass ganze Spalten Null werden)
		\end{itemize}
	\end{definition}
	
	\begin{satz}[Konvergenz zu strukturierter Sparsity]
		Die modifizierte Loss-Funktion fördert Gewichtsmatrizen mit folgenden Eigenschaften:
		\begin{enumerate}
			\item Spalten-Sparsity: Ganze Spalten (korrespondierend zu Input-Neuronen) werden Null
			\item Dies ist optimal für spalten-priorisierte Inferenz
			\item Theoretischer Speedup: $1 / (1 - p_{\text{column}})$ wobei $p_{\text{column}}$ der Anteil Null-Spalten
		\end{enumerate}
	\end{satz}
	
	\begin{beispiel}[LSTM für Sprachmodellierung]
		Training eines LSTM (4 Schichten, 1024 Hidden Units) auf Penn Treebank:
		
		\textbf{Training-Konfiguration}:
		\begin{itemize}
			\item Baseline: Standardtraining ohne Sparsity-Regularisierung
			\item Sparse: $\lambda_1 = 10^{-5}$, $\lambda_2 = 10^{-4}$
			\item Structured-Sparse: $\lambda_1 = 10^{-5}$, $\lambda_2 = 10^{-3}$ (höherer Gruppierungs-Term)
		\end{itemize}
		
		\textbf{Ergebnisse nach Konvergenz}:
		\begin{center}
			\begin{tabular}{|l|c|c|c|c|}
				\hline
				\textbf{Methode} & \textbf{Sparsity} & \textbf{Perplexity} & \textbf{Inferenzzeit} & \textbf{Speedup} \\
				\hline
				Baseline & 0\% & 72.3 & 18.7 ms & 1.0$\times$ \\
				Sparse & 68\% & 73.8 & 9.2 ms & 2.0$\times$ \\
				Structured-Sparse & 71\% & 74.1 & 4.3 ms & 4.3$\times$ \\
				\hline
			\end{tabular}
		\end{center}
		
		Structured-Sparse erzielt mehr als doppelten Speedup bei nur minimal höherer Perplexity, da die Sparsity gezielt spalten-orientiert ist.
	\end{beispiel}
	
	\section{Finanzmathematik und Risikomanagement}
	
	Finanzmärkte sind hochgradig vernetzte Systeme mit charakteristischer Asymmetrie: Wenige systemisch wichtige Institutionen beeinflussen viele andere.
	
	\subsection{Portfolio-Optimierung mit sparse Korrelationsmatrizen}
	
	\begin{definition}[Markowitz Mean-Variance-Optimierung]
		Die klassische Portfolio-Opti\-mierung löst:
		\[
		\min_{w} w^T \Sigma w \quad \text{s.t.} \quad w^T \mu = \mu_{\text{target}}, \quad \sum_i w_i = 1
		\]
		
		wobei:
		\begin{itemize}
			\item $w \in \mathbb{R}^n$: Portfolio-Gewichte (Anteil jedes Assets)
			\item $\Sigma \in \mathbb{R}^{n \times n}$: Kovarianzmatrix der Returns
			\item $\mu \in \mathbb{R}^n$: Erwartete Returns
			\item $\mu_{\text{target}}$: Ziel-Return
		\end{itemize}
	\end{definition}
	
	\begin{satz}[Sparsity in Korrelationsmatrizen]
		Empirische Korrelationsmatrizen großer Asset-Universen zeigen charakteristische Struktur:
		\begin{itemize}
			\item Sektorale Cluster: Assets innerhalb eines Sektors stark korreliert
			\item Geringe Cross-Sektor-Korrelation: Typisch $|r_{ij}| < 0.1$ für $i, j$ in verschiedenen Sektoren
			\item Sparse Approximation: Setze $\Sigma_{ij} = 0$ für $|r_{ij}| < \tau$ (z.B. $\tau = 0.05$)
		\end{itemize}
		
		Für S\&P 500 mit 11 Sektoren: Sparsity $\approx 85\%$ nach Thresholding.
	\end{satz}
	
	\begin{beispiel}[S\&P 500 Portfolio-Optimierung]
		Optimierung eines Portfolios über alle S\&P 500 Aktien:
		
		\textbf{Daten}:
		\begin{itemize}
			\item Assets: $n = 500$
			\item Beobachtungsperiode: 5 Jahre (1260 Handelstage)
			\item Kovarianzmatrix: $500 \times 500 = 250.000$ Einträge
		\end{itemize}
		
		\textbf{Sparse Approximation}:
		\begin{itemize}
			\item Schwellwert: $\tau = 0.05$ (Korrelationen $< 0.05$ werden Null gesetzt)
			\item Resultierende Sparsity: 87.3\%
			\item Nicht-Null-Einträge: 31.750 (statt 250.000)
		\end{itemize}
		
		\textbf{Optimierungslaufzeit}:
		\begin{center}
			\begin{tabular}{|l|c|c|}
				\hline
				\textbf{Methode} & \textbf{Zeit} & \textbf{Speedup} \\
				\hline
				Dense Kovarianzmatrix & 2.37 s & 1.0$\times$ \\
				Sparse Kovarianzmatrix & 0.31 s & 7.6$\times$ \\
				\hline
			\end{tabular}
		\end{center}
		
		\textbf{Out-of-Sample Performance}:
		\begin{itemize}
			\item Dense: Sharpe Ratio = 1.23
			\item Sparse: Sharpe Ratio = 1.21
			\item Minimaler Performanceverlust ($< 2\%$) bei 7.6$\times$ Speedup
		\end{itemize}
	\end{beispiel}
	
	\subsection{Systemisches Risiko in Bankennetzwerken}
	
	\begin{definition}[Banken-Expositions-Netzwerk]
		Ein Bankensystem mit $n$ Banken wird als gerichtetes Netzwerk modelliert:
		\begin{itemize}
			\item Knoten: Banken
			\item Kante $(i,j)$ mit Gewicht $E_{ij}$: Bank $i$ hat Forderungen in Höhe $E_{ij}$ gegenüber Bank $j$
			\item Systemmatrix: $A_{ij} = E_{ij} / C_i$ (Exposition relativ zu Eigenkapital von Bank $i$)
		\end{itemize}
		
		Schock-Propagation:
		\[
		\ell_{t+1} = \ell_t + A \ell_t
		\]
		wobei $\ell_i(t)$ der Verlust von Bank $i$ zum Zeitpunkt $t$ ist.
	\end{definition}
	
	\begin{satz}[Hub-Struktur in Bankennetzwerken]
		Reale Bankennetzwerke zeigen extreme Broadcasting-Struktur:
		\begin{itemize}
			\item Wenige systemisch wichtige Banken (SIFIs): Out-Degree $\gg$ Durchschnitt
			\item Viele kleine Banken: Out-Degree $\approx$ konstant
			\item Power-Law-Verteilung mit Exponent $\gamma \approx 2.5$
		\end{itemize}
		
		Konsequenz: Ein Schock bei einer SIFI propagiert zu vielen Banken gleichzeitig (spalten-orientierte Ausbreitung optimal).
	\end{satz}
	
	\begin{beispiel}[Europäisches Bankennetzwerk]
		Simulation eines Schock-Szenarios im europäischen Bankensystem:
		
		\textbf{Netzwerk-Charakteristika}:
		\begin{itemize}
			\item Banken: $n = 50$ (30 G-SIBs, 20 kleinere Institute)
			\item Gesamtexpositionen: €8.7 Billionen
			\item Sparsity: 73\% (meiste Banken interagieren nur mit wenigen anderen)
		\end{itemize}
		
		\textbf{Schock-Szenario}:
		\begin{itemize}
			\item Initialschock: Eine G-SIB erleidet 30\% Kapitalverlust
			\item Propagation: Über Expositions-Netzwerk
			\item Zeitskala: 10 Perioden (z.B. Quartale)
		\end{itemize}
		
		\textbf{Simulationsergebnisse}:
		\begin{center}
			\begin{tabular}{|l|c|c|}
				\hline
				\textbf{Metrik} & \textbf{Dense} & \textbf{Sparse (spalten-orientiert)} \\
				\hline
				Rechenzeit (10 Perioden) & 0.73 s & 0.09 s \\
				Speedup & 1.0$\times$ & 8.1$\times$ \\
				Anzahl insolventer Banken & 7 & 7 \\
				Gesamtverlust & €340 Mrd. & €340 Mrd. \\
				\hline
			\end{tabular}
		\end{center}
		
		Die sparse Methode liefert identische Ergebnisse bei 8$\times$ Speedup. Für Monte-Carlo-Stress-Tests mit 10.000 Szenarien:
		\begin{itemize}
			\item Dense: $0.73 \times 10.000 = 7.300$ s $\approx$ 2 Stunden
			\item Sparse: $0.09 \times 10.000 = 900$ s $\approx$ 15 Minuten
		\end{itemize}
	\end{beispiel}
	
	\begin{korollar}[Praktische Implikationen für Regulierung]
		Die effiziente Berechnung systemischer Risikomaße ermöglicht:
		\begin{enumerate}
			\item Echtzeit-Überwachung systemischer Risiken
			\item Häufigere Stress-Tests (monatlich statt jährlich)
			\item Detailliertere Szenario-Analysen
			\item Quantifizierung von Tail-Risiken via Monte-Carlo
		\end{enumerate}
		
		Dies verbessert die Fähigkeit von Regulatoren, systemische Krisen frühzeitig zu erkennen.
	\end{korollar}
	
	\section{Zusammenfassung der Anwendungen}
	
	Die vier Fallstudien demonstrieren die Vielseitigkeit der entwickelten Methoden:
	
	\begin{enumerate}
		\item \textbf{Robotersteuerung}: Hierarchische Struktur → Faktor 20 Speedup → Echtzeit-MPC möglich
		\item \textbf{Netzwerkdynamik}: Broadcasting-Struktur → Faktor 40+ Speedup → Monte-Carlo-Simulation praktikabel
		\item \textbf{Neuronale Netze}: Kombinierte Sparsity → Faktor 8+ Speedup → Edge-Device-Inferenz
		\item \textbf{Finanzrisiko}: Sparse Korrelationen → Faktor 8 Speedup → Häufigere Stress-Tests
	\end{enumerate}
	
	Gemeinsame Erfolgsfaktoren:
	\begin{itemize}
		\item Hohe Sparsity ($> 70\%$) in System- oder Inputstrukturen
		\item Klare Asymmetrie zwischen In- und Out-Degree
		\item Wiederholte Berechnungen (amortisieren Strukturerkennung)
		\item Große Systemgrößen ($n > 100$)
	\end{itemize}
	
	\chapter{Theoretische Vertiefung: Spektrale Eigenschaften und Langzeitverhalten}
	
	Die bisherige Analyse fokussierte auf algorithmische Effizienz. Dieses Kapitel vertieft das theoretische Verständnis durch spektrale Analyse und untersucht das Langzeitverhalten dynamischer Systeme.
	
	\section{Spektrale Analyse asymmetrischer Systemmatrizen}
	
	Die Eigenwerte einer Matrix kodieren fundamentale Eigenschaften des Systems. Für asymmetrische Matrizen besteht ein direkter Zusammenhang zwischen Degree-Verteilung und Eigenwertspektrum.
	
	\subsection{Gershgorin-Kreise und Degree-Schranken}
	
	\begin{satz}[Gershgorin-Kreise für In- und Out-Degree]
		Für Systemmatrix $A \in \mathbb{R}^{n \times n}$ liegen alle Eigenwerte in der Vereinigung der Gershgorin-Kreise:
		\[
		\lambda \in \bigcup_{i=1}^n \left\{ z \in \mathbb{C} : |z - A_{ii}| \leq R_i \right\}
		\]
		wobei $R_i = \sum_{j \neq i} |A_{ij}|$ der Zeilenradius ist.
		
		\textbf{Zusammenhang mit In-Degree}:
		Für sparse Matrix mit In-Degree $d_{\text{in}}(i)$ gilt:
		\[
		R_i \approx \bar{w} \cdot d_{\text{in}}(i)
		\]
		wobei $\bar{w}$ das durchschnittliche Gewicht ist.
		
		\textbf{Implikation}: Zeilen mit hohem In-Degree führen zu größeren Kreisen → potenziell größere Eigenwerte.
	\end{satz}
	
	\begin{proof}
		Nach dem Gershgorin-Kreis-Theorem liegt jeder Eigenwert $\lambda$ in mindestens einem Kreis:
		\[
		|Ax - \lambda x| = 0 \implies |A_{ii} - \lambda| \leq \sum_{j \neq i} |A_{ij}| \frac{|x_j|}{|x_i|}
		\]
		
		Für die entsprechende Komponente $i$ mit $|x_i| = \max_k |x_k|$ folgt:
		\[
		|\lambda - A_{ii}| \leq \sum_{j \neq i} |A_{ij}| = R_i
		\]
		
		Für sparse Matrix ist $R_i$ proportional zur Anzahl nicht-null Einträge in Zeile $i$, also dem In-Degree.
	\end{proof}
	
	\begin{korollar}[Spektralradius-Schranke]
		Der Spektralradius $\rho(A) = \max_i |\lambda_i|$ ist beschränkt durch:
		\[
		\rho(A) \leq \max_i \left( |A_{ii}| + \sum_{j \neq i} |A_{ij}| \right) \leq \|A\|_{\infty} = \max_i \sum_j |A_{ij}|
		\]
		
		Dies ist die maximale Zeilensumme, direkt abhängig vom maximalen In-Degree.
	\end{korollar}
	
	\begin{beispiel}[Broadcasting-Matrix]
		Für Broadcasting-Matrix mit Hub-Struktur:
		\begin{itemize}
			\item Normale Zeilen: $R_i \approx 5$ (wenige Einflüsse)
			\item Hub-Zeilen: $R_i \approx 5000$ (viele Einflüsse von Hub)
			\item Gershgorin-Kreise: Extreme Varianz in Größe
			\item Spektralradius: Dominiert durch Hub-Zeilen
		\end{itemize}
		
		Dies erklärt, warum Broadcasting-Systeme oft einen dominanten Eigenwert haben, der weit vom Rest des Spektrums entfernt ist.
	\end{beispiel}
	
	\subsection{Perron-Frobenius-Theorie}
	
	\begin{satz}[Perron-Frobenius für nicht-negative Matrizen]
		Für nicht-negative, irreduzible Matrix $A \geq 0$ existiert ein dominanter Eigenwert $\lambda_1 > 0$ mit:
		\begin{enumerate}
			\item $\lambda_1$ ist einfach (algebraische Multiplizität 1)
			\item $\lambda_1 > |\lambda_i|$ für alle $i \neq 1$
			\item Der zugehörige Eigenvektor $v_1$ hat nur positive Komponenten
			\item $\lambda_1 \leq \max_i \sum_j A_{ij}$ (maximale Zeilensumme)
			\item $\lambda_1 \geq \min_i \sum_j A_{ij}$ (minimale Zeilensumme)
		\end{enumerate}
	\end{satz}
	
	\begin{korollar}[Asymmetrie und dominanter Eigenwert]
		Für Broadcasting-Matrix mit extremer In-Degree-Asymmetrie gilt:
		\begin{itemize}
			\item Obere Schranke: $\lambda_1 \leq \max_i \text{in-deg}(i) \cdot \bar{w}$ (Hub-Zeilen)
			\item Untere Schranke: $\lambda_1 \geq \min_i \text{in-deg}(i) \cdot \bar{w}$ (normale Zeilen)
			\item Verhältnis: $\frac{\lambda_1^{\text{upper}}}{\lambda_1^{\text{lower}}} \approx \frac{\max \text{in-deg}}{\min \text{in-deg}}$ kann sehr groß sein
		\end{itemize}
		
		Dies impliziert: Broadcasting-Systeme haben typischerweise große spektrale Lücke ($\lambda_1 \gg \lambda_2$).
	\end{korollar}
	
	\begin{beispiel}[PageRank-Matrix]
		Die Google PageRank-Matrix ist ein klassisches Beispiel:
		\[
		M = \alpha S + (1-\alpha) \frac{1}{n} \mathbf{1}\mathbf{1}^T
		\]
		wobei $S$ die normalisierte Link-Matrix und $\alpha \approx 0.85$.
		
		\textbf{Spektrale Eigenschaften}:
		\begin{itemize}
			\item Dominanter Eigenwert: $\lambda_1 = 1$ (stationäre Verteilung existiert)
			\item Spektrale Lücke: $\lambda_2 \approx 0.85$ (schnelle Konvergenz zur stationären Verteilung)
			\item Konvergenzrate: $O(\lambda_2^t)$ → nach $t \approx 50$ Iterationen konvergiert
		\end{itemize}
		
		Die spalten-orientierte Power-Iteration nutzt die Sparsity von $S$ aus und beschleunigt die Berechnung um Faktor $\approx 100$ für Web-Graphen mit Milliarden Seiten.
	\end{beispiel}
	
	\section{Langzeitverhalten und Konvergenz}
	
	Das asymptotische Verhalten dynamischer Systeme wird durch das Eigenwertspektrum bestimmt.
	
	\subsection{Konvergenz zu stationären Zuständen}
	
	\begin{definition}[Stationäre Verteilung]
		Für Markov-Kette mit Übergangsmatrix $P$ ist $\pi \in \mathbb{R}^n$ eine stationäre Verteilung, falls:
		\[
		\pi = P^T \pi, \quad \sum_i \pi_i = 1, \quad \pi_i \geq 0 \quad \forall i
		\]
		
		Physikalisch: $\pi$ ist der Gleichgewichtszustand des Systems.
	\end{definition}
	
	\begin{satz}[Konvergenz für ergodische Ketten]
		Für irreduzible, aperiodische Markov-Kette mit Übergangsmatrix $P$ gilt:
		\[
		\lim_{t \to \infty} P^t = \mathbf{1} \pi^T
		\]
		wobei $\pi$ die eindeutige stationäre Verteilung ist.
		
		Die Konvergenzrate ist exponentiell:
		\[
		\|P^t - \mathbf{1}\pi^T\| \leq C \lambda_2^t
		\]
		wobei $\lambda_2$ der zweitgrößte Eigenwert (in Betrag) ist.
	\end{satz}
	
	\begin{beispiel}[Random Walk auf sozialem Netzwerk]
		Betrachten Sie einen Random Walk auf einem sozialen Netzwerk:
		\begin{itemize}
			\item Knoten: $n = 10.000$ Nutzer
			\item Kanten: Freundschaftsverbindungen
			\item Übergangsmatrix: $P_{ij} = 1/d_j$ falls $(i,j)$ Kante, sonst 0
		\end{itemize}
		
		\textbf{Spektrale Analyse}:
		\begin{itemize}
			\item $\lambda_1 = 1$ (stationäre Verteilung existiert)
			\item $\lambda_2 \approx 0.92$ (moderate spektrale Lücke)
			\item Mischungszeit: $t_{\text{mix}} \approx \frac{\log(n)}{1 - \lambda_2} \approx \frac{9.2}{0.08} \approx 115$ Schritte
		\end{itemize}
		
		\textbf{Berechnung der stationären Verteilung}:
		\begin{itemize}
			\item Power-Iteration: $\pi^{(t+1)} = P^T \pi^{(t)}$
			\item Spalten-orientiert: Nutzt Sparsity von $P$ ($\approx 95\%$)
			\item Konvergenz nach $\approx 150$ Iterationen
			\item Zeit pro Iteration: 0.3 ms (spalten-orientiert) vs. 12 ms (dense)
			\item Gesamtzeit: 45 ms vs. 1.8 s → Speedup 40$\times$
		\end{itemize}
	\end{beispiel}
	
	\subsection{Mischungszeiten und Asymmetrie}
	
	\begin{definition}[Mischungszeit]
		Die $\epsilon$-Mischungszeit einer Markov-Kette ist:
		\[
		t_{\text{mix}}(\epsilon) = \min \left\{ t : \max_{\text{Startverteilung } x_0} \|P^t x_0 - \pi\|_{\text{TV}} \leq \epsilon \right\}
		\]
		wobei $\|\cdot\|_{\text{TV}}$ die Total-Variation-Distanz ist.
	\end{definition}
	
	\begin{satz}[Mischungszeit und spektrale Lücke]
		Für reversible Markov-Kette gilt:
		\[
		t_{\text{mix}}(\epsilon) \leq \frac{1}{1-\lambda_2} \log\left(\frac{1}{\epsilon \pi_{\min}}\right)
		\]
		wobei $\pi_{\min} = \min_i \pi_i$ die kleinste stationäre Wahrscheinlichkeit.
		
		\textbf{Zusammenhang mit Asymmetrie}:
		Für Broadcasting-Netzwerke gilt typischerweise:
		\begin{itemize}
			\item Hubs haben hohe stationäre Wahrscheinlichkeit: $\pi_{\text{hub}} \gg \bar{\pi}$
			\item Normale Knoten: $\pi_{\text{normal}} \approx \bar{\pi}$
			\item Minimale Wahrscheinlichkeit: $\pi_{\min} \ll \bar{\pi}$
			\item Konsequenz: Längere Mischungszeiten
		\end{itemize}
	\end{satz}
	
	\begin{beispiel}[Einfluss der Hub-Struktur]
		Vergleich zweier Netzwerke mit $n = 1000$ Knoten:
		
		\textbf{Netzwerk A (gleichverteilt)}:
		\begin{itemize}
			\item Alle Knoten: Degree $\approx 10$
			\item $\lambda_2 \approx 0.95$
			\item $\pi_{\min} \approx 1/1000$
			\item $t_{\text{mix}}(0.01) \approx 140$ Schritte
		\end{itemize}
		
		\textbf{Netzwerk B (Broadcasting)}:
		\begin{itemize}
			\item 10 Hubs: Degree $\approx 500$
			\item 990 normale: Degree $\approx 5$
			\item $\lambda_2 \approx 0.98$ (größer!)
			\item $\pi_{\min} \approx 1/100.000$ (viel kleiner!)
			\item $t_{\text{mix}}(0.01) \approx 580$ Schritte (4$\times$ länger!)
		\end{itemize}
		
		Broadcasting-Struktur verlangsamt paradoxerweise die Konvergenz zur stationären Verteilung, obwohl Information schnell propagiert.
	\end{beispiel}
	
	\section{Sensitivitätsanalyse}
	
	Die Robustheit der entwickelten Algorithmen gegenüber Störungen ist entscheidend für praktische Anwendungen.
	
	\subsection{Perturbationstheorie für Eigenwerte}
	
	\begin{satz}[Bauer-Fike-Theorem]
		Sei $A = X \Lambda X^{-1}$ diagonalisierbar mit Eigenwerten $\lambda_1, \ldots, \lambda_n$. Für gestörte Matrix $\tilde{A} = A + E$ liegt jeder Eigenwert $\tilde{\lambda}$ von $\tilde{A}$ in mindestens einem der Kreise:
		\[
		|\tilde{\lambda} - \lambda_i| \leq \kappa(X) \|E\|
		\]
		wobei $\kappa(X) = \|X\| \|X^{-1}\|$ die Konditionszahl der Eigenvektormatrix.
		
		\textbf{Implikation}: Schlecht konditionierte Eigenvektoren führen zu hoher Sensitivität.
	\end{satz}
	
	\begin{korollar}[Konditionierung asymmetrischer Matrizen]
		Für asymmetrische Matrizen (nicht-normal) kann $\kappa(X) \gg 1$ sein:
		\begin{itemize}
			\item Normale Matrizen ($AA^T = A^T A$): $\kappa(X) = 1$ (optimal)
			\item Symmetrische Matrizen: $\kappa(X) = 1$ (Eigenvektoren orthogonal)
			\item Broadcasting-Matrizen: $\kappa(X)$ kann $10^6$ oder größer sein
		\end{itemize}
		
		Dies bedeutet: Kleine Störungen in $A$ können große Änderungen in Eigenwerten verursachen.
	\end{korollar}
	
	\begin{beispiel}[Sensitivität von Broadcasting-Systemen]
		Für Broadcasting-Matrix mit extremer Asymmetrie:
		\begin{itemize}
			\item Exakte Matrix $A$: $\lambda_1 = 5.0$, $\lambda_2 = 0.3$
			\item Konditionszahl: $\kappa(X) = 10^5$
			\item Störung: $\|E\| = 10^{-6}$ (numerische Rundungsfehler)
			\item Eigenwert-Änderung: $|\tilde{\lambda}_1 - \lambda_1| \leq 10^5 \cdot 10^{-6} = 0.1$
		\end{itemize}
		
		Der dominante Eigenwert kann sich um bis zu 2\% ändern durch winzige Störungen!
		
		\textbf{Praktische Konsequenz}: Für schlecht konditionierte Systeme sollte die Eigenwertmethode vorsichtig verwendet werden. Iterative Methoden (Power-Iteration) sind oft robuster.
	\end{beispiel}
	
	\subsection{Robustheit der Sparsity-Ausnutzung}
	
	\begin{satz}[Stabilität spalten-orientierter Algorithmen]
		Für Matrix-Vektor-Multiplikation $y = Ax$ mit gestörten Daten:
		\begin{itemize}
			\item Gestörte Matrix: $\tilde{A} = A + \Delta A$ mit $\|\Delta A\| \leq \epsilon_A \|A\|$
			\item Gestörter Vektor: $\tilde{x} = x + \Delta x$ mit $\|\Delta x\| \leq \epsilon_x \|x\|$
		\end{itemize}
		
		Der relative Fehler ist beschränkt durch:
		\[
		\frac{\|\tilde{y} - y\|}{\|y\|} \leq \left( \kappa(A) + 1 \right) \max(\epsilon_A, \epsilon_x) + O(\epsilon^2)
		\]
		
		Dies ist \textbf{unabhängig} von der Berechnungsmethode (zeilen- vs. spalten-orientiert).
		
		\textbf{Konsequenz}: Sparsity-Ausnutzung ändert Effizienz, nicht Stabilität.
	\end{satz}
	
	\begin{proof}
		Standardresultat aus numerischer linearer Algebra. Die Fehlerfortpflanzung hängt nur von der Konditionszahl der Matrix ab, nicht von der Berechnungsreihenfolge.
	\end{proof}
	
	\begin{beispiel}[Numerische Experimente zur Robustheit]
		Test mit sparse Matrix $A \in \mathbb{R}^{1000 \times 1000}$, Sparsity 90\%:
		
		\textbf{Störungen}:
		\begin{itemize}
			\item Rundungsfehler: double precision ($\epsilon \approx 2.2 \times 10^{-16}$)
			\item Zusätzliches Rauschen: $\Delta A_{ij} \sim \mathcal{N}(0, 10^{-10})$
		\end{itemize}
		
		\textbf{Gemessene Fehler} (1000 Durchläufe):
		\begin{center}
			\begin{tabular}{|l|c|c|}
				\hline
				\textbf{Methode} & \textbf{Relativer Fehler (Mittel)} & \textbf{Standardabweichung} \\
				\hline
				Dense & $1.3 \times 10^{-15}$ & $3.2 \times 10^{-16}$ \\
				Sparse Row-oriented & $1.4 \times 10^{-15}$ & $3.5 \times 10^{-16}$ \\
				Sparse Column-oriented & $1.3 \times 10^{-15}$ & $3.3 \times 10^{-16}$ \\
				\hline
			\end{tabular}
		\end{center}
		
		Alle Methoden haben praktisch identische Fehler → Robustheit bestätigt.
	\end{beispiel}
	
	\section{Verallgemeinerung auf Tensorsysteme}
	
	Viele moderne Anwendungen (Deep Learning, Quantenchemie, Data Mining) verwenden höherdimensionale Arrays (Tensoren). Die Asymmetrie-Konzepte lassen sich erweitern.
	
	\subsection{Tensor-Vektor-Multiplikation}
	
	\begin{definition}[Mode-$k$-Produkt]
		Für Tensor $\mathcal{A} \in \mathbb{R}^{n_1 \times n_2 \times \cdots \times n_d}$ und Vektor $v \in \mathbb{R}^{n_k}$ ist das Mode-$k$-Produkt:
		\[
		(\mathcal{A} \times_k v)_{i_1, \ldots, i_{k-1}, i_{k+1}, \ldots, i_d} = \sum_{i_k=1}^{n_k} \mathcal{A}_{i_1, \ldots, i_d} \cdot v_{i_k}
		\]
		
		Dies verallgemeinert Matrix-Vektor-Multiplikation auf Tensoren.
	\end{definition}
	
	\begin{satz}[Asymmetrie in Tensor-Modi]
		Analog zum Matrix-Fall gibt es für Tensoren $d$ verschiedene "Richtungen" der Multiplikation (Modi). Für sparse Tensor mit Sparsity-Muster:
		\begin{itemize}
			\item Mode-$k$-Sparsity: Anteil der Null-Fasern entlang Mode $k$
			\item Asymmetrie: Verschiedene Modi können drastisch unterschiedliche Sparsity haben
		\end{itemize}
		
		Optimale Berechnungsreihenfolge: Beginne mit dem sparsesten Mode.
	\end{satz}
	
	\begin{beispiel}[3D Convolutional Neural Networks]
		Ein 3D-CNN verarbeitet Video-Daten (Zeit $\times$ Höhe $\times$ Breite):
		
		\textbf{Activation Tensor}: $\mathcal{A} \in \mathbb{R}^{T \times H \times W}$ mit $T=16$, $H=W=224$
		
		\textbf{Sparsity-Analyse} (nach ReLU):
		\begin{itemize}
			\item Zeit-Mode: 45\% Sparsity (viele Frames haben ähnliche Aktivierungen)
			\item Raum-Mode: 78\% Sparsity (viele räumliche Regionen inaktiv)
			\item Unterschied: Faktor 1.7 in Sparsity
		\end{itemize}
		
		\textbf{Optimale Berechnungsreihenfolge}:
		Verarbeite zuerst räumliche Modi (höhere Sparsity), dann zeitlichen Mode.
		
		\textbf{Speedup}: 3.2$\times$ gegenüber naiver Reihenfolge
	\end{beispiel}
	
	\subsection{Tucker-Zerlegung und Asymmetrie}
	
	\begin{definition}[Tucker-Zerlegung]
		Die Tucker-Zerlegung approximiert einen Tensor als:
		\[
		\mathcal{A} \approx \mathcal{G} \times_1 U^{(1)} \times_2 U^{(2)} \times_3 \cdots \times_d U^{(d)}
		\]
		wobei:
		\begin{itemize}
			\item $\mathcal{G} \in \mathbb{R}^{r_1 \times r_2 \times \cdots \times r_d}$: Kern-Tensor (klein)
			\item $U^{(k)} \in \mathbb{R}^{n_k \times r_k}$: Faktor-Matrizen (orthogonal)
			\item $r_k \ll n_k$: Rang entlang Mode $k$
		\end{itemize}
	\end{definition}
	
	\begin{satz}[Asymmetrische Ränge]
		Für Tensoren aus realen Anwendungen sind die Ränge oft stark asymmetrisch:
		\begin{itemize}
			\item Video-Daten: Zeitlicher Rang $r_t \ll$ räumlicher Rang $r_s$ (temporale Kohärenz)
			\item fMRI-Gehirn-Daten: Räumlicher Rang $r_{\text{space}} \gg$ temporaler Rang $r_{\text{time}}$ (viele Hirnregionen, wenige Zeitskalen)
		\end{itemize}
		
		Optimale Kompressionsstrategie: Komprimiere zuerst Modi mit niedrigem Rang.
	\end{satz}
	
	\begin{beispiel}[Video-Kompression]
		Video-Tensor: $\mathcal{V} \in \mathbb{R}^{1920 \times 1080 \times 1000}$ (Full HD, 1000 Frames)
		
		\textbf{Tucker-Zerlegung}:
		\begin{itemize}
			\item Zeitlicher Rang: $r_t = 5$ (nur wenige Bewegungsmuster)
			\item Räumlicher Rang: $r_h = r_w = 50$ (komplexere räumliche Strukturen)
		\end{itemize}
		
		\textbf{Kompressionsverhältnis}:
		\begin{itemize}
			\item Original: $1920 \times 1080 \times 1000 = 2.07 \times 10^9$ Werte
			\item Tucker: $5 \times 50 \times 50 + 1920 \times 50 + 1080 \times 50 + 1000 \times 5 = 117.500$ Werte
			\item Verhältnis: $\approx 17.600:1$
		\end{itemize}
		
		Die extreme Asymmetrie ($r_t \ll r_h, r_w$) ermöglicht diese massive Kompression.
	\end{beispiel}
	
	\section{Zusammenfassung der theoretischen Vertiefung}
	
	Die spektrale Analyse und Langzeituntersuchung liefern tiefere Einsichten:
	
	\begin{enumerate}
		\item \textbf{Gershgorin-Kreise}: In-Degree-Asymmetrie manifestiert sich in Eigenwertlokalisierung
		\item \textbf{Perron-Frobenius}: Broadcasting-Systeme haben große spektrale Lücke
		\item \textbf{Mischungszeiten}: Hub-Strukturen verlangsamen paradoxerweise Konvergenz
		\item \textbf{Sensitivität}: Asymmetrische Matrizen können schlecht konditioniert sein
		\item \textbf{Robustheit}: Sparsity-Ausnutzung ändert Effizienz, nicht numerische Stabilität
		\item \textbf{Tensor-Verallgemeinerung}: Asymmetrie-Konzepte übertragen sich auf höhere Dimensionen
	\end{enumerate}
	
	Diese theoretischen Erkenntnisse komplementieren die algorithmischen Resultate und liefern ein umfassendes Verständnis asymmetrischer Systemmatrizen.
	
	\chapter{Überarbeitete Zusammenfassung und Ausblick}
	
	\section{Synthese der Kernerkenntnisse}
	
	Diese Arbeit hat die Asymmetrie zwischen Zeilen und Spalten in Systemmatrizen auf drei Ebenen untersucht: theoretisch, algorithmisch und experimentell. Die zentrale These – dass diese Asymmetrie ein fundamentales Organisationsprinzip ist, das systematisch ausgenutzt werden kann – wurde umfassend validiert.
	
	\subsection{Theoretische Fundierung}
	
	Die semantische Unterscheidung zwischen Aggregation (Zeilen) und Ausbreitung (Spalten) ist mehr als eine konzeptionelle Betrachtung. Sie manifestiert sich konkret in:
	
	\begin{enumerate}
		\item \textbf{Algorithmischer Komplexität}: Spalten-priorisierte Berechnung nutzt sparse Eingabevektoren aus und erreicht Laufzeit $O(s \cdot \bar{d}_{\text{out}})$ statt $O(n \cdot \bar{d}_{\text{in}})$, wobei typischerweise $s \ll n$.
		
		\item \textbf{Broadcasting-Systemen}: Hub-Strukturen zeigen extreme Out-Degree-Varianz ($\text{Var}[\text{out-deg}] \gg \bar{d}^2$), die spalten-orientierte Methoden optimal ausnutzen.
		
		\item \textbf{Hierarchischen Strukturen}: Nilpotenz mit Index $h$ (Ebenenanzahl) führt zu deterministischer transienter Dynamik und ermöglicht Beschleunigungen um Faktor $h$.
		
		\item \textbf{Spektralen Eigenschaften}: Gershgorin-Kreise verbinden In-Degree mit Eigenwertlokalisierung; Perron-Frobenius-Theorie erklärt dominante Eigenwerte in Broadcasting-Systemen.
	\end{enumerate}
	
	\subsection{Experimentelle Validierung und praktische Grenzen}
	
	Die sieben durchgeführten Experimente lieferten sowohl Bestätigungen als auch wichtige Erkenntnisse über praktische Grenzen:
	
	\textbf{Bestätigte theoretische Vorhersagen}:
	\begin{itemize}
		\item Sparse Matrix-Vektor-Multiplikation: 4-23$\times$ Speedup bei 90-99\% Input-Sparsity (Experiment 2)
		\item Neuronale Netze: Bis zu 39.5$\times$ Speedup bei kombinierter Gewichts- und Aktivierungs-Sparsity (Experiment 3)
		\item Hierarchische Systeme: Nilpotenz-Eigenschaften exakt wie vorhergesagt (Experiment 6)
		\item Azyklische Systeme: Transiente Dynamik mit Konvergenz zum Nullzustand (Experiment 7)
	\end{itemize}
	
	\textbf{Entdeckte praktische Grenzen}:
	\begin{itemize}
		\item Block-Diagonal-Optimierung: Overhead dominiert bei $n < 100$ → Speedup nur 0.21$\times$ statt theoretisch $m^2$ (Experiment 5)
		\item Hochoptimierte BLAS: Für dichte Matrizen mittlerer Größe schwer zu schlagen durch strukturelle Optimierung
		\item Amortisierung: Strukturerkennung ($O(n^3)$) lohnt sich nur bei $> 100$ wiederholten Berechnungen
	\end{itemize}
	
	\subsection{Anwendungsbreite}
	
	Die vier detaillierten Fallstudien (Kapitel 10) demonstrierten die Vielseitigkeit:
	
	\begin{center}
		\begin{tabular}{|l|l|c|l|}
			\hline
			\textbf{Anwendung} & \textbf{Struktur} & \textbf{Speedup} & \textbf{Impakt} \\
			\hline
			Roboter-MPC & Hierarchisch & 20$\times$ & Echtzeit-MPC @ 1 kHz \\
			Netzwerkdynamik & Broadcasting & 40+$\times$ & Monte-Carlo praktikabel \\
			Neuronale Netze & Sparse & 8+$\times$ & Edge-Device-Inferenz \\
			Finanzrisiko & Sparse & 8$\times$ & Häufigere Stress-Tests \\
			\hline
		\end{tabular}
	\end{center}
	
	Gemeinsamer Erfolgsfaktor: Hohe Sparsity ($> 70\%$) kombiniert mit klarer struktureller Asymmetrie.
	
	\section{Praktischer Implementierungsleitfaden}
	
	Basierend auf allen theoretischen und experimentellen Erkenntnissen ergibt sich folgender Leitfaden:
	
	\subsection{Entscheidungsbaum}
	
	\begin{algorithm}[H]
		\caption{Strategie-Auswahl für Matrix-basierte Systeme}
		\label{alg:final-decision-tree}
		\textbf{Eingabe:} Systemmatrix $A$, geplante Anzahl Berechnungen $N_{\text{ops}}$ \\
		\textbf{Ausgabe:} Empfohlene Implementierungsstrategie
		\begin{algorithmic}[1]
			\State Messe Systemparameter: $n$ (Größe), $m$ (Nicht-Null-Einträge)
			\State $\text{sparsity} \gets 1 - m/n^2$
			
			\If{$n < 100$}
			\State \Return \texttt{USE\_BLAS} \Comment{Overhead dominiert}
			\EndIf
			
			\If{$\text{sparsity} < 0.5$}
			\If{$N_{\text{ops}} > 100$ \textbf{and} $n > 500$}
			\State Berechne Eigenwertzerlegung (einmalig)
			\State \Return \texttt{USE\_EIGENVALUE\_METHOD}
			\Else
			\State \Return \texttt{USE\_DENSE\_BLAS}
			\EndIf
			\EndIf
			
			\State Analysiere Struktur: $\text{type} \gets $ Strukturerkennung($A$)
			
			\If{$\text{type} = $ \texttt{BLOCK\_DIAGONAL} \textbf{and} $\text{block\_size} > 500$}
			\State \Return \texttt{USE\_BLOCK\_DECOMPOSITION}
			\ElsIf{$\text{type} = $ \texttt{HIERARCHICAL}}
			\State \Return \texttt{USE\_LAYER\_ADAPTIVE}
			\ElsIf{$\text{sparsity} > 0.7$}
			\If{erwarte sparse Eingabevektoren}
			\State \Return \texttt{USE\_COLUMN\_ORIENTED\_SPARSE}
			\Else
			\State \Return \texttt{USE\_ROW\_ORIENTED\_SPARSE}
			\EndIf
			\Else
			\State \Return \texttt{USE\_ADAPTIVE} \Comment{Wechselt dynamisch}
			\EndIf
		\end{algorithmic}
	\end{algorithm}
	
	\subsection{Implementierungs-Checkliste}
	
	Für erfolgreiche Implementierung:
	
	\begin{enumerate}
		\item \textbf{Profilierung}: Messe reale Sparsity-Muster auf Produktionsdaten, nicht synthetischen Beispielen
		\item \textbf{Schwellwerte}: Kalibriere numerische Schwellwerte ($\tau \approx 10^{-12}$ für double precision)
		\item \textbf{Speicherformat}: Wähle CSC für spalten-orientiert, CSR für zeilen-orientiert
		\item \textbf{Cache-Optimierung}: Berücksichtige Cache-Linien (typisch 64 Bytes)
		\item \textbf{Fallback}: Implementiere dense Fallback für Fälle, wo Sparsity-Annahmen verletzt werden
		\item \textbf{Validierung}: Vergleiche Ergebnisse mit Referenzimplementierung (Genauigkeit)
		\item \textbf{Benchmarking}: Messe auf Zielarchitektur, nicht Entwicklungsmaschine
	\end{enumerate}
	
	\section{Forschungsagenda}
	
	Die Arbeit öffnet mehrere vielversprechende Forschungsrichtungen:
	
	\subsection{Kurzfristige Forschungsziele (1-2 Jahre)}
	
	\begin{enumerate}
		\item \textbf{Hardware-spezifische Optimierung}
		\begin{itemize}
			\item GPU-Implementierung spalten-orientierter Algorithmen mit CUDA
			\item TPU-Mapping für strukturierte Sparsity
			\item ARM NEON SIMD-Optimierung für Edge-Devices
			\item Benchmarking über diverse Architekturen
		\end{itemize}
		
		\item \textbf{Automatische Codegenerierung}
		\begin{itemize}
			\item Compiler-Extension für strukturbewusste Optimierung
			\item JIT-Kompilierung basierend auf Runtime-Sparsity
			\item Integration in MLIR (Multi-Level Intermediate Representation)
		\end{itemize}
		
		\item \textbf{Erweiterte Strukturerkennung}
		\begin{itemize}
			\item Machine Learning für Struktur-Klassifikation
			\item Automatische Parameterwahl ($\tau$, Algorithmus) via Bayessche Optimierung
			\item Online-Lernen: Adaptation an zeitvariante Sparsity
		\end{itemize}
	\end{enumerate}
	
	\subsection{Mittelfristige Forschungsziele (2-5 Jahre)}
	
	\begin{enumerate}
		\item \textbf{Theoretische Vertiefung}
		\begin{itemize}
			\item Beweis optimaler Speedup-Faktoren für spezifische Strukturklassen
			\item Untere Schranken: Kann man besser als $O(m)$ für sparse Matrizen?
			\item Approximationstheorie: Trade-off zwischen Genauigkeit und Geschwindigkeit
		\end{itemize}
		
		\item \textbf{Nichtlineare Verallgemeinerung}
		\begin{itemize}
			\item Asymmetrie in nichtlinearen dynamischen Systemen
			\item Neuronale ODEs mit adaptiver Sparsity
			\item Koopman-Operator-Theorie für nichtlineare Systeme
		\end{itemize}
		
		\item \textbf{Stochastische Systeme}
		\begin{itemize}
			\item Kalman-Filterung mit sparse Systemmatrizen
			\item Particle Filter: Spalten-orientierte Resampling-Strategien
			\item Uncertainty Quantification via sparse Polynomial Chaos
		\end{itemize}
		
		\item \textbf{Quantencomputing}
		\begin{itemize}
			\item Quantenalgorithmen für sparse Matrix-Operationen
			\item HHL-Algorithmus (Harrow-Hassidim-Lloyd) mit Strukturausnutzung
			\item Fehlerkorrektur in Quantensystemen
		\end{itemize}
	\end{enumerate}
	
	\subsection{Langfristige Vision (5+ Jahre)}
	
	\begin{enumerate}
		\item \textbf{Hardware-Software-Co-Design}
		\begin{itemize}
			\item Spezialisierte Prozessoren für spalten-orientierte Operationen
			\item Neuromorphe Hardware mit inhärenter Sparsity
			\item Optische Computer: Nutzung von Asymmetrie in photonischen Systemen
		\end{itemize}
		
		\item \textbf{Foundation Models für Systemdynamik}
		\begin{itemize}
			\item Pre-trained Models für Struktur-Prädiktion
			\item Transfer Learning zwischen verschiedenen Systemklassen
			\item Few-Shot Learning: Strukturerkennung mit minimalen Daten
		\end{itemize}
		
		\item \textbf{Autonome Systeme}
		\begin{itemize}
			\item Selbst-optimierende Algorithmen mit Online-Strukturlernen
			\item Adaptive Computing: Hardware rekonfiguriert sich basierend auf Struktur
			\item Bio-inspirierte Architekturen: Gehirn als Vorbild für asymmetrische Verarbeitung
		\end{itemize}
	\end{enumerate}
	
	\section{Offene theoretische Probleme}
	
	Mehrere fundamentale Fragen bleiben unbeantwortet:
	
	\begin{frage}[Optimale Schwellwerttheorie]
		Existiert eine optimale Strategie zur Wahl des Sparsity-Schwellwerts $\tau$, die den Trade-off zwischen Genauigkeit und Geschwindigkeit minimiert?
		
		Formal: Gegeben gewünschte Genauigkeit $\epsilon$ und Matrix $A$, finde:
		\[
		\tau^* = \arg\min_{\tau} T(\tau) \quad \text{s.t.} \quad \text{Error}(\tau) \leq \epsilon
		\]
		
		Dies ist ein Optimierungsproblem mit unbekannten Zielfunktionen (Laufzeit und Fehler hängen komplex von $\tau$ ab).
	\end{frage}
	
	\begin{frage}[Universelle Speedup-Schranken]
		Existieren universelle obere Schranken für den Speedup spalten-orientierter Methoden?
		
		Vermutung: Für sparse Matrix mit Sparsity $p$ und Eingabevektor-Sparsity $p_x$ gilt:
		\[
		\text{Speedup} \leq \frac{1}{(1-p)(1-p_x)} \cdot C(\text{arch})
		\]
		wobei $C(\text{arch})$ eine architekturspezifische Konstante ist.
		
		Beweis steht aus.
	\end{frage}
	
	\begin{frage}[Fill-in-Vermeidung]
		Für welche Klassen von sparse Matrizen bleiben die Potenzen $A^k$ sparse?
		
		Bekannt: Tridiagonale, Band-Matrizen haben kontrolliertes Fill-in.
		
		Unbekannt: Gibt es eine allgemeine Charakterisierung? Kann man Fill-in für beliebige Strukturen vorhersagen?
	\end{frage}
	
	\begin{frage}[Nichtlineare Asymmetrie]
		Lässt sich das Asymmetrie-Konzept auf nichtlineare Systeme verallgemeinern?
		
		Für $\dot{x} = f(x)$ mit Jacobi-Matrix $J(x) = \nabla f(x)$: Die Asymmetrie von $J(x)$ variiert mit $x$. Wie kann man dies ausnutzen?
		
		Möglicher Ansatz: Lokale Linearisierung + adaptive Algorithmuswahl basierend auf $J(x(t))$.
	\end{frage}
	
	\section{Abschließende philosophische Reflexion}
	
	Die Asymmetrie zwischen Zeilen und Spalten ist ein bemerkenswertes Beispiel dafür, wie mathematische Struktur physikalische Realität widerspiegelt:
	
	\textbf{Zeilen als Aggregatoren}: In Naturgesetzen, Physik und Biologie ist Aggregation ubiquitär. Elektrisches Potential aggregiert Beiträge aller Ladungen, Populationswachstum aggregiert alle ökologischen Einflüsse, neuronale Aktivierung aggregiert alle synaptischen Inputs. Die Zeilen-Perspektive kodiert diese "many-to-one"-Struktur.
	
	\textbf{Spalten als Broadcaster}: Gleichzeitig propagieren Einflüsse von wenigen Quellen zu vielen Zielen. Ein Influencer erreicht Millionen, eine Mutation breitet sich in einer Population aus, ein Hub-Neuron aktiviert viele nachgelagerte Neuronen. Die Spalten-Perspektive kodiert diese "one-to-many"-Struktur.
	
	Diese Dualität ist nicht nur mathematisch elegant, sondern fundamentales Organisationsprinzip komplexer Systeme. Die vorliegende Arbeit hat gezeigt, dass die systematische Ausnutzung dieser Asymmetrie transformative Auswirkungen auf Berechnungseffizienz haben kann – mit Speedups von 40× und mehr in realen Anwendungen.
	
	Blickt man weiter in die Zukunft, so deutet sich ein Paradigmenwechsel an: Von symmetrischer, gleichförmiger Verarbeitung (wie in klassischen Architekturen) hin zu asymmetrischer, strukturbewusster Verarbeitung (wie im Gehirn). Das menschliche Gehirn ist das ultimative asymmetrische System: Wenige Neuronen haben tausende Verbindungen, die meisten nur hunderte. Diese Asymmetrie ist kein Fehler, sondern Feature – sie ermöglicht effiziente Verarbeitung komplexer Aufgaben mit begrenzten Ressourcen.
	
	Die Erkenntnisse dieser Arbeit sind ein kleiner Schritt in Richtung solcher bio-inspirierter, strukturbewusster Berechnungssysteme. Sie zeigen: Asymmetrie ist nicht zu vermeiden – sie ist zu umarmen, zu verstehen und systematisch auszunutzen.
	
	\section{Schlusswort}
	
	Diese Arbeit begann mit der einfachen Beobachtung, dass Zeilen und Spalten einer Matrix unterschiedliche semantische Rollen spielen. Sie endet mit der Erkenntnis, dass diese Asymmetrie ein fundamentales Prinzip ist, das von der Theorie über Algorithmen bis hin zu konkreten industriellen Anwendungen durchgängig relevant ist.
	
	Die erzielten Speedups – bis zu 40× in realen Szenarien – sind nicht nur akademisch interessant. Sie ermöglichen Anwendungen, die zuvor nicht praktikabel waren: Echtzeit-MPC für komplexe Roboter, Monte-Carlo-Simulation epidemiologischer Modelle mit Millionen Szenarien, Edge-Device-Inferenz tiefer neuronaler Netze, und häufige Stress-Tests systemischer Finanzrisiken.
	
	Gleichzeitig zeigt die Arbeit auch die Grenzen: Overhead dominiert bei kleinen Systemen, hochoptimierte Bibliotheken sind schwer zu schlagen, und nicht jede strukturelle Eigenschaft lohnt die Ausnutzung. Diese Erkenntnisse sind genauso wertvoll wie die Erfolge.
	
	Die Zukunft gehört strukturbewussten, adaptiven Systemen, die automatisch erkennen, wie Berechnungen optimal durchzuführen sind. Diese Arbeit liefert theoretische Grundlagen, algorithmische Werkzeuge und praktische Leitlinien für diesen Weg. Sie zeigt: Die Asymmetrie zwischen Zeilen und Spalten ist kein Detail – sie ist Schlüssel zu effizienter Berechnung im 21. Jahrhundert.
	
	\vspace{1cm}
	
	\begin{center}
		\rule{0.5\textwidth}{0.4pt}
	\end{center}
	
	\vspace{0.5cm}
	
	\textit{Die vorliegende Arbeit hat gezeigt, dass fundamentale mathematische Strukturen – wie die Asymmetrie zwischen Zeilen und Spalten – nicht nur theoretische Eleganz besitzen, sondern praktische Auswirkungen von enormer Tragweite haben. Von der Robotik über Epidemiologie bis zur Künstlichen Intelligenz: Überall, wo Matrizen die Welt beschreiben, lauert Potenzial für dramatische Effizienzgewinne.}

\begin{thebibliography}{50}

% ── Matrixmultiplikation und Komplexität ────────────────────────────────────
\bibitem{strassen1969}
V.~Strassen.
\newblock Gaussian elimination is not optimal.
\newblock \textit{Numerische Mathematik}, 13(4):354--356, 1969.

\bibitem{coppersmith1990}
D.~Coppersmith und S.~Winograd.
\newblock Matrix multiplication via arithmetic progressions.
\newblock \textit{Journal of Symbolic Computation}, 9(3):251--280, 1990.

\bibitem{legall2014}
F.~Le~Gall.
\newblock Powers of tensors and fast matrix multiplication.
\newblock In \textit{Proceedings of the 39th International Symposium on Symbolic and Algebraic Computation (ISSAC)}, Seiten 296--303, 2014.

\bibitem{williams2012}
V.~V.~Williams.
\newblock Multiplying matrices faster than Coppersmith-Winograd.
\newblock In \textit{Proceedings of the 44th ACM Symposium on Theory of Computing (STOC)}, Seiten 887--898, 2012.

\bibitem{pan1984}
V.~Ya.~Pan.
\newblock How to multiply matrices faster.
\newblock \textit{Lecture Notes in Mathematics}, Vol.~179, Springer, Berlin, 1984.

% ── Sparse Matrizen und effiziente Algorithmen ──────────────────────────────
\bibitem{davis2006}
T.~A.~Davis.
\newblock \textit{Direct Methods for Sparse Linear Systems}.
\newblock SIAM, Philadelphia, 2006.

\bibitem{golub2013}
G.~H.~Golub und C.~F.~Van~Loan.
\newblock \textit{Matrix Computations}, 4.~Auflage.
\newblock Johns Hopkins University Press, Baltimore, 2013.

\bibitem{barrett1994}
R.~Barrett, M.~Berry, T.~F.~Chan, J.~Demmel, J.~Donato, J.~Dongarra, V.~Eijkhout, R.~Pozo, C.~Romine und H.~Van~der~Vorst.
\newblock \textit{Templates for the Solution of Linear Systems: Building Blocks for Iterative Methods}.
\newblock SIAM, Philadelphia, 1994.

\bibitem{saad2003}
Y.~Saad.
\newblock \textit{Iterative Methods for Sparse Linear Systems}, 2.~Auflage.
\newblock SIAM, Philadelphia, 2003.

\bibitem{eisenstat1988}
S.~C.~Eisenstat und J.~W.~H.~Liu.
\newblock Exploiting structural symmetry in a sparse partial factorization scheme.
\newblock \textit{SIAM Journal on Matrix Analysis and Applications}, 14(1):253--257, 1988.

% ── Matrixexponential ───────────────────────────────────────────────────────
\bibitem{moler2003}
C.~Moler und C.~Van~Loan.
\newblock Nineteen dubious ways to compute the exponential of a matrix, twenty-five years later.
\newblock \textit{SIAM Review}, 45(1):3--49, 2003.

\bibitem{higham2009}
N.~J.~Higham.
\newblock The scaling and squaring method for the matrix exponential revisited.
\newblock \textit{SIAM Journal on Matrix Analysis and Applications}, 26(4):1179--1193, 2005.

\bibitem{higham2008}
N.~J.~Higham.
\newblock \textit{Functions of Matrices: Theory and Computation}.
\newblock SIAM, Philadelphia, 2008.

\bibitem{al2009}
M.~Al-Mohy und N.~J.~Higham.
\newblock A new scaling and squaring algorithm for the matrix exponential.
\newblock \textit{SIAM Journal on Matrix Analysis and Applications}, 31(3):970--989, 2009.

\bibitem{ward1977}
R.~C.~Ward.
\newblock Numerical computation of the matrix exponential with accuracy estimate.
\newblock \textit{SIAM Journal on Numerical Analysis}, 14(4):600--610, 1977.

% ── Dynamische Systeme und Systemtheorie ────────────────────────────────────
\bibitem{kailath1980}
T.~Kailath.
\newblock \textit{Linear Systems}.
\newblock Prentice-Hall, Englewood Cliffs, NJ, 1980.

\bibitem{antsaklis2006}
P.~J.~Antsaklis und A.~N.~Michel.
\newblock \textit{A Linear Systems Primer}.
\newblock Birkhäuser, Boston, 2006.

\bibitem{chen1984}
C.-T.~Chen.
\newblock \textit{Linear System Theory and Design}.
\newblock Oxford University Press, New York, 1984.

\bibitem{khalil2002}
H.~K.~Khalil.
\newblock \textit{Nonlinear Systems}, 3.~Auflage.
\newblock Prentice-Hall, Upper Saddle River, NJ, 2002.

\bibitem{horn2013}
R.~A.~Horn und C.~R.~Johnson.
\newblock \textit{Matrix Analysis}, 2.~Auflage.
\newblock Cambridge University Press, Cambridge, 2013.

% ── Graphentheorie und Boolean-Matrizen ─────────────────────────────────────
\bibitem{diestel2017}
R.~Diestel.
\newblock \textit{Graph Theory}, 5.~Auflage.
\newblock Springer, Berlin, 2017.

\bibitem{boolean_mm}
H.~S.~Warren.
\newblock \textit{Hacker's Delight}, 2.~Auflage.
\newblock Addison-Wesley, Boston, 2013.

\bibitem{cormen2009}
T.~H.~Cormen, C.~E.~Leiserson, R.~L.~Rivest und C.~Stein.
\newblock \textit{Introduction to Algorithms}, 3.~Auflage.
\newblock MIT Press, Cambridge, MA, 2009.

\bibitem{tarjan1972}
R.~Tarjan.
\newblock Depth-first search and linear graph algorithms.
\newblock \textit{SIAM Journal on Computing}, 1(2):146--160, 1972.

\bibitem{kahn1962}
A.~B.~Kahn.
\newblock Topological sorting of large networks.
\newblock \textit{Communications of the ACM}, 5(11):558--562, 1962.

% ── Neuronale Netze und Deep Learning ──────────────────────────────────────
\bibitem{lecun2015}
Y.~LeCun, Y.~Bengio und G.~Hinton.
\newblock Deep learning.
\newblock \textit{Nature}, 521(7553):436--444, 2015.

\bibitem{goodfellow2016}
I.~Goodfellow, Y.~Bengio und A.~Courville.
\newblock \textit{Deep Learning}.
\newblock MIT Press, Cambridge, MA, 2016.

\bibitem{han2015}
S.~Han, J.~Pool, J.~Tran und W.~J.~Dally.
\newblock Learning both weights and connections for efficient neural networks.
\newblock In \textit{Advances in Neural Information Processing Systems (NeurIPS)}, Seiten 1135--1143, 2015.

\bibitem{frankle2019}
J.~Frankle und M.~Carlin.
\newblock The lottery ticket hypothesis: Finding sparse, trainable neural networks.
\newblock In \textit{International Conference on Learning Representations (ICLR)}, 2019.

\bibitem{glorot2011}
X.~Glorot, A.~Bordes und Y.~Bengio.
\newblock Deep sparse rectifier neural networks.
\newblock In \textit{Proceedings of the 14th International Conference on Artificial Intelligence and Statistics (AISTATS)}, Seiten 315--323, 2011.

% ── Numerische lineare Algebra und HPC ─────────────────────────────────────
\bibitem{dongarra1990}
J.~J.~Dongarra, J.~Du~Croz, S.~Hammarling und I.~Duff.
\newblock A set of level 3 basic linear algebra subprograms.
\newblock \textit{ACM Transactions on Mathematical Software}, 16(1):1--17, 1990.

\bibitem{anderson1999}
E.~Anderson, Z.~Bai, C.~Bischof, S.~Blackford, J.~Demmel, J.~Dongarra, J.~Du~Croz, A.~Greenbaum, S.~Hammarling, A.~McKenney und D.~Sorensen.
\newblock \textit{LAPACK Users' Guide}, 3.~Auflage.
\newblock SIAM, Philadelphia, 1999.

\bibitem{van2011}
S.~van~der~Walt, S.~C.~Colbert und G.~Varoquaux.
\newblock The NumPy array: A structure for efficient numerical computation.
\newblock \textit{Computing in Science \& Engineering}, 13(2):22--30, 2011.

\bibitem{scipy2020}
P.~Virtanen et al.
\newblock SciPy 1.0: Fundamental algorithms for scientific computing in Python.
\newblock \textit{Nature Methods}, 17(3):261--272, 2020.

\bibitem{intel_mkl}
Intel Corporation.
\newblock \textit{Intel Math Kernel Library Reference Manual}.
\newblock Intel Press, Santa Clara, CA, 2020.

% ── Modellprädiktive Regelung ───────────────────────────────────────────────
\bibitem{rawlings2017}
J.~B.~Rawlings, D.~Q.~Mayne und M.~M.~Diehl.
\newblock \textit{Model Predictive Control: Theory, Computation, and Design}, 2.~Auflage.
\newblock Nob Hill Publishing, Madison, WI, 2017.

\bibitem{mayne2000}
D.~Q.~Mayne, J.~B.~Rawlings, C.~V.~Rao und P.~O.~M.~Scokaert.
\newblock Constrained model predictive control: Stability and optimality.
\newblock \textit{Automatica}, 36(6):789--814, 2000.

\bibitem{qin2003}
S.~J.~Qin und T.~A.~Badgwell.
\newblock A survey of industrial model predictive control technology.
\newblock \textit{Control Engineering Practice}, 11(7):733--764, 2003.

% ── Strukturausnutzung und Algorithmendesign ────────────────────────────────
\bibitem{brent1970}
R.~P.~Brent.
\newblock Algorithms for matrix multiplication.
\newblock Technical Report CS-TR-70-157, Stanford University, 1970.

\bibitem{ibarra1982}
O.~H.~Ibarra und S.~Moran.
\newblock Probabilistic algorithms for deciding equivalence of straight-line programs.
\newblock \textit{Journal of the ACM}, 30(1):217--228, 1983.

\bibitem{duff1986}
I.~S.~Duff, A.~M.~Erisman und J.~K.~Reid.
\newblock \textit{Direct Methods for Sparse Matrices}.
\newblock Oxford University Press, Oxford, 1986.

\bibitem{george1981}
A.~George und J.~W.~H.~Liu.
\newblock \textit{Computer Solution of Large Sparse Positive Definite Systems}.
\newblock Prentice-Hall, Englewood Cliffs, NJ, 1981.

\bibitem{parter1961}
S.~V.~Parter.
\newblock The use of linear graphs in Gauss elimination.
\newblock \textit{SIAM Review}, 3(2):119--130, 1961.

% ── Parallelrechnung und GPU-Computing ─────────────────────────────────────
\bibitem{nvidia2023}
NVIDIA Corporation.
\newblock \textit{CUDA C++ Programming Guide}, Version 12.0.
\newblock NVIDIA Corporation, Santa Clara, CA, 2023.

\bibitem{kirk2016}
D.~B.~Kirk und W.~W.~Hwu.
\newblock \textit{Programming Massively Parallel Processors}, 3.~Auflage.
\newblock Morgan Kaufmann, Burlington, MA, 2016.

\bibitem{bell2009}
N.~Bell und M.~Garland.
\newblock Implementing sparse matrix-vector multiplication on throughput-oriented processors.
\newblock In \textit{Proceedings of the 2009 ACM/IEEE Conference on Supercomputing (SC09)}, Artikel~18, 2009.

% ── Anwendungsgebiete ───────────────────────────────────────────────────────
\bibitem{varga2000}
A.~Varga.
\newblock \textit{Matrix Analysis for Scientists and Engineers}.
\newblock SIAM, Philadelphia, 2004.

\bibitem{strang2016}
G.~Strang.
\newblock \textit{Introduction to Linear Algebra}, 5.~Auflage.
\newblock Wellesley-Cambridge Press, Wellesley, MA, 2016.

\bibitem{trefethen1997}
L.~N.~Trefethen und D.~Bau~III.
\newblock \textit{Numerical Linear Algebra}.
\newblock SIAM, Philadelphia, 1997.

\bibitem{higham2002}
N.~J.~Higham.
\newblock \textit{Accuracy and Stability of Numerical Algorithms}, 2.~Auflage.
\newblock SIAM, Philadelphia, 2002.

\bibitem{demmel1997}
J.~W.~Demmel.
\newblock \textit{Applied Numerical Linear Algebra}.
\newblock SIAM, Philadelphia, 1997.

\bibitem{watkins2010}
D.~S.~Watkins.
\newblock \textit{Fundamentals of Matrix Computations}, 3.~Auflage.
\newblock Wiley, Hoboken, NJ, 2010.

\bibitem{anderson2010}
B.~D.~O.~Anderson und J.~B.~Moore.
\newblock \textit{Optimal Control: Linear Quadratic Methods}.
\newblock Dover Publications, Mineola, NY, 2007.

\bibitem{mezard2009}
M.~Mézard und A.~Montanari.
\newblock \textit{Information, Physics, and Computation}.
\newblock Oxford University Press, Oxford, 2009.

\end{thebibliography}

\end{document}